\subsection{Ordinary inert primes on the retained cut}
\label{subsec:inert-rows}

At the ordinary inert primes the split column condition disappears. The row calculation must instead follow the signed jumps around a complete Frobenius orbit in the unramified local field.

Use the ordinary prime range of Section~\ref{subsec:ordinary-split-length}.
An inert prime in this range has \(\ell\ge11\) once the family scale exceeds its lower threshold. Every place \(\lambda\mid\ell\) is a row. Its completion
\(F=E_\lambda\) is unramified over \(\Q_\ell\), with degree \(d_\lambda\).
The cyclic index at \(\lambda\) runs over
\[
 \operatorname{Hom}_{\mathbb Q_\ell}(E_\lambda,\overline{\mathbb Q}_\ell),\qquad \#\operatorname{Hom}_{\mathbb Q_\ell}(E_\lambda,\overline{\mathbb Q}_\ell)=d_\lambda.
\]
The ordering comes from Frobenius on the actual host, including its quadratic factors.

We describe the Wach levels and compare the arcs at these
inert primes.

Use the marked row surjection \(T_1\twoheadrightarrow T_2\) of
Section~\ref{subsec:split-local-length}. The target is the rank-\(r\)
Tate dual at the selected point and \(T_1\) is a lattice in a
subrepresentation of a finite sum of the Tate duals of the normalized
retained representations. The rational comparison is crystalline.
The local jump windows hold on all of the summands on this cut.
Strictness for the rational crystalline filtration passes them to
\(T_1\) and to \(T_2\). For the marked rank-\(r\) target the central multiplicity bound is pointwise in the letter:
\[
 \#\{\text{central target jumps at }i\}\le m\qquad(\text{every }i),
\]
where multiplicities are counted.
Section~\ref{subsec:cut-wach-models} gives the signed integral
Wach inclusions on these same two lattices and their local finite
extensions, at the labeled places of the retained tuple.

The common word is obtained by restricting the host weight chart
to the Frobenius cycle of \(\lambda\). Its signs are the two CM
signs for the Tate dual. At an inert prime Frobenius exchanges
these signs, which accounts for the alternating words and the
even local degree. Rounding the upper interval and its width
gives local bounds \(b_0,w,S,\eta\) as below, within the common bounds
\[
 w\le w_P,\qquad S\le S_c,\qquad \eta\le\eta_3,\qquad b_0+S+\eta\le H.
\]
The shifted place weights obey the same
central gap at every retained point. These comparisons concern
the surjection and jump data used by the congruence determinant,
including the case of a ramified extension of the coefficient DVR.
\label{cor:inert-row-lengths-1}

\paragraph{The inert contributions to the weighted prime sum.}
\label{sec:cut-length-rows}
Use \(H,G_P,w_P,\eta_3,z\) from
Section~\ref{subsec:prime-cost-sums}. Write \(A_{\rm sec}\) for
the secondary threshold denoted \(A\) there and \(o_\ell\) for
the order in the host's real cyclotomic quotient. Case B has
\(\ell<A_{\rm sec}\) and \(o_\ell>4L_w\). Set
\[
 j_\ell=\max(1,\lceil\log_{\ell-1}A_{\rm sec}\rceil),\qquad
 U_B=(d_\lambda+1)r
     \left(\frac{H}{(\ell-1)^{j_\ell+1}}+
          (G_P+1)(j_\ell+1)+z\right).
\]
The remaining ordinary inert primes are case C if
\(\ell\ge A_{\rm sec}\), and case D if \(\ell<A_{\rm sec}\) and
\(o_\ell\le4L_w\). The local comparison distinguishes two
cases. Larger jumps occur at every position precisely when
Frobenius acts as complex conjugation and the starting
position has larger jumps.

To see both alternatives, first consider \(o_\ell\le4L_w\).
Two visits to the large set with labels \(s_1,s_2\) satisfy
\[
 D\mid o_\ell t_0(s_2-s_1),\qquad
 n_0\mid o_\ell(s_2-s_1),\qquad
 |o_\ell(s_2-s_1)|<n_0 .
\]

The first relation follows from \eqref{eq:translated-dyadic-visits},
whose coordinate formula holds at every order. Thus \(s_1=s_2\).
Each label specifies a single real place, represented by the two
conjugate embeddings of \(E\). Consecutive larger positions would
then give the same real place, so Frobenius fixes an embedding of
\(E^+\). Since \(E/\mathbb Q\) is abelian, it belongs to
\(\operatorname{Gal}(E/E^+)=\{1,c\}\). Inertness in
\(\mathbb Q(\sqrt{-d})\) rules out 1. For Frobenius \(c\),
a larger starting place remains larger throughout its cycle.
In the long-order case the three-visit argument of
Lemma~\ref{lem:long-orbit-spacing} applies for any ordinary inert
prime, since it uses only \(o_\ell>4L_w\) and the designated
labels. As \(L_w\ge3\), at most two consecutive positions are larger.

Use the place indicator \(B_\lambda\)
that a good point over this place exhibits the short band of
Section~\ref{subsec:ordinary-split-length}. Define
\[
\begin{aligned}
 \mathcal P_\ell&=1+H/\ell+\log_\ell H,&
 F_{\ell,\lambda}&=(w_P+\eta_3)/\ell+S_c+z+\eta_3B_\lambda,\\
 U_H&=(d_\lambda+1)rH/\ell^2,&
 U_F&=(d_\lambda+1)r F_{\ell,\lambda},&
 U_M&=d_\lambda m\mathcal P_\ell .
\end{aligned}
\]
The following estimates yield one rational bound at each place,
valid simultaneously for its good points and coefficient lattices.
The case B entry is \(O(U_B)\), and the common case C and D
entry is \(O(U_H+U_F+U_M)\). Their constants are absolute once
the host constants are fixed. In particular they are independent of
the degree of the coefficient field and the jump heights.
\label{cor:inert-row-lengths-2}

Every \(d_\lambda\ge1\), so the sum over all row places satisfies
\[
 \sum_{\lambda\mid\ell}(d_\lambda+1)\le2\sum_{\lambda\mid\ell} d_\lambda=4f.
\]
The global indicator \(I_\ell\) in
Section~\ref{subsec:prime-cost-sums} bounds every \(B_\lambda\)
by the band transport there. Thus the \(U_B\), \(U_H\) and \(U_F\)
terms give the three inert costs in that section, after
enlarging the fixed common constant. The \(U_M\) terms contribute
at most \(2fm\mathcal P_\ell\) before that constant and go into
the same sum as the split lifting term. The factor in the
global determinant is absorbed by the common constant.
At these inert primes the column set is empty, so the remainder
is the boundary term with \(k_c=0\) in the split calculation.

For case B, it remains to obtain the required row colength from the length of the Frobenius orbit; this is Proposition~\ref{prop:inert-long-colength}.

\subsubsection{The long-orbit row bound}\label{sec:inert-long-orbit}
\hypertarget{source-inert-long-orbit}{}

Among the inert orbits, a long Frobenius orbit allows the marked jumps to be cut into shorter arcs. The colength can be bounded using those arcs only if the correction at each cut is controlled integrally.

\begin{established}[Case B and the cyclic signed jump windows]\label{hyp:long-orbit-data}
Let \(\ell<A_{\rm sec}\) be an ordinary inert prime whose Frobenius
has order \(o_\ell>4L_w\) in the real cyclotomic quotient
(case B of Section~\ref{sec:cut-length-rows}). Write \(F=E_\lambda\) for the
unramified completion of the auxiliary field and
\[
 d=[F:\Q_\ell]=n_\lambda.
\]

Work over the integers \(O\) of a local coefficient field splitting \(F\).
Write \(T_2\) for the marked dual Tate lattice, of rank \(r\), and
\(T_1\) for a lattice in a subrepresentation of a finite sum of the
twisted duals from the same cut. The maps to be estimated, including the marked matrix row, are precisely the indicated integral Galois surjections
\[
 f:T_1\twoheadrightarrow T_2,\qquad \operatorname{rank}_O T_2=r .
\]

The source and target are crystalline with central jumps \(|a|\le\eta_0\) and one outer signed interval of width \(S_0\) at each Frobenius embedding. The small letters have height at most \(G_P\) and all letters have height at most \(H\), in the notation of Section~\ref{sec:cut-length-rows}.
\end{established}
\begin{proof}
The ordinary inert branch excludes all primes ramifying in the field or in the rows. Conjugating its marked row preserves the coefficient lattice and takes its dual Tate jumps to the reflected signed intervals. Apply Lemma~\ref{lem:marked-row-column-lattices} to obtain the displayed surjection. Lemma~\ref{lem:marked-signed-jumps} gives the central and outer bands before passage to its quotient. Both rational representations are crystalline with the
local jump windows from the cut. Enumerate the embeddings of \(F\)
cyclically by Frobenius. The residue characteristic is \(\ell\ge11\)
in this inert branch.

Use \(b_0,w_0,S_0,\eta_0\) for the local size parameters, with
\(S_0,\eta_0\ge1\), and \(s_i\) for the size at position \(i\).
The source jumps at \(i\) lie in \([-\eta_0,\eta_0]\) together
with either \([s_i,s_i+S_0]\) or its negative. The larger range is
\[
 b_0\le s_i\le b_0+\eta_0.
\]

At every other position,
\(s_i\le w_0\). Put \(g=w_0+S_0+\eta_0\) and set
\[
 h=\begin{cases}b_0+\eta_0+S_0,&\text{if some position has larger jumps},\\
                  g,&\text{otherwise}.\end{cases}
\]
Thus the positions with smaller jumps have jumps in \([-g,g]\) and all positions
have jumps in \([-h,h]\). The target satisfies the same bounds by
strict exactness of the rational crystalline filtration.
The two upper endpoints in Section~\ref{sec:cut-length-rows} bound \(b_0+\eta_0+S_0\le H\)
and \(g\le G_P\), with \(H,G_P\) as in Section~\ref{sec:cut-length-rows}.
The length of the cycle here is the unramified local degree, including
the quadratic components of Frobenius.
\end{proof}

\begin{auxlemma}[Spacing of the larger jumps]\label{lem:long-orbit-spacing}
On the unramified cycle of Lemma~\ref{hyp:long-orbit-data}, let
\(J\) be the positions with larger jumps and \(L_w\) the length
parameter for \(T_2\) in Lemma~\ref{lem:second-set-progressions}.
Write \(\theta_\ell\in\mathbb Z/D\mathbb Z\) for the action of
arithmetic Frobenius on \(C_e\). If
\(o_\ell=\operatorname{ord}(\theta_\ell)>4L_w\), then
\[
 \#(J\cap I)\le2\qquad\text{for every interval of }L_w
                  \text{ consecutive positions }I .
\]
If \(1\le j_\ell\) and \(2j_\ell+2\le L_w\), the cycle admits
cuts with \(j_\ell\) smaller positions following each cut, at most
two larger positions per arc, and maximum arc length \(3j_\ell+1\).
\end{auxlemma}
\begin{proof}
In \(\operatorname{Gal}(C_e/\mathbb Q)\), the step \(\theta_\ell\)
is the class of \(\ell\bmod 2^{e+2}\) up to sign, expressed in the
generator 5. Thus
\[
   o_\ell=\min\{a\ge1:\ell^a\equiv\pm1\pmod{2^{e+2}}\},
    \qquad o_\ell\mid [E_\lambda:\mathbb Q_\ell] .
\]
The divisibility holds because the order of the full Frobenius
equals the unramified local degree and restriction preserves powers.

By Lemma~\ref{hyp:translate-weights} the larger places are the
translated labels \(e_1(s+u-t)\) and their complex conjugates
at the two embeddings of a real place when the coordinates are
centred at the selected top \(t\). Restriction to the real
field \(C_e\) kills complex conjugation. Let \(w\) be the
coordinate of the translate, equal to \(\rho(u-t)\) with this
origin or \(\rho(u)\) before recentering. With \(s_i\) the integer representative of the larger label at
position \(i\), \eqref{eq:dyadic-coordinate-orders} gives
\begin{equation}\label{eq:translated-dyadic-visits}
 x_0+i\theta_\ell=w+t_0 s_i\pmod D,\qquad
                  0\le 4L_w s_i<n_0\quad(i\in J).
\end{equation}

Here \(x_0\) is the starting coordinate. Reversing the orientation
replaces \(\theta_\ell\) by its negative and preserves its order.
If \(a<b<c\) in one interval of length \(L_w\) are larger
positions, their coordinates are distinct by the order bound.
Put \(u'=b-a\), \(v'=c-a\). Eliminating the step gives
\[
 t_0\big(v'(s_b-s_a)-u'(s_c-s_a)\big)\equiv0\pmod D,
 \qquad
 v'(s_b-s_a)-u'(s_c-s_a)\equiv0\pmod{n_0}.
\]

The second implication uses \(D=t_0n_0\). The four terms have
absolute value less than \(n_0/4\) by \eqref{eq:translated-dyadic-visits},
so the last congruence is an integer equality between three
distinct selected labels. The sphere construction in
Lemma~\ref{lem:second-set-progressions} excludes this equality.
Taking integer lifts of the positions includes intervals crossing
the end of the cycle.

Let \(j=j_\ell\). The numerical host choice allows some exponent
\(j'\) with \(A_{\rm sec}\le(\ell-1)^{j'}\) and \(2j'+2\le L_w\).
Here \(\ell<A_{\rm sec}\), so the least exponent \(j\) is at least two
and satisfies
\[
 2j+2\le L_w.
\]
Thus the local cycle has \(d>2j\).
Join successive positions with larger jumps when their forward gap is at most
\(j\). A chain of three would occupy at most \(2j+1\) consecutive
positions and is excluded. The clusters thus have one or two
positions, with at least \(j\) intervening positions with smaller jumps between
clusters. Cut just after the last position with larger jumps of each cluster,
using that position as the cut index \(c\). Each of \(c+1,\dots,c+j\)
then has smaller jumps.

If the next cluster starts at forward distance \(a\),
insert cuts at the multiples of \(j\) up to \(a-j-1\), measured
from the current cut. Include the original cut at distance zero. If
\(k=\lfloor(a-j-1)/j\rfloor\), the remaining gap is
\[
 a-kj=j+1+\bigl((a-j-1)\bmod j\bigr)\in[j+1,2j].
\] Adding the span of the next
cluster gives an arc length at most \(3j\). Each arc is taken
half-open, with its terminal cut included.

Each such arc contains at most two
positions with larger jumps. When there are no positions with larger jumps, partition the
cycle into arcs of length \(j\) with the last remainder added to one
arc. Each arc has length at least \(j\), since inserted cuts are
spaced by \(j\) and each gap to a cluster of positions with larger jumps starts with its
\(j\)-letter buffer of positions with smaller jumps. The maximum length \(A\) satisfies
\(j\le A\le3j+1\).
\end{proof}

\paragraph{The exponent in the Wach lifting estimate.}\label{par:inert-cut-exponents}
The correction at a cut in Paragraph~\ref{par:wach-cut-correction} and the integral level tests of Proposition~\ref{prop:wach-arc-membership} apply to the
cyclotomic Wach modules of a surjection of crystalline lattices over
an unramified field. Its arc hypotheses are a cyclic partition with
\(j\) positions with smaller jumps after each cut, at most two positions with larger jumps per
arc, and specified size inequalities when their signs differ.
Lemma~\ref{lem:long-orbit-spacing} supplies the partition.

We now compute the exponents for the symmetric jump boxes of
Lemma~\ref{hyp:long-orbit-data}.
Put \(q_0=\ell-1\), \(D_0=q_0^{j+1}\), and
\[
 \begin{split}
 M&=\left\lfloor\frac{h+2g}{q_0}\right\rfloor+1,\\
 \sigma&=\left\lceil 1+\frac{\ell}{q_0^2}
                      \left(3g+\frac{M}{\ell^{j-1}}\right)\right\rceil,\\
 H_i&=\begin{cases}2h,&i\text{ with larger jumps},\\2g,&i\text{ with smaller jumps},\end{cases}\\
 c_H&=\left\lceil 1+\log_\ell(1+\max_i H_i)\right\rceil,\\
 S(c)&=\sum_{a\ge0}\frac{H_{c+a+1}+c_H+3}{q_0^a},\\
 R&=\max_c\left(c_H+\left\lfloor\frac{\ell}{q_0^2}
           \min\{\lfloor S(c)\rfloor,\max_i H_i\}\right\rfloor\right).
 \end{split}
\]
The initial exponent is \(a_0=2\sigma+R+2gA\). Define \(\kappa\)
as the maximum of
\[
       A+\sum_{u\in J\setminus\{0\}}(A+1+v_\ell(u))
\]
over unions \(J\) of three intervals of at most \(3g\) consecutive
integers in \([-2h,2h]\).

In Theorem E the inputs for the source and
target can both be taken to be \(H_i\). If the positions with larger jumps of
opposite sign have magnitudes \(x\) at backward offset \(s_1\) and
\(y\) at offset \(s_0<s_1\), the required comparisons are
\(y<\ell x\) when the earlier position in the cycle is negative,
and \(2y\le\ell x\) and \(2a_0<x\) when it is positive. With these
inputs the theorem lifts every integral finite class of \(T_2\) after
multiplication by \(\ell^e\), where
\begin{equation}\label{eq:long-orbit-exponent}
                 e=2R+4\sigma+2a_0+2\kappa.
\end{equation}

For the proof with cyclotomic evaluation coordinates, extend the
coefficient field to contain the primitive \(\ell^d\)-th roots of
unity. This extends the coefficients while keeping \(F\) and its
Frobenius cycle fixed. The embeddings over the larger coefficient
field are the same after composition, since the original field
already splits \(F\). The finite cohomology group of a lattice
is the preimage of the rational finite group.

A basis of the
extended coefficient ring over \(O\) identifies its continuous
cochains with finitely many copies of the original complex.
Rationalization and the crystalline period comparison also commute
with this finite direct sum. Thus integral finite cohomology, its
map and the cokernel extend by tensor product. For the cokernel \(C\) and extended coefficient DVR \(O'\), faithful flatness gives
\[
 \ell^e(C\otimes_O O')=0\quad\Longrightarrow\quad\ell^e C=0.
\]


The geometric estimate for \(R\) exhibits the dependence on the
buffer of positions with smaller jumps. The first \(j\) terms of the sum at each cut have
\(H_{c+a+1}=2g\), hence
\[
 S(c)\le \frac{q_0}{q_0-1}
          \left(2g+c_H+3+\frac{2h}{q_0^j}\right).
\]
Since \(\ell\ge11\) and \(h\ge1\), we have
\(c_H\le3+\log_\ell h\) and
\[
  R\le8g+3h/D_0+2\log_\ell h,
  \qquad \sigma\le3g+2h/D_0.
\]
Indeed \(\ell/(q_0(q_0-1))\le11/90\) and
\(2\ell/(q_0-1)\le22/9\) give the first inequality.

For the second, \(\ell/q_0^2\le11/100\), and the coefficient
of \(h+2g\) is at most \((11/10)/D_0\) for \(j\ge1\).
Choose the common constant \(C_2\ge3200\) in the size condition
\(b_0\ge C_2(j+1)(g+1)\). This condition holds uniformly for
the rounded local word. In fact \(b_0\ge b-(t+1)/2\ge b/2\)
at large scale, \(g+1\le G_P+1\le11b'\), and
\(j+1\le3+\log A_{\rm sec}=O(\log x)\). Hence
\[
 C_2(j+1)(g+1)/b_0=O(C_2\log x\, b'/b),
\]
which tends
to zero since
\[
 C_2\log x\,b'/b=C_2x\log x\exp(-x^{\varepsilon/3})\longrightarrow0.
\]

If there are positions with larger jumps,
\(g\le b_0\), \(h\le2b_0\), and both magnitudes are between
\(b_0\) and \(b_0+\eta_0\). Therefore \(2y\le4b_0<\ell x\).
The bounds on \(\sigma,R\) in the expression \(a_0=2\sigma+R+2gA\) give
\[
 2a_0\le40g(j+1)+14h/D_0+4\log_\ell h<b_0.
\]
For example, the three terms are bounded by \(0.04b_0\),
\(0.28b_0\), and \(4\sqrt{2b_0}\) using \(b_0\ge2000\) and
\(\log_\ell h\le\sqrt h\). The two mixed-sign comparisons follow.
When all positions have smaller jumps these comparisons are vacuous.

To bound \(\kappa\), each interval with at most \(3g\) integers
contains at most \(3g/\ell^v+1\) nonzero multiples of \(\ell^v\),
and such multiples require \(v\le\log_\ell(2h)\). Summing over
the three intervals gives
\(\kappa\le25gA+3\log_\ell h\) since \(A,g\ge1\).
Using \eqref{eq:long-orbit-exponent},
\[
 e\le110gA+28h/D_0+14\log_\ell h
   \le C_2\left(h/D_0+(g+1)(j+1)+\log_\ell h\right).
\]
In particular the high jump bound carries the factor
\((\ell-1)^{-j-1}\) even when the cycle has many clusters. The
exponents use a maximum over the arcs. The factor for the host
degree enters only in the dimension of finite cohomology.

\paragraph{Correcting the lattice at a Frobenius cut.}\label{par:wach-cut-correction}
In the ordinary cyclotomic variable \(\pi\), let \(\Delta\)
be the tame group of order \(\ell-1\) and take
\(X=-\prod_{\delta\in\Delta}\delta(\pi)\).
\begin{lemma}[The tame-invariant parameter]
The tame-invariant coefficient ring is \(O[[X]]\). Put \(q=X+\ell\)
and \(p_s=\phi^s(q)\). Up to a unit, \(p_s\) is a distinguished
polynomial of degree \(\ell^s\) with constant term associated to
\(\ell\), and for each \(j\ge1\),
\[
 (\phi^j X)=(X)\prod_{0\le s<j}(p_s)^{\ell-1}.
\]
\end{lemma}
\begin{proof}
The defining product is a unit times \(\pi^{\ell-1}\). On the invariant ring
write \(\mathcal R=O[[X]]\). Indeed the action on degree \(k\) in the \(\pi\)-adic
filtration is by the \(k\)-th power of the tame character.
An invariant series can only have its first nonzero term in
degree a multiple of \(\ell-1\). Successively subtracting
a constant multiple of the corresponding power of \(X\)
expresses it as a series in \(X\).
At a primitive \(\ell\)-th root of unity minus one,
\(X\) evaluates to \(-\ell\), by the product formula for
all primitive \(\ell\)-th roots. Thus \(q=X+\ell\) is
associated to the usual cyclotomic factor in \(O[[\pi]]\)
by Weierstrass preparation.

As a result
\[
 \phi(X)=uXq^{\ell-1},\qquad
 \gamma(X)=\omega_X X\pmod{X^2},\qquad
 v_\ell(\omega_X-1)=1 ,
\]
where \(u\) is a tame-invariant unit and
\(\omega_X=\chi(\gamma)^{\ell-1}\).
Here \(\phi\) acts on coefficients in
\(\mathcal R\) and \(\Phi\) is the corresponding semilinear operator
between letters. Put \(p_s=\phi^s(q)\). Iteration gives
\(\phi^j(X)\) associated to \(X\prod_{s<j}p_s^{\ell-1}\), and
\(p_s\) lies in \((\ell,X^{\ell^s})\). The latter follows by
reducing the Frobenius formula modulo \(\ell\). By Weierstrass
preparation, \(p_s\) is a unit times a monic polynomial of degree
\(\ell^s\) whose other coefficients are divisible by \(\ell\)
and whose constant coefficient is associated to \(\ell\).
Different \(s\) have different cyclotomic zero sets and their
distinguished polynomials are relatively prime over
\(\operatorname{Frac}(O)\).

\end{proof}
For the free Wach lattices \(N_i\), use the height bounds
\(q^h\Phi N_i\subset N_{i+1}\) everywhere and
\(q^g\Phi N_i\subset N_{i+1}\) in the first \(j\) steps after a
cut. The cyclotomic action is trivial on \(N_i/XN_i\). Write
\(P_c=\mathcal R\Phi^j N_c,\ D=X^M\prod_{s<j}p_s^{3g}\).

\begin{lemma}[Correction at a cut]
Use the chosen positive integers \(M,j\) and nonnegative integer
\(\sigma\) with \((\ell-1)M>h,\ (\ell-1)M\ge3g\) and
\((\ell-1)^2\sigma/\ell-3g\ge M/\ell^{j-1}\). In the odd-prime
range of the cut, each \(z\in XN_c\) has a correction \(\tau\in P_c\)
satisfying
\[
 \ell^\sigma\Phi^j z-(\gamma-1)\tau\in DP_c\subset X^M N_{c+j}.
\]
Its class modulo \(DP_c\) is uniquely determined by the two
congruences for \(v,\tau\) below, is additive in \(z\), and is
compatible with maps of the Wach lattices.
\end{lemma}
\begin{proof}
The height bounds and the parameter formula give
\[
 \Phi(X^M N_i)\subset q^{(\ell-1)M-h}X^MN_{i+1},\qquad
 \Phi^j(X^MN_c)\subset DP_c\subset X^MN_{c+j}.
\]
The first inclusion yields contraction of \(\Phi\) on the cyclic
product of these deep lattices. For the second use
\[
 (\prod_{s<j}p_s^g)P_c\subset N_{c+j},
\]
by iteration.
Frobenius sends a basis to an independent family after inverting
the cyclotomic factors by the two-sided height comparison. It thus
gives a basis of \(P_c\). The image of \(D\) modulo a uniformizer
is a nonzero power series. Componentwise reduction in this basis
then shows that \(P_c/DP_c\) has no \(\ell\)-torsion.
Both lattices are \(\gamma\)-stable by commutation with \(\Phi\)
and stability of the ideals \((X),(p_s)\).
Given \(z\in XN_c\), first solve the equation for
\(\gamma-1\) modulo \(X^M\).

On degree \(r\) it acts by
\(\omega_X^r-1\). With
\(A_r=\sum_{k=1}^r(1+v_\ell(k))\) and \(K=A_{M-1}\), there are
\(d_r\in N_c\) such that
\[
 v=\sum_{1\le r<M}\ell^{K-A_r}X^r d_r,\qquad
 (\gamma-1)v-\ell^K z\in X^M N_c.
\]
To construct \(v\), at degree \(r\) multiply the previous
approximation by \(\ell^{1+v_\ell(r)}\) and subtract an
\(X^r\)-term given by the inverse of the unit part of
\(\omega_X^r-1\) times the degree-\(r\) residual coefficient.
This eliminates the residual at that degree. Legendre's formula
gives \(A_r\le\ell r/(\ell-1)\).

For \(1\le r<M\) one can divide
\(\ell^\sigma(\phi^j X)^r\) by \(\ell^{A_r}\) modulo \(D\).
For \(A_r\le\sigma\) this is literal division. Otherwise let
\(d_j=\ell^{j-1}\) and
\(m=(\ell-1)r-3g-(A_r-\sigma)\). From the bound on \(A_r\)
and the definition of \(\sigma\),
\[
 m\ge\left(\frac{(\ell-1)^2}{\ell}-1\right)A_r-3g+\sigma
   \ge\frac{(\ell-1)^2}{\ell}\sigma-3g\ge M/d_j.
\]

The binomial expansion of \(p_{j-1}^{A_r-\sigma+m}\), using
\(p_{j-1}\in(\ell,X^{d_j})\), lies in
\((\ell^{A_r-\sigma},X^{d_j m})\).
After multiplying by the other factors in \((\phi^j X)^r\),
the second summand is divisible by \(D\). Choose \(b_r\in
\mathcal R\) with
\(\ell^{A_r} b_r\equiv\ell^\sigma(\phi^jX)^r\pmod{D}\).
Then
\[
 \tau=\sum_{1\le r<M} b_r\Phi^j d_r\in P_c,\qquad
 \ell^K\tau-\ell^\sigma\Phi^j v\in DP_c .
\]
Commuting \(\gamma\) past \(\Phi^j\) and using the equation for
\(v\) shows that
\[
 \ell^K(\ell^\sigma\Phi^j z-(\gamma-1)\tau)\in DP_c.
\]

Saturation removes \(\ell^K\), and
\[
 DP_c\subset X^M N_{c+j}
\]
gives the required correction. If two corrections solve the
two displayed congruences for the same \(z\), take their
difference and bring their powers of \(\ell\) to a common power.
The associated difference of the \(v\)'s is in \(XN_c\) and
has \(\gamma-1\) image in \(X^MN_c\).

If \(v'\) denotes that difference, reduction in each successive degree gives
\[
 (\omega_X^r-1)[v']_r=0\ \Longrightarrow\ [v']_r=0
 \qquad(1\le r<M).
\]
The scalar is nonzero on a free graded piece, so induction gives \(v'\in X^MN_c\). Saturation gives a difference of corrections in
\(DP_c\). This also proves compatibility with maps and addition
of the congruences, by applying the equations before comparing
corrections. For a change of the cut vector by \(u\in XN_c\) the
correction to \((\gamma-1)u\) can be taken to be
\(\ell^\sigma\Phi^j u\), since one may use \(v=u\), \(K=0\)
in those congruences. This is the comparison modulo \(XN_c\)
needed in the Herr assembly.

\end{proof}
\paragraph{The cyclotomic level models.}\label{par:wach-level-models}
Fix the same crystalline lattice map \(f:T_1\to T_2\) as in
Section~\ref{sec:inert-long-orbit}. Write \(K_O=\operatorname{Frac}(O)\)
for the coefficient field, and retain the choice containing the
cyclotomic roots. Consider an arc of length \(a\) from \(c-a\) to
\(c\). At backward offset \(s\), the letter is \(c-s\). A jump of
the decreasing Hodge filtration at \(u\) means that
\(\operatorname{Fil}^u\ne\operatorname{Fil}^{u+1}\).
Let \(G_0\le g\) bound the absolute jumps at thin letters. We
can take \(G_0=g\) for the given thin box.

Write \(b_s=s_{c-s}\)
for the size of a position with larger jumps. Upper and lower bounds \(m_s^+\)
and \(-m_s^-\) for the jumps, and the exponent of the arc
denominator, are
\[
\begin{array}{c|ccc}
 & m_s^+ & m_s^- & H_s^{\mathrm{arc}}=g+m_s^+\\ \hline
 \text{thin} &G_0&G_0&g+G_0\\
 \text{up}   &b_s+S_0&\eta_0&g+b_s+S_0\\
 \text{down} &\eta_0&b_s+S_0&g+\eta_0 .
\end{array}
\]
Here up means that the large jumps are in \([b_s,b_s+S_0]\).
For down they are in \([-b_s-S_0,-b_s]\).
The arc exponents are used in
\[
 \Lambda=\prod_{s<a}p_s^{H_s^{\mathrm{arc}}},\qquad
 \Theta=\prod_{s<a}p_s^g .
\]
The symmetric widths \(H_i\) from the recurrence for \(R\) control
the ordinary Wach cokernel.

For the evaluation at offset \(s\), use the product over all
primitive \(\ell^{s+1}\)-th roots \(\zeta\) of the space of Laurent
series in \(t\) with coefficients in the terminal crystalline
space \(D_{i,c}\). In each component expand the ordinary
cyclotomic parameter at \(1+\pi=\zeta\exp(t)\).
The scalar evaluation map is denoted by \(\rho_s\).

At every primitive root indexing \(\rho_s\) there is an expansion
\[
 \rho_s(p_s)=t\,u_s(t),\qquad u_s(0)\ne0.
\]
Indeed the primitive roots are simple
roots of the cyclotomic polynomial in characteristic zero and
\(\zeta\exp(t)\) has a nonzero derivative at zero.
The tame parameter \(X\) and all \(p_{s'}\) with \(s'\ne s\)
evaluate to units. Frobenius on the series sends
\(\zeta\) to \(\zeta^\ell\) and rescales \(t\) by \(\ell\),
and on the crystalline coefficients it transports them to the
next letter. Cyclotomic action on the degree \(k\) coefficients
is the permutation of the roots given by the cyclotomic value
times its \(k\)-th power. The level space consists of the tuples
fixed by the tame group for this action.

If \(\operatorname{Fil}^u_{c-s}\) denotes the filtration at the starting letter, the filtration at offset \(s\) is
\[
 \mathcal F_{i,s}^u=\Phi_{\rm cris}^s(\operatorname{Fil}^u_{c-s})\subseteq D_{i,c}.
\]
This transport is the de Rham comparison at that level. We label a Wach
step by its terminal Hodge letter. To translate
the de Rham convention using \(\varphi^{-(s+1)}\), label a
Wach component by the embedding one step before its
ordinary coefficient embedding and transport its crystalline
coefficients by \(\Phi_{\rm cris}^{-1}\). An evaluation at
level \(s=0\) then has the Hodge filtration of its own letter.
At level \(s\) the remaining \(s\) steps give the filtration of
\(c-s\), transported by \(\Phi_{\rm cris}^s\).
This convention is common to the source, target and extension.
The Wach evaluation generates the completed rational module
of tuples having coefficient of degree \(k\) in
\(\mathcal F_{i,s}^{-k}\). We denote this condition by
\(J_{i,s}\).

\begin{lemma}[The completed Wach lattice at a cyclotomic root]\label{lem:wach-evaluated-hodge}
Label a Wach step by its terminal Hodge letter and use
\(\Phi_{\rm cris}^{s}\) to transport the filtration from \(c-s\).
If \(\widehat N_{i,s}^{\rm rat}\) denotes the completed rational Wach lattice at cyclotomic level \(s\), then
\[
 \widehat N_{i,s}^{\rm rat}=J_{i,s}\qquad(0\le s<a),
\]
with \(J_{i,s}\) as in Paragraph~\ref{par:wach-level-models}. At the starting letter
reduction at \(X=0\), after inverting \(\ell\), gives the crystalline
space. Both comparisons commute with the coefficient-field action.
\end{lemma}

\begin{proof}
First use a common Tate twist to put the
crystalline weights in the positive range of Berger's Wach
comparison \cite[Propositions II.2.1, II.3.1 and III.2.1]{berger-wach}.
Over the ordinary cyclotomic parameter, the comparison at a
root is the de Rham embedding indexed by that root, with
\(t\) rescaled and crystalline coefficients transported by
Frobenius. If \(d_j\) is an adapted basis with jumps \(u_j\), the completed de Rham lattice is
\[
 \bigoplus_j K_O[[t]]\,t^{-u_j}d_j
\]
in a root component with coefficient field \(K_O\). Indeed these elements lie in the de Rham
lattice and their reductions form a basis by the Hodge--Tate
comparison for the de Rham representation. The completed
Wach module maps into this lattice through
\(N(V)\subset D^+(V)\) in the positive range. Relative to the
lattice generated by the constant crystalline basis, its
elementary-divisor exponents at the root are \(u_j\) by
Berger's two elementary-divisor propositions.

Let \(L\subseteq L'\) denote the inclusion at the root. Their determinant orders relative to the same constant basis are both \(\sum_j u_j\), so
\[
 \operatorname{length}(L'/L)
  =\operatorname{ord}\det L-\operatorname{ord}\det L'=0.
\] This comparison
can first be made on the underlying \(\Q_\ell\)-representation.
It commutes with the coefficient-field action and hence with
the idempotents splitting the coefficients and the embeddings.

Undoing the Tate twist shifts the jumps and the Laurent orders
by the same amount, since the change of Wach basis and of
crystalline basis differ by a power of \(t/\pi\). At a primitive root \(\zeta\ne1\) the local parameter \(t\) vanishes and \(\pi(0)=\zeta-1\ne0\). Hence
\[
 \operatorname{ord}_t(t/\pi)=1.
\]
This is the shift in Laurent order when the Tate twist is undone.

Passage to tame invariants
uses the averaging idempotent for the group of order
\(\ell-1\). Averaging the ordinary basis preserves its reduction
modulo \(\pi\), where the action is trivial. The averaged basis
therefore remains a basis over \(O[[\pi]]\), and its invariant
coefficients are in \(\mathcal R\). This yields the asserted
coefficient condition on the tame-invariant spaces. At the root one, the starting-letter identification is
\[
 N_{i,c-a}[1/\ell]/XN_{i,c-a}[1/\ell]\simeq D_{i,c-a}.
\]
It is the crystalline comparison at the origin
\cite[Corollary III.4.5]{berger-wach}. In crystalline
coordinates the ordinary expansion at one is regular, and
the factor \(t/\pi\) for the twist is a unit there.
\end{proof}

\begin{proposition}[An integral membership test along an arc]\label{prop:wach-arc-membership}
Write \(w_s\) for the level-\(s\) evaluation and put
\[
 W_i=\Lambda^{-1}\operatorname{span}_{\mathcal R}\Phi^a N_{i,c-a}.
\]
Using \(m_s^\pm\) and \(\Lambda\) from
Paragraph~\ref{par:wach-level-models}, one has
\[
 \Bigl(\prod_{s<a}p_s^{m_s^+}\Bigr)N_{i,c}
    \subset\operatorname{span}_{\mathcal R}\Phi^aN_{i,c-a},
 \qquad
 w\in N_{i,c}\ \Longleftrightarrow\
 (\forall s<a)\ w_s\in J_{i,s}\quad(w\in W_i).
\]
The signed step bound applies also to the crystalline extension of
\(O\) by \(T_2\).
\end{proposition}

\begin{proof}
\proofstep{Clearing the one-step denominators.}
These comparisons give integral step and level tests. For
one step ending at \(i\), the span \(\mathcal R\Phi N_{i-1}\)
evaluates to the regular lattice at the first cyclotomic
level. This follows by transport of the comparison at one.
The condition at that level on \(N_i\) is the Hodge filtration
at the ending letter. At the first level the two factors move in opposite directions:
\[
 q^{m^+}N_i\ \subseteq\ \mathcal R\Phi N_{i-1},\qquad q^{m^-}\mathcal R\Phi N_{i-1}\ \subseteq\ N_i,
\]
initially in the evaluated modules. For integral inclusions use the coarse Wach height bounds.
After a common twist the ordinary positive Wach module is
stable under Frobenius, with linearized cokernel killed by
the cyclotomic factor to the power of the width of the
jump box \cite[Definition III.4.1 and Proposition III.4.2]{berger-wach}.
Undoing the twist multiplies the Frobenius matrix by a
power of that factor since \(\phi(\pi)/\pi\) is the ordinary
factor.

Tame averaging gives the corresponding inclusions
on invariant modules. Thus the coordinates on either basis
have an a priori denominator a power of \(q\). The
local test removes that denominator integrally. In fact the
remainder on division by the distinguished associate of
\(q^m\) vanishes when its Taylor coefficients of order
less than \(m\) vanish at the roots. Indeed let \(D\) be the distinguished polynomial and \(R\) the remainder. Over \(K_O\) the Taylor conditions imply
\[
 D\mid R,\qquad \deg R<\deg D\quad\Longrightarrow\quad R=0.
\]
The quotient by \(D\) is free over \(O\), so this vanishing descends integrally. Iterating the resulting signed bound gives
\[
 \Bigl(\prod_{s<a}p_s^{m_s^+}\Bigr)N_{i,c}
    \subset\operatorname{span}_{\mathcal R}\Phi^aN_{i,c-a}.
\]
The same proof applies to the crystalline extension
\(0\to T_2\to T'\to O\to0\). Exactness of the crystalline filtration
adds only the jump at zero, which lies in each of the stated boxes.

\proofstep{Testing the product of the denominators.}
Put
\[
 W_i=\Lambda^{-1}\operatorname{span}_{\mathcal R} \Phi^aN_{i,c-a}.
\]
Using the opposite step bound, there is a product
\(Q_*\) of powers of the factors \(p_s\) with \(Q_*W_i\subset N_{i,c}\).
For \(w\in W_i\), expand \(Q_*w\) on a Wach basis at \(c\).
At each level the Wach basis is a basis of the completed Hodge
lattice at each root. Thus the evaluation of \(w\) is in \(J_{i,s}\)
exactly when its multiplied coordinates for \(Q_*w\) vanish to the
order of \(\rho_s(Q_*)\) at every root of level \(s\).

For the scalar
remainder test over \(O[[X]]\), the values of \(X\) at primitive
roots identify the distinct tame orbits, with one simple root of the
distinguished associate of \(p_s\) per orbit. Weierstrass division
and the same root test therefore give divisibility by the entire
product \(Q_*\) in every coordinate. Consequently
\[
       w\in N_{i,c}\quad\Longleftrightarrow\quad
       w_s\in J_{i,s}\text{ for every }s<a .
\]
The units at each level give the two tests used below. The inverse
image of \(N_{i,c}\) under \(\Theta\) lies in \(W_i\) by the signed
step bound. In \(W_i\) its test is \(p_s^g w_s\in J_{i,s}\)
for every \(s\). Also \(XN_{i,c}\) consists of the vectors of
\(XW_i\) having evaluations in every \(J_{i,s}\); apply the
previous test after dividing by \(X\), a unit at the evaluated
levels.
\end{proof}

Define \(D_{i,s}\) to be the diagonal crystalline constants.
For an up letter write \(U^K_{i,s}=\mathcal F_{i,s}^{b_s}\),
and for down write \(U^K_{i,s}=\mathcal F_{i,s}^{-\eta_0}\).
Define \(B_{i,s}\) by the following increasing degree filtration
on its coefficients. In a thin letter it is zero below degree
zero and full from zero. In an up letter the breaks are \(-b_s\)
and zero. In a down letter they are zero and \(b_s\). In the
last two cases it is zero before the first break, \(U^K_{i,s}\)
from the first to the second break, and full from the second.
Each coefficient condition is imposed at all the roots.

\begin{lemma}[Comparison with the rounded degree filtration]\label{lem:wach-rounded-breaks}
For the coefficient spaces \(B_{i,s}\) with the rounded breaks
just defined, one has
\[
 p_s^gB_{i,s}\subset J_{i,s},\qquad p_s^gJ_{i,s}\subset B_{i,s}.
\]
In coordinates the rounded filtration gives the divisibility conditions
of \eqref{eq:wach-level-condition}, compatibly with the crystalline map.
\end{lemma}

\begin{proof}
Between the two jump windows the Hodge filtration is constant,
with middle subspace precisely \(U^K_{i,s}\). Its two groups
of degrees differ from the rounded breaks by at most \(g\)
since \(G_0,S_0,\eta_0\le g\). In the thin case the one
group differs from zero by at most \(G_0\). It follows,
using multiplication by the unit part of \(\rho_s(p_s)\) on
the increasing filtrations, that
\[
       p_s^g B_{i,s}\subset J_{i,s},\qquad
       p_s^g J_{i,s}\subset B_{i,s}.
\]

Fix a Wach basis \(e_i\) at the start and write \(\lambda_i\)
for its reduction at one transported by \(\Phi^a\). This is
a basis of \(D_{i,c}\). Viewing the basis as a row, the coordinate evaluation at offset \(s\) is
\[
 y\longmapsto\rho_s(\Lambda)^{-1}\lambda_i\rho_s(y),\qquad y\in\mathcal R^{r_i}.
\]
The saturated inverse image \(U_{i,s}\) under \(\lambda_i\) has an integral complement \(C_{i,s}\). Thus
\[
 O^{r_i}=U_{i,s}\oplus C_{i,s},\qquad U_{i,s}=\lambda_i^{-1}(U^K_{i,s})\cap O^{r_i}.
\]
Multiplying
the rounded conditions by \(\rho_s(\Lambda)\) shifts the breaks
by \(H_s^{\mathrm{arc}}\). In bases of the saturated submodule
and its complement, Weierstrass division gives exactly the
three divisibility tests in \eqref{eq:wach-level-condition}.

For a constant plus this lattice, subtract the series
\(\Lambda d\) for the corresponding crystalline constant
in \(K_O^{r_i}\) and clear its finitely many coefficient
denominators. This gives the test with \(\varpi^k y-\Lambda d\)
in that paragraph. The crystalline map on the transported
bases has the matrix of the Wach map reduced modulo \(X\), because
reduction, crystalline comparison and Frobenius are functorial.
Thus both evaluation maps commute with the specified maps.
For the crystalline map \(F\), compatibility of the up submodules comes from
\[
 F(\mathcal F_{1,s}^{b_s})\subseteq\mathcal F_{2,s}^{b_s}.
\]

\end{proof}

\paragraph{One integral partner at all levels of an arc.}\label{par:wach-partners}
The truncation modulus after the arc is
\[
 D'=\phi^{a-j}\left(X^M\prod_{s<j}p_s^{3g}\right).
\]
Its orders at the evaluated levels are
\[
 B_s^*=\begin{cases}
 (\ell-1)M,&s<a-j,\\
 3g,&a-j\le s<a .
 \end{cases}
\]

\begin{lemma}[One integral representative for all levels of an arc]\label{lem:wach-arc-partners}
For the cut and arc of the construction, assume \((\ell-1)M>h+2g\),
\(3g\le(\ell-1)M\), and the stated buffer of thin letters.
The comparison uses the filtration
\[
 F_{i,s}=\{z:\operatorname{ord}(z)\ge B_s^*-H_s^{\rm arc}\}.
\] After the scalar loss \(\ell^\sigma\)
of the cut congruences, a vector in either the Wach or the coordinate
description has one integral partner in the other description
agreeing modulo \(F_{i,s}\) at every \(s<a\).
\end{lemma}

\begin{proof}
\proofstep{Orders of the truncation factors.}
The polynomial \(\phi^{a-j}(X)\) has a simple \(X\)-factor and
\(p_s\)-orders \(\ell-1\) for \(s<a-j\). Define \(F_{i,s}\)
to consist of tuples of order at least \(B_s^*-H_s^{\mathrm{arc}}\).
At each of the last \(j\) backward offsets the letter is thin
by the buffer after the starting cut. At thin offsets the order
is at least \(g\), and at up offsets it is at least \(g+1\).

At down offsets it is at least \(b_s+S_0+1\).
Here we use
\[
 (\ell-1)M>h+2g
\]
and \(b_s+S_0\le h\).
All these degrees are in the full-coefficient part of
both filtrations. Thus \(F_{i,s}\) is contained in
\(J_{i,s}\cap B_{i,s}\). The same order is used for the two
representations and the crystalline map preserves it.
If a coordinate evaluation is in \(F_{i,s}\), its
numerator is divisible by \(p_s^{B_s^*}\) by the scalar
remainder test, in particular by \(p_s^{H_s^{\mathrm{arc}}}\).

\proofstep{A basis fixed modulo the cut.}
There is a basis change \(A_i(X)\) modulo \(X^M\), with
constant term the identity, making \(e_i A_i\)
\(\gamma\)-invariant modulo \(X^M\) over \(K_O\).
Construct it degree by degree, dividing at degree \(r\) by
\(\omega_X^r-1\), just as in the cut correction.
Write
\[
 A_i(X)=1+\sum_{r=1}^{M-1}A_{i,r}X^r.
\]

If \(C_{i,r}\) is the degree-\(r\) coefficient of the remaining
invariance defect of the previous truncation in the basis \(e_i\),
the next correction solves
\[
  (\omega_X^r-1)A_{i,r}=-C_{i,r},\qquad
  A_{i,r}\in\ell^{-a_r}\operatorname{Mat}(O),\qquad
  a_r=\sum_{k=1}^r(1+v_\ell(k)).
\]
For the inverse, every term in degree \(r\) of the geometric series
has the form \(\pm A_{i,r_1}\cdots A_{i,r_t}\), where
\(r_1+\cdots+r_t=r\). For a product of coefficients with degrees \(r_1+\cdots+r_t=r\), the extra valuation is controlled by
\[
 \sum_j v_\ell(r_j!)\le v_\ell(r!),\qquad \frac{r!}{\prod_j r_j!}\in\mathbb Z.
\]

In subsequent bounds write \(A_r=a_r\). Use polynomial representatives of the two
truncated matrices. At the root one the ordinary parameter
satisfies \(X\sim t^{\ell-1}\). In crystalline coordinates the
normalized basis at one is constant modulo \(t^{(\ell-1)M}\):
if \(u_k\) is a coefficient in an intermediate degree, cyclotomic invariance gives
\[
 (\omega^k-1)u_k=0,\qquad 0<k<(\ell-1)M ,
\]
and \(\omega^k\ne1\) forces \(u_k=0\). Every evaluated root reaches one after
\(a\) Frobenius steps, which transport this constant to
\(\lambda_i\) and preserve the orders of the nonconstant terms.

\proofstep{The two integral representatives.}
Applying \(\phi^{a-j}\) to the scalar cut congruences shows that
\(\ell^\sigma\phi^a(A_i)\) and
\(\ell^\sigma\phi^a(A_i^{-1})\) are congruent to integral
matrices modulo \(D'\) over \(\mathcal R[1/\ell]\).
Here each coefficient uses the bound with \(A_r\) at its own
degree. Also \(\phi^a(X)^M\) is divisible by \(D'\)
by the upper bound
\[
 3g\le(\ell-1)M.
\]
Consequently
the product of the two transported matrices is the identity
modulo \(D'\) before multiplying by the scalar.
Write an arc vector as \(\Lambda^{-1}(\Phi^a e_i)v\).
Its coordinate partner is obtained by applying an integral
representative of \(\ell^\sigma\phi^a(A_i^{-1})\) to \(v\).

Conversely, the arc partner of a coordinate vector \(y\)
uses \(v\) obtained by applying an integral representative
of \(\ell^\sigma\phi^a(A_i)\) to \(y\).
In either direction the transported invariant basis gives
the required agreement after multiplying by \(\Lambda\), modulo
\(t^{B_s^*}\) at every level. Dividing by \(\Lambda\) gives
the congruence modulo \(F_{i,s}\). Each partner uses a single
integral vector for the whole arc.
\end{proof}

\paragraph{The two cyclotomic polynomial actions.}\label{par:wach-polynomial-actions}
Let \(\omega=\chi(\gamma)\) and let \(\tau_s\) be the permutation
of the primitive roots at level \(s\) given by \(\omega\).
On degree \(k\), \(\gamma\) acts by \(\omega^k\tau_s\).
The principal units act through a cyclic group of order
\(\ell^s\), generated by \(\tau_s\), and together with
the tame units act transitively on the roots.

\begin{lemma}[Polynomial corrections at all evaluated levels]\label{lem:wach-level-polynomials}
On the tame-invariant space,
\[
 (\gamma-1)w_s\in J_{i,s}\quad\Longrightarrow\quad w_s\in D_{i,s}+J_{i,s}.
\] For \(E=a-1\) the polynomials
\(P_{E,I}\) constructed from the respective target and source
degree intervals below send \(p_s^{-g}B_{2,s}\) into
\(D_{2,s}+B_{2,s}\) and \(p_s^{-g}J_{1,s}\) into \(D_{1,s}+J_{1,s}\)
at all offsets. The degrees used lie in at most three intervals
of length \(3g\) inside the prescribed range \([-2h,2h]\).
\end{lemma}

\begin{proof}
Suppose \((\gamma-1)w_s\in J_{i,s}\). In degree \(k\ne0\),
let \(\bar w_{s,k}\) be the fixed class in the degree-\(k\)
quotient. The permutation of the roots has order \(\ell^s\), so
\[
   \bar w_{s,k}=\gamma^{\ell^s}\bar w_{s,k}
               =\omega^{k\ell^s}\bar w_{s,k}.
\]
For \(k\ne0\), \(\omega^{k\ell^s}\ne1\), hence
\(\bar w_{s,k}=0\) over the coefficient field. In degree zero the class is both tame and
principal invariant, thus constant across the roots. A lift
of that constant in \(D_{i,c}\) shows \(w_s\in D_{i,s}+J_{i,s}\).
The action fixes diagonal constants and preserves \(J_{i,s}\).

\proofstep{Averaging the degree-zero coefficient.}
Use the polynomials \(P_{E,I}\) of the arc calculation with
\(E=a-1\), so \(E\ge s\) at each level.
A factor \(Z^{\ell^E}-\omega^{k\ell^E}\) kills degree \(k\).
In degree zero the sum \(\sum_{v=0}^{\ell^E-1}Z^v\) averages
over the principal orbit with a scalar multiplicity.
For a tame-invariant degree-zero tuple \(b\) on the roots, its
\(\gamma\)-orbit has period dividing \(\ell^s\). Hence
\[
 \sum_{v=0}^{\ell^E-1}\gamma^v b
   =\ell^{E-s}\sum_{j=0}^{\ell^s-1}\gamma^j b
   \qquad(E\ge s).
\]
The sum is fixed by \(\gamma\) and by the tame action. Their transitivity
makes it diagonal constant, a space preserved by the remaining factors.

\proofstep{The rounded target and the filtered source.}
For the target polynomial, write \(\mathcal B_s\) for the rounded breaks defining \(B_{2,s}\). Take precisely the integer degrees
\[
 I=\bigcup_{e\in\mathcal B_s}\bigl([e-g,e)\cap\mathbb Z\bigr).
\]

A vector in \(p_s^{-g}B_{2,s}\) satisfies in degree \(k\)
the rounded condition for \(k+g\).
When this condition is larger than that for \(k\), one of
the selected intervals contains \(k\). Every such \(k\)
is nonzero, since \(b_s>g\) on positions with larger jumps.
The polynomial kills these degrees, and in the other degrees
it preserves the required subspace. It fixes the space of
constants up to scalar, giving the required map to
\(D_{2,s}+B_{2,s}\).

On the source side use the separate set
\[
 I_{\rm src}=\bigcup_{u}\bigl([-u-g,-u)\cap\mathbb Z\bigr),
\]
where \(u\) runs over the jumps in the prescribed source boxes.
The degree test for \(p_s^{-g}J_{1,s}\) compares
\(\mathcal F_{1,s}^{-k-g}\) with \(\mathcal F_{1,s}^{-k}\).
Every change occurs in a selected interval. The factors kill
all offending nonzero degrees. In degree zero the principal
sum gives a diagonal constant, which is allowed in
\(D_{1,s}+J_{1,s}\). At degrees where the subspace is
unchanged it is preserved. This proves the second action
on the same evaluation model.

For each polynomial the union lies in three integer intervals
of length at most \(3g\). After enlarging \(G_0\) to include \(\eta_0\), the central interval and its length satisfy
\[
 I_{\rm cent}\subseteq[-G_0-g,G_0-1],\qquad |I_{\rm cent}|\le2G_0+g\le3g.
\]
The intervals around a
thick up or down break lie within the corresponding one-sided
translated window of length at most \(g+S_0\).
The arc contains at most two positions with larger jumps. All degrees lie
in \([-2h,2h]\). These are precisely the window and order
bounds used for the scalar valuation at one. The powers of
\(\ell\) in the partners and polynomial values use only the
maximum length of an arc and the prescribed boxes.
\end{proof}


\paragraph{A level comparison for the mixed-sign arc.}
\label{par:wach-mixed-constant}
Use the arc length \(a\le A\), the powers \(H_s^{\mathrm{arc}}\),
and the level data of Paragraph~\ref{par:wach-level-models}.
Distinguish the transported ranks \(r_i\), the coordinate map
\(L:O^{r_1}\to O^{r_2}\) on the starting bases modulo \(X\),
and its coefficientwise extension \(L_{\mathcal R}\).
The polynomial comparison uses, at each \(s\), evaluation
maps from \(W_i\) and from \(\mathcal R^{r_i}\) to a common level
space. Denote their values by \(w_s\) and \(y_s\), respectively.
The latter refers to a vector \(y\) of coordinate series, not to
its coefficients at \(X=0\). The following are the lattice
properties used in this comparison.

\begin{itemize}
\item There are two local lattices \(B_{i,s},J_{i,s}\) with
\(p_s^g B_{i,s}\subset J_{i,s}\) and
\[
 p_s^g J_{i,s}\subset B_{i,s}.
\]
A submodule \(F_{i,s}\) lies
in both. Every vector on either side of the evaluation has a partner
on the other side, whose evaluation agrees with \(\ell^\sigma\)
times the original modulo \(F_{i,s}\), simultaneously at all
levels. The two evaluation maps commute with the corresponding
maps induced by \(f\) and \(L_{\mathcal R}\). The level map
preserves the \(F_{i,s}\).
\item The preimage of \(N_{i,c}\) under multiplication by
\(\Theta\) lies in \(W_i\). On \(W_i\) this preimage is tested
by
\[
 p_s^g w_s\in J_{i,s}
\]
for all \(s\).
The map \(f\) lifts on this preimage with cost \(\ell^R\).
If \(y_s\in F_{2,s}\), then \(p_s^{H_s^{\mathrm{arc}}}\) divides
\(\varpi^k y\) for some \(k\), where \(\varpi\) is a
uniformizer of \(O\).
\item There are saturated constant submodules \(U_{i,s}\) of
\(O^{r_i}\) and subgroups \(D_{i,s}\) of crystalline constants
in the level spaces. Membership \(y_s\in B_{i,s}\) is equivalent,
over all levels, to the conditions
\[
 y\in
 \begin{cases}
 p_s^{H_s^{\mathrm{arc}}}\mathcal R^{r_i} &\text{smaller jumps},\\
 p_s^{H_s^{\mathrm{arc}}-b_s}\mathcal R U_{i,s}+p_s^{H_s^{\mathrm{arc}}}\mathcal R^{r_i}
                          &\text{up},\\
 p_s^{H_s^{\mathrm{arc}}}\mathcal R U_{i,s}+p_s^{H_s^{\mathrm{arc}}+b_s}\mathcal R^{r_i}
                          &\text{down},
 \end{cases}
 \tag{\(*\)}\label{eq:wach-level-condition}
\]
after allowing multiplication of \(y\) by a power of \(\varpi\)
at each level. Membership \(y_s\in D_{i,s}+B_{i,s}\) is tested
by the same conditions on \(\varpi^k y-\Lambda d\) at each level,
for some \(k\ge0\), \(d\in O^{r_i}\). Write \(\mathcal N_i\)
and \(\mathcal I_i\) for these two modules of coordinate series.
\end{itemize}

\begin{lemma}[The up subspace after a down letter]\label{lem:up-after-down}
Suppose the two positions with larger jumps are a down letter
at \(s_0\) and an up letter at \(s_1>s_0\).
Assume \(L(U_{1,s_1})\subset U_{2,s_1}\),
\(2b_{s_0}\le\ell b_{s_1}\) and \(2a_0<b_{s_1}\), where
\(a_0=2\sigma+R+2gA\) is the initial exponent of
Section~\ref{sec:inert-long-orbit} and \(a\le A\) is the arc length. Then
\[
 \ell^{2a_0}U_{2,s_1}\subset L(U_{1,s_1}).
\]
\end{lemma}

\begin{proof}
Put \(Q=p_{s_1}^{b_{s_1}}\), \(P=p_{s_0}^{b_{s_0}}\) and
\(\Lambda_1=\Lambda/Q\). If \(v\in U_{2,s_1}\), the series
\(y_2=P\Lambda_1 v\) belongs to \(\mathcal N_2\), directly
by \eqref{eq:wach-level-condition} at each of the two positions with larger jumps and at every position with smaller jumps.

Choose its partner \(w_2\in W_2\). The error at each level lies
in \(F_{2,s}\) and \((y_2)_s\in B_{2,s}\), so
\[
 p_s^g(w_2)_s\in J_{2,s}.
\]
The preimage test gives a lift
\(w_1\) on the \(\Theta\)-divisible lattice, with
\(f(w_1)=\ell^R w_2\). Choose the coordinate partner \(y_1\)
of \(w_1\). Each \(p_s^{2g}(y_1)_s\) belongs to \(B_{1,s}\),
so \(\Theta^2 y_1\in\mathcal N_1\).
The partner errors and compatibility with \(f\) give
\[
 (L_{\mathcal R}y_1-\ell^{2\sigma+R}y_2)_s\in F_{2,s}.
\]

Every \(p_s^{H_s^{\mathrm{arc}}}\) therefore divides the difference after
multiplication by a power of \(\varpi\). These divisibilities
hold integrally: each quotient by a distinguished polynomial
power is free over \(O\) by Weierstrass division. Their product
\(\Lambda\) divides by relative primality after inverting
\(\varpi\), followed by the same saturation for the product.
This uses no integral Chinese remainder decomposition.

\proofstep{Testing cyclic quotients of the image.}
Take an integral linear form \(\lambda\) on \(O^{r_2}\)
such that \(\lambda(v)=1\) and
\(\lambda L(U_{1,s_1})\subset tO\).
For any \(y\in\mathcal N_1\),
\[
                 \lambda L_{\mathcal R}(y)
                   \in \Lambda_1(t,Q).
\]
Indeed split \(O^{r_1}\) along the saturated \(U_{1,s_1}\).
On its complement the up condition implies divisibility by
\(p_{s_1}^{H_{s_1}}\), and on the summand \(U_{1,s_1}\) it
implies divisibility by \(p_{s_1}^{H_{s_1}-b_{s_1}}\).
All coordinates are divisible by \(p_s^{H_s^{\mathrm{arc}}}\) at the other
levels. Multiplying the powers of the distinct cyclotomic factors, the complementary
summand gives a multiple of \(\Lambda\) and the up summand gives
a multiple of \(t\Lambda_1\). Applying this to
\(\Theta^2 y_1\) and using the congruence for
\(L_{\mathcal R}y_1\) yields, after cancelling \(\Lambda_1\),
\begin{equation}\label{eq:wach-gauss-ideal}
                 \ell^{2\sigma+R}\Theta^2P\in(t,Q).
\end{equation}

\proofstep{The Gauss valuation and Smith factors.}
Normalize \(v_\varpi(\varpi)=1\), and put \(e_O=v_\varpi(\ell)\).
For \(x>0\) the Gauss valuation of \(\sum a_m X^m\) is
\(\min_m(v_\varpi(a_m)+mx)\). If \(i,j\) are the least minimizing indices for \(f,g\), then in
the coefficient of \(X^{i+j}\) only \(f_i g_j\) attains the minimum.
Indeed a different pair with the same index sum uses an index below
\(i\) or below \(j\). Therefore
\[
 v_\varpi([X^{i+j}]fg)+(i+j)x=v_x(f)+v_x(g)=v_x(fg).
\] Preparation gives
\(v_x(p_s)=\min(e_O,\ell^s x)\).

If \(t\ne0\), set
\(k=v_\varpi(t)\); for \(t=0\) take \(k=+\infty\).
The assertion for \(k=0\) is immediate. Otherwise set
\(b_1=b_{s_1}\), \(b_0=b_{s_0}\),
\(m=\min(k,e_O b_1)\), and \(x=m/(b_1\ell^{s_1})\). Then
\(v_x(Q)=m\), and \(v_x(P)\le b_0 m/(\ell b_1)\).
The ideal relation and \(v_x(p_s)\le e_O\) give
\[
 m\le e_O(2\sigma+R+2ga)+\frac{b_0}{\ell b_1}m
   \le e_O a_0+m/2 .
\]
Thus \(m\le2e_Oa_0<e_Ob_1\), which excludes \(t=0\) and
proves \(t\mid\ell^{2a_0}\).

Since \(U_{2,s_1}\) is a direct
summand, Smith coordinate forms extend to the ambient
\(O^{r_2}\). Apply the divisibility assertion to its Smith
basis vectors for the inclusion of \(L(U_{1,s_1})\).
A missing image direction would give \(t=0\); the invariant
factors in the image directions all divide \(\ell^{2a_0}\).
Consequently
\[
                   \ell^{2a_0}U_{2,s_1}\subset L(U_{1,s_1}).
\]
The use of \(s_0<s_1\) gives the factor \(1/\ell\) in the
valuation comparison. The argument measures divisibility by a
scalar in a ramified coefficient ring and so is uniform under
the coefficient extension.
\end{proof}

\paragraph{From constant coordinates to the arc lattice.}
\begin{lemma}[Lifting constant coordinates across an arc]\label{lem:wach-constant-lift}
With the integral lattices and scalar \(a_0\) of this comparison,
each \(y\in\mathcal I_2\) has a lift modulo \(\mathcal N_2\):
\begin{equation}\label{eq:wach-const-coord-lift}
 \ell^{2a_0}y=L_{\mathcal R}y_1+\ell^{2a_0}\nu,\qquad
 y\in\mathcal I_2,\quad y_1\in\mathcal I_1,\quad
 \nu\in\mathcal N_2 .
\end{equation}
\end{lemma}
\begin{proof}
Here is the constant comparison for all the words. For either
representation, omit its index and let \(V=O^r\). If there is
no down letter, the level conditions are unchanged by adding a
local constant times \(\Lambda\), so \(\mathcal I=\mathcal N\).
This includes a word with two up letters. Otherwise there is at
most one up letter. Write \(Q=p_s^{b_s}\) at the up letter and
\(Q=1\) if there is none. Divisibility at every level gives
\(y=(\Lambda/Q)Y\), with \(Y\in\mathcal R^r\).
At an up letter the condition for \(Y\) is
\[
             Y\in\mathcal R U_s+Q\mathcal R^r .
\]

At a down letter with \(P=p_s^{b_s}\) it is the condition that
the image of \(Y\) in \((\mathcal R/(P))\otimes_O(V/U_s)\)
lie in the rational constant line factor \(\overline Q\) times
\((V/U_s)\otimes_O\operatorname{Frac}(O)\). To see this,
divide by the full local power \(p_s^{H_s^{\mathrm{arc}}}\), then cancel the
other level factors modulo \(P\) after inverting \(\varpi\).
The difference between \(\mathcal I\) and \(\mathcal N\) is
that \(\mathcal N\) requires the image at each down letter to
be zero. The quotients in this test are free over \(O\) by
Weierstrass division and saturation of \(U_s\).

With a single down letter and no up letter, primitivity means the coefficient map splits the constant inclusion:
\[
 O\xrightarrow{\ a\mapsto a\ }\mathcal R/(P)\xrightarrow{\lambda}O,\qquad \lambda(1)=1.
\]
In fact, if the projected class is \(1\otimes v\), choose an integral
linear form \(\lambda\) on \(\mathcal R/(P)\) with \(\lambda(1)=1\).
Then
\[
 (\lambda\otimes1)(1\otimes v)=v\in V/U_s .
\]
This recovers the constant integrally in the saturated quotient. If the up letter
has smaller offset than the down one, its power \(Q\) reduces
modulo \(\varpi\) to a unit times a monomial of smaller degree
than the distinguished polynomial for \(P\), by the inequality
\(b_{\rm up}<\ell b_{\rm down}\). The class of \(Q\) modulo
\(P\) is again primitive, which proves integrality of the
constant by the same linear-form test. In both cases lift the
constant to \(d\in V\). Then
\[
 y-\Lambda d\in\mathcal N.
\]


\proofstep{Two down letters.}
For two down letters let \(P,Q\) be their powers in increasing
offset order, and \(U_0,U_1\) the corresponding submodules.
There are integral constants \(d_0,d_1\in V\) by the separate
primitive test for \(1\) in each quotient. Their difference
belongs to the ideal \((P,Q)\) in
\(\mathcal R\otimes_O V/(U_0+U_1)\). Let \(\rho\) be least with
\[
 \varpi^\rho\in(P,Q),
\]
which exists by coprimality after
inverting \(\varpi\). The constant components in a cyclic
decomposition of \(V/(U_0+U_1)\) belong to
\(\varpi^\rho V/(U_0+U_1)\). On a free summand this follows
from the definition of \(\rho\).

On a summand of order
\(\varpi^k\), if an integral representative \(a\) in
\((P,Q,\varpi^k)\) had valuation below \(\min(k,\rho)\), then
\(a-\varpi^k f\in(P,Q)\) for some \(f\in\mathcal R\).
The factor \(1-(\varpi^k/a)f\) is a unit, a contradiction
to the choice of \(\rho\).
Write
\[
 d_0-d_1=u_0+u_1+\varpi^\rho d'
\]
with
\(u_i\in U_i\). Choose \(f_\rho\in(Q)\) with
\(f_\rho-\varpi^\rho\in(P)\). Put \(Y'=Y-d_1-u_1-f_\rho d'\).
The two reductions are
\[
 \begin{aligned}
  Y'&\equiv d_0-d_1-u_1-\varpi^\rho d'=u_0
        &&\pmod {P\mathcal R^r+\mathcal R U_0},\\
  Y'&\equiv d_1-d_1-u_1=-u_1
        &&\pmod {Q\mathcal R^r+\mathcal R U_1}.
 \end{aligned}
\]
Both images vanish. Multiplying back by \(\Lambda\) gives the
two factors \(\Lambda\) and \(\Lambda f_\rho\).

\proofstep{A down letter before an up letter.}
For a down letter at the smaller offset and an up letter at
the larger one, write \(U=U_{s_1}\) and \(V_0=U_{s_0}\),
with \(P,Q,\Lambda_1\) as in the mixed-sign estimate. Write
\(Q=\varpi^k\alpha+P\beta\), with \(k\) largest possible.
There is such a finite \(k\) because the image of \(Q\) in
the finite free \(\mathcal R/(P)\) is nonzero. The class
of \(\alpha\) is primitive. The down condition for \(Y\)
therefore gives
\[
 \overline Y=\overline\alpha d_0
\]
in
\((\mathcal R/(P))\otimes_O(V/V_0)\) for some \(d_0\in V/V_0\).
The up condition, projected further to \(V/(U+V_0)\),
shows that \(\overline\alpha d_0\) belongs to
\(\varpi^k\overline\alpha((\mathcal R/(P))\otimes_O V/(U+V_0))\).

A linear form taking value one on \(\overline\alpha\)
gives \(d_0=v+\varpi^k d\) modulo \(V_0\), with
\(v\in U,d\in V\). This test works on the tensor product
with \(V/(U+V_0)\) even if that module has torsion. Now
\[
 \overline{Y-\alpha v-Qd}
    = \overline\alpha(d_0-v-\varpi^k d)=0
     \quad\text{in }(\mathcal R/(P))\otimes_O(V/V_0).
\]
The same difference still lies in the required up lattice:
\(\alpha v\in\mathcal R U\) and \(Qd\in Q\mathcal R^r\).
Upon multiplying back by the coordinate factors, the resulting
normal form uses \(\Lambda\) times an ordinary constant and
\(\alpha\Lambda_1\) times a vector of \(U\).
In each of the cases the displayed representatives themselves
satisfy the \(\mathcal I\) conditions, by reducing at the same
levels.

On those normal forms, the ordinary constants lift using the
cost \(\ell^R\) on the starting constant map and \(R\le2a_0\).
In the mixed case the additional constant in \(U_{2,s_1}\)
lifts by the Gauss valuation and Smith argument. This gives
\eqref{eq:wach-const-coord-lift}.
\end{proof}

\begin{lemma}[Integral lift between the two polynomial actions]
Let \(x'\) be a target arc vector evaluating in \(D_{2,s}+B_{2,s}\)
at each level of the cut. There is one source vector \(z\) for the arc map \(f\) such that
\[
 f(z)=\ell^{2\sigma+2a_0+R}x',\qquad z_s\in D_{1,s}+p_s^{-g}J_{1,s}\quad\text{at every level }s,
\]
where \(z_s\) is its evaluation.
\end{lemma}
\begin{proof}
Choose the coordinate partner \(y\) of \(x'\). Its error belongs
to \(F_{2,s}\subset B_{2,s}\), so \(y\in\mathcal I_2\).
Apply \eqref{eq:wach-const-coord-lift} to \(y\) and choose the
source arc partner \(x_1\) of \(y_1\). Put
\(\Delta=\ell^{2\sigma+2a_0}x'-f(x_1)\).
In its evaluation the target partner error has coefficient
\(-\ell^{\sigma+2a_0}\), the mapped source partner error has
coefficient \(-1\), and the remaining term is
\(\ell^{\sigma+2a_0}\nu_s\). All belong to \(B_{2,s}\);
the preimage test therefore puts \(\Delta\) in the
\(\Theta\)-divisible lattice. Lift \(\ell^R\Delta\) there to
\(\Delta_1\). Then
\[
 f(\ell^Rx_1+\Delta_1)
 =\ell^R(f(x_1)+\Delta)=\ell^{2\sigma+2a_0+R}x'.
\]
The source partner congruence and
\(B_{1,s}\subset p_s^{-g}J_{1,s}\) give the stated evaluation
at every level, with the same exponent.
\end{proof}
\begin{lemma}[The lift on invariant classes]
The invariant first-coordinate test used in the cut comparison is
\[
 (\gamma-1)z\in XN.
\]
The source and target polynomials in \(\gamma\) and the intervening lattice lift preserve this test in the sense of that comparison.
\end{lemma}
\begin{proof}
To apply this lift, use the level properties of
Paragraphs~\ref{par:wach-level-models} and~\ref{par:wach-polynomial-actions}. The level actions commute with the
evaluation of arc vectors. If the difference under \(\gamma-1\)
is in \(J_{i,s}\) at a level, the vector there lies in
\(D_{i,s}+J_{i,s}\). Writing \(P_2,P_1\) for the target and source polynomials, the actions at a level are
\[
 \begin{aligned}P_2(\gamma)(D_{2,s}+p_s^{-g}B_{2,s})&\subseteq D_{2,s}+B_{2,s},\\ P_1(\gamma)(D_{1,s}+p_s^{-g}J_{1,s})&\subseteq D_{1,s}+J_{1,s}.\end{aligned}
\]
The action fixes \(D_{i,s}\) and preserves
\(J_{i,s}\). Finally \(XN_{i,c}\) consists of the vectors in
\(XW_i\) whose evaluations belong to all the \(J_{i,s}\).
These are statements on the two level models at the same cut.

The action on \(W_i/XW_i\) is trivial. To see this, use the
trivial action on the starting Wach lattice modulo \(X\) and
the commutation with \(\Phi\), together with
\(\phi^a(X)\in(X)\). Division by \(\Lambda\) preserves the trivial action at the origin, since
\[
 \gamma(\Lambda)/\Lambda\equiv1\pmod X.
\]
For an arc vector \(w\) with
\((\gamma-1)w\in XN_{2,c}\), the first level property gives
evaluations in \(D_{2,s}+J_{2,s}\).

Since
\(J_{2,s}\subset p_s^{-g}B_{2,s}\), the target polynomial
produces the input \(x'\) of the intervening lattice lift.
The output, after applying the source polynomial, has
evaluations in \(D_{1,s}+J_{1,s}\). Its \(\gamma-1\) image
therefore has evaluations in all \(J_{1,s}\) and belongs to
\(XW_1\). The last level property puts the image in
\(XN_{1,c}\). An operator polynomial on a vector fixed
modulo \(XN_{2,c}\) agrees with its value at one modulo that
submodule. Writing \(P_2,P_1\) for the target and source
polynomials and \(c_{\rm lift}=2\sigma+2a_0+R\) for the intervening exponent,
choose \(y\) with
\[
 f(y)=\ell^{c_{\rm lift}} P_2(\gamma)w.
\]
Then
\[
 f(P_1(\gamma)y)
 =\ell^{c_{\rm lift}} P_1(\gamma)P_2(\gamma)w
 \equiv\ell^{c_{\rm lift}} P_1(1)P_2(1)w\pmod{XN_{2,c}} .
\]

\end{proof}
\paragraph{The colength at the cuts of the long orbit.}\label{par:wach-cut-colength}

\begin{lemma}[The scalar annihilator at a cut]\label{lem:wach-cut-annihilator}
Let \(f\) be the integral crystalline surjection on the local
arc of Paragraph~\ref{par:wach-level-models}. For \(R\) defined
in Paragraph~\ref{par:inert-cut-exponents}, write \(f_N\) for the Wach map at the cut and \(f_N^{\rm tame}\) for its map on tame invariants. Then
\[
 \ell^R\operatorname{coker} f_N=0,\qquad \ell^R\operatorname{coker} f_N^{\rm tame}=0.
\]
In the extension it gives a lift of \(\ell^R e_0\), and gives
the corresponding constant-map bound after reduction modulo \(X\).
\end{lemma}

\begin{proof}
For an integral crystalline surjection, let \(s_i\) be the
least exponent of the ordinary cyclotomic parameter \(\pi\)
that kills the Wach cokernel at letter \(i\), before tame
invariants. The global height box in this case is symmetric
and has width
\[
 H_{\max}=\max_i H_i=2h.
\]
A common Tate
twist moves both representations into the positive range
of Berger's convention. This multiplies both Wach lattices
by the same power of \(\pi\) in the twisted coordinates and
preserves the cokernel. Berger's integral comparison
\cite[Theorem III.3.1]{berger-wach} places the base change of
each positive Wach lattice between the corresponding base
change of \(T_i\) and its sublattice multiplied by
\(\pi^{H_{\max}}\). Apply the comparison over
\[
 \widetilde{\mathbf A}^{+}=W(\widetilde{\mathbf E}^{+}),
\]
which contains the integral period ring of that comparison.

Since the integral Galois lattice map is onto,
\(\pi^{H_{\max}}N_2\) lies in the image after this base change.
This base change is faithfully flat over
\(\mathcal O_F[[\pi]]\). Modulo \(\ell\) the tilt valuation
ring is torsion-free over the DVR \(k_F[[\overline\pi]]\),
and \(\overline\pi\) is a nonunit in it. The ring of Witt
vectors is \(\ell\)-torsion-free and \(\ell\)-adically complete,
so the adic flatness criterion gives flatness
\cite[Lemma 15.28.5]{stacks-flat}. The residue of the maximal
ideal is proper after extension, giving faithfulness.
Extending by \(O\) and taking the embedding factors preserve
faithful flatness. In the ordinary cokernel
\[
 C_i=N_{2,i}/f(N_{1,i}),
\]
the inclusion
after base change reads
\[
 (\pi^{H_{\max}}C_i)\otimes_{\mathcal O_F[[\pi]]}
      \widetilde{\mathbf A}^+=0.
\]
Faithful flatness implies \(\pi^{H_{\max}}C_i=0\), hence \(s_i\le H_{\max}\). The same argument applies to every crystalline
extension of the trivial lattice by \(T_2\), since its
jumps are the union of the target jumps and zero.

\proofstep{The scalar loss from the cyclotomic action.}
Use a topological generator \(\gamma_*\) of the full
cyclotomic group whose cyclotomic value \(u_*\) generates
\(\Z_\ell^\times\). This action is trivial on Wach lattices
modulo \(\pi\). Suppose \(\ell^k\pi^r N_{2,i}\subset f(N_{1,i})\)
for \(r>0\). For any \(v\in N_{2,i}\), choose \(x\) lifting
\(\ell^k\pi^r v\). The vector \((\gamma_*-1)x/\pi\) is
Wach-integral. Its image, modulo \(\ell^k\pi^rN_{2,i}\), is
\[
                   \ell^k\pi^{r-1}(u_*^r-1)v .
\]
This follows by reducing \(\gamma_*(\pi)/\pi\) and the action
on \(v\) modulo \(\pi\).

The error is already in the image.
One therefore lowers \(r\) by one at the cost
\(v_\ell(u_*^r-1)\). The cost is zero unless
\[
 \ell-1\mid r,
\]
and it is \(1+v_\ell(r/(\ell-1))\)
otherwise. Thus the initial parameter exponent \(s\) gives
a scalar exponent at most \(\ell s/(\ell-1)^2\) by
Legendre's formula.

\proofstep{The recurrence in the letter exponents.}
To pass through Frobenius, put \(d_0=\ell-1\) and let
\(s=s_i\). From \(s\) down to \(r<s\), the exact cost is
\(n-m+v_\ell(n!/m!)\), with \(n=\lfloor s/d_0\rfloor\)
and \(m=\lfloor r/d_0\rfloor\). For \(n=0\) this is zero.
Otherwise it is at most
\[
         \frac{\ell(s-r)}{d_0^2}+\frac{\ell}{d_0}
                         +\lfloor\log_\ell n\rfloor .
\]
Indeed each term in the factorial valuation differs from
\((n-m)/\ell^v\) by less than one, and there are at most
\(\lfloor\log_\ell n\rfloor\) nonzero terms.

Take
\(a=\lceil s/d_0\rceil+c_H+2\). Since
\(q_{\rm cyc}\in(\ell,\pi^{d_0})\), the power
\(\phi(\pi)^a=\pi^a q_{\rm cyc}^a\) is in the ideal
generated by \(\ell^{a-j'}\pi^{a+d_0j'}\) for
\(0\le j'\le a\).
Those terms with \(r=a+d_0j'\ge s\) kill the cokernel.
For \(r<s\),
\[
 a-j'-\frac{\ell(s-r)}{d_0^2}
 =a-\frac{\ell(s-a)}{d_0^2}+\frac{j'}{d_0}
 \ge\frac{a}{d_0^2}+\frac{\ell(c_H+2)}{d_0}.
\]
This exceeds \(\ell/d_0+\lfloor\log_\ell n\rfloor\) since
\(n\le H_{\max}\). Hence each term kills the cokernel by
the cost bound, and so does \(\phi(\pi)^a\).

Write \(H_i=M_i+h_i\) for the letter's width, with the two ordinary one-step height inclusions
\[
 q_{\rm cyc}^{M_i}\Phi N_{2,i-1}\subset N_{2,i},\qquad q_{\rm cyc}^{h_i}N_{1,i}\subset\operatorname{span}(\Phi N_{1,i-1}).
\]

Lift \(\phi(\pi)^a q_{\rm cyc}^{M_i}\Phi v\) for
\(v\in N_{2,i-1}\) and multiply its lift by
\(q_{\rm cyc}^{h_i}\pi^{H_i}\).
It now lies in the source Frobenius span and maps to
\(\Phi(\pi^{a+H_i}v)\).
Frobenius linearization is injective on both lattices,
as it is invertible on the ambient \((\Phi,\gamma)\)-modules
after localization. The scalar ring is free over its
Frobenius image with a basis starting at \(1\), by
Weierstrass division by \(\phi(\pi)-Y\).
Compare the coefficient of \(1\) in the Frobenius-span
expression for the lift.

This constant coefficient gives a lift of
\(\pi^{a+H_i}v\) in the preceding letter. Consequently
\[
                  s_{i-1}\le H_i+c_H+3+\frac{s_i}{\ell-1}.
\]
Iterating the recurrence leaves a remainder bounded by a
geometric fraction of \(H_{\max}\), which tends to zero.
This gives \(s_i\le S(i)\), in addition to
\(s_i\le H_{\max}\). The scalar annihilator above therefore
has exponent at most \(R\) at each cut.

\proofstep{Passage to tame invariants.}
Averaging a lift of an invariant vector by the tame group costs
a unit. Thus the ordinary-module annihilator yields the desired
annihilator on the tame-invariant Wach lattices. This gives
the lift of \(\ell^R e_0\) in the extension Wach lattice and
the ordinary constant-map bound at the preceding cut after
reducing the surjection modulo \(X\). Dividing by \(\Theta\)
commutes with the map and is possible in the ambient module
because the factors of \(\Theta\) are units there. It gives
the required lift on the \(\Theta\)-divisible lattices at
the terminal cut.
\end{proof}


\begin{proposition}[Finite-class colength in the long-orbit range]\label{prop:inert-long-colength}
Use the marked crystalline surjection and signed jump windows
of Lemma~\ref{hyp:long-orbit-data}, with the cyclic spacing
and arc exponents of Paragraph~\ref{par:inert-cut-exponents}.
Then \eqref{eq:long-orbit-exponent} gives a normalized colength on
integral finite classes at most \(3C_2U_B\), with \(U_B\) as in
Section~\ref{sec:cut-length-rows}.
There exists one rational place entry \(c_w\) for all good characters at \(w\), satisfying
\[
 3C_2U_B\le c_w\le(3C_2+1)U_B.
\]
\end{proposition}

\begin{proof}
The lifting argument behind \eqref{eq:long-orbit-exponent} is
in the two operators \(\Phi-1\) and \(\gamma-1\) of the Herr
complex. Here \(\Phi\) advances one letter and \(\gamma\) acts
within a letter. Work in the tame-invariant Wach coordinates,
with parameter \(X\) and Wach lattices \(N_i\) at the letters.
The first filtration estimate places \(\Phi(X^M N)\) in
\(X^M(\ell,X)N\) on the product \(N\) of the letters. On \(X^MN\) contraction therefore gives the convergent inverse
\[
 (\Phi-1)^{-1}=-\sum_{a\ge0}\Phi^a.
\]
The coefficients are complete and
\(\Phi\) preserves the ideal \((\ell,X)\).

\proofstep{Represent the scaled class by a restricted Herr pair and lift its arcs.}
A finite class in the target has a Herr representative attached to
a crystalline extension of the trivial lattice by \(T_2\). If \(E_0\) lifts the trivial generator in period coordinates, the target pair representing the connecting class is
\[
 ((\Phi-1)E_0,(\gamma-1)E_0).
\]
The
cut colength bound provides, at each cut \(c\), a vector \(x_c\)
in the extension Wach lattice over \(\ell^R\) times the trivial
generator. The difference \((\gamma-1)x_c\) belongs to \(XN_c\)
of the target.

In fact it belongs to \(X\) times the extension
Wach lattice by triviality modulo \(X\), and to the target
period module because its image in the trivial quotient vanishes.
Since \(X\) is regular on the trivial quotient and the
extension Wach lattice intersects the target period module in the
target Wach lattice, this gives the assertion. After advancing
\(j\) letters and multiplying by \(\ell^\sigma\), the correction
of Paragraph~\ref{par:wach-cut-correction} gives
\[
       Q_c=\ell^\sigma\Phi^j x_c-\tau_c,
       \qquad (\gamma-1)Q_c\in X^M N_{c+j}.
\]
Use the target Wach lattice as its image in the extension.

For consecutive cuts \(c',c\), write \(a\ge j\) for their
forward distance. At the letters from \(c'+j\) up to the
letter before \(c+j\), set \(E_{c'+j+t}=\Phi^tQ_{c'}\).
This gives a global vector \(E\) with trivial quotient
\(\ell^{R+\sigma}\) times the generator. The first component
\(\ell^\sigma(\Phi-1)E\) vanishes at all the letters within
each propagated arc. At a shifted cut it is
\[
 \begin{split}
   y_c&=\ell^\sigma(\Phi^a Q_{c'}-Q_c)
       =\ell^\sigma(\Phi^j w_c+\tau_c),\\
   w_c&=\Phi^{a-j}Q_{c'}-\ell^\sigma x_c.
 \end{split}
\]
The two terms defining \(w_c\) have the same trivial
quotient, so \(w_c\) is in the target. Since \(\gamma\)
commutes with \(\Phi\), the deep-lattice defects of the
\(Q_c\)'s propagate along each arc and
\((\gamma-1)y_c\in X^M N\).

Set \(z_c=(\Phi-1)^{-1}(\gamma-1)y_c\) there.
Uniqueness of the inverse on \(X^MN\) gives
\[
 \sum_c(y_c,z_c)
    =\bigl(\ell^\sigma(\Phi-1)E,\ell^\sigma(\gamma-1)E\bigr).
\]
The difference from \(\ell^{R+2\sigma}\) times the original
pair is the Herr coboundary of
\(\ell^\sigma E-\ell^{R+2\sigma}E_0\). This is a target
vector because its two terms have equal trivial quotient.

Let \(c'\) be the preceding cut and \(a\) the length of its arc
to \(c\). The propagated \(Q_{c'}\) at \(c\) belongs to the span
of \(\Phi^a\) of the extension Wach lattice at \(c'\), since its
correction at \(c'+j\) is in the span of \(\Phi^jN_{c'}\).
The signed bound into the Frobenius span in
Paragraph~\ref{par:wach-level-models} gives the same membership
after multiplying \(x_c\) by
\[
 \Lambda=\prod_{s<a}p_s^{H_s^{\mathrm{arc}}}.
\]

The filtration of this extension has the same bounds as the target
with the jump at zero included.

For \(\Lambda w_c\), membership in the extension Frobenius
span and in the target period module already holds. To
intersect these two modules, put \(\mathcal R=O[[X]]\)
and use a basis \((b_t)_t\) of \(\mathcal R\) over
\(\phi^a(\mathcal R)\). Such a basis comes from Weierstrass
division by \(\phi^a(X)-Y\) over \(O[[Y]]\), since
\(\phi^a(X)\) first has a unit coefficient in degree \(\ell^a\).
Group the extension expression as \(\sum_t b_t\Phi^a u_t\).
If \(d_t\) is the trivial quotient of \(u_t\), then
\[
       \sum_t b_t\Phi^a u_t\longmapsto
                    \left(\sum_t b_t\phi^a(d_t)\right)e_0=0 .
\]
Independence of the basis and injectivity of \(\phi^a\)
give \(d_t=0\) for every \(t\). Therefore every \(u_t\)
lies in the target Wach lattice at \(c'\). This proves
\(w_c\in W_2\). The two defects to be subtracted belong respectively to \(X^MN_c\) and \(XN_c\). Therefore
\[
 (\gamma-1)w_c\in X^M N_c+XN_c\subseteq XN_c.
\]


For a change of cut vector by \(u\in XN_c\), the correction
of Paragraph~\ref{par:wach-cut-correction} changes by
\(\ell^\sigma\Phi^j u\) modulo the transported error.
It cancels the change of the uncorrected advanced vector.
Hence the corrected first component changes by a vector
\[
 v\in X^M N.
\]
The cocycle equation determines the change
of the second component as \((\Phi-1)^{-1}(\gamma-1)v\).
Together these give the coboundary of \((\Phi-1)^{-1}v\).
Thus comparison of cut vectors modulo \(XN_c\) suffices
for the finite classes of the corrected arc pairs.

Choose a product \(D_{\rm prev}\) of cyclotomic factors
clearing \(\Lambda w_c\) into \(N_c\), by iterating the
one-step height inclusions on the preceding arc. Choose
\(D_{\rm adv}\) clearing \(\Phi^j N_c\) into \(N_{c+j}\)
and \(D_\tau\) clearing \(P_c\) into \(N_{c+j}\).
Then the product
\[
            D=\phi^j(D_{\rm prev}\Lambda)\,D_{\rm adv}D_\tau
\]
clears \(y_c=\ell^\sigma(\Phi^j w_c+\tau_c)\) into \(N_{c+j}\)
by semilinearity and \(\tau_c\in P_c\).
The second component \(z_c\) and the defect
\((\gamma-1)y_c=(\Phi-1)z_c\) are in \(X^MN\).
Thus the two components satisfy the actual finite
cyclotomic-denominator and deep-lattice tests.

Pass to the ordinary cyclotomic parameter \(\pi\). Put
\(V=T\otimes_O K\), with \(K=\operatorname{Frac}(O)\), and set
\[
 q_{\rm cyc}=\phi(\pi)/\pi,\qquad t=\log(1+\pi),\qquad
 \lambda=t/\pi .
\]
The factors \(p_s\) occurring in \(D\) are associates of
\(\phi^s(q_{\rm cyc})\). Also \((X^M)=(\pi^{(\ell-1)M})\).
Thus, for some finitely supported family of integers \(b_s\ge0\),
the Herr pair \((y,z)\) satisfies
\begin{equation}\label{eq:long-arc-herr-lattices}
 \Bigl(\prod_s\phi^s(q_{\rm cyc})^{b_s}\Bigr)y\in N(T),\qquad
 z\in\pi N(T),\qquad (\gamma-1)y=(\Phi-1)z .
\end{equation}
The integral Herr isomorphism gives a class in \(H^1(F,T)\).
Its rational image is represented by the same pair in the Robba-ring
Herr complex \cite[Theorems 2.3 and 2.8]{liu-herr}.
In the corresponding extension \(0\to V\to V'\to K\to0\), choose
a tame-invariant lift \(e_0\) of the trivial basis. The extension
actions are \(\Phi e_0=e_0+y\) and \(\gamma e_0=e_0+z\).

Let \(\mathcal B^+\) denote the analytic functions on \(|\pi|<1\).
Use crystalline coordinates in the kernel. For a positive crystalline
representation the Wach comparison and its elementary divisors put the
change of basis and its inverse over \(\mathcal B^+[1/\lambda]\)
\cite[Propositions II.2.1 and III.2.1]{berger-wach}.
A Tate twist gives the same conclusion for \(V\): the ratio of
the two twist generators is
\[
 \pi^{-k}/t^{-k}=\lambda^k.
\]

Since \(\lambda(0)=1\), the change of basis is regular at zero
in both directions. Writing brackets for crystalline coordinates gives
\[
 \lambda=\prod_{s\ge0}\frac{\phi^s(q_{\rm cyc})}{\ell},
 \qquad \lambda^B[y]\in(\mathcal B^+)^{\dim V}\quad(B\gg 0),
 \qquad [z](0)=0 .
\]
By \eqref{eq:long-arc-herr-lattices}, a sufficiently high power of the product absorbs all finitely
many cyclotomic denominator factors of \(y\) and the poles of the
change of basis. The factor \(\pi\) in \(z\) gives the last equality.

In the crystalline basis the equation \((\gamma-1)v=z\)
has a solution analytic near zero. Use \(t\) as local
coordinate and put \(\omega=\chi(\gamma)\). If
\(z(t)=\sum_{k\ge1} z_k t^k\), the solution with zero
constant is
\[
          v(t)=\sum_{k\ge1}\frac{z_k}{\omega^k-1}t^k .
\]
The valuations \(1+v_\ell(k)\) of the denominators grow
logarithmically in \(k\). Thus any strictly smaller positive disk of convergence of
\(z\) is a disk of convergence of \(v\). For \(u=y-\Phi v+v\)
the Herr equation and \(\gamma\Phi=\Phi\gamma\) give
\[
 (\gamma-1)u=(\Phi-1)z-(\Phi-1)(\gamma-1)v=0 .
\]

In the crystalline basis the coefficient \(u_k\) in degree \(k\)
is therefore killed by \(\chi(\gamma)^k-1\). For \(k>0\) this
factor is nonzero, so the coordinates of \(u\) are constant. For every \(m\ge1\)
the identity of germs is
\[
 v=\Phi^m v-\sum_{i=0}^{m-1}\Phi^i y+\sum_{i=0}^{m-1}\Phi^i u.
\]
On every disk of fixed radius less than one,
\(\phi^m(\pi)=(1+\pi)^{\ell^m}-1\) tends uniformly to zero.
Hence \(\Phi^m v\) is analytic there for sufficiently large
\(m\). Also \(\lambda^B\Phi^i y\) is analytic on the
whole open disk for every \(i\), since
\[
 \lambda/\phi^i\lambda=\prod_{k<i}\phi^k(q_{\rm cyc}/\ell).
\]

The displayed identity therefore extends \(\lambda^B v\)
to the whole disk. The continuations agree on overlaps by
their Taylor series at zero. It follows that \(v\) belongs
to the Robba module with \(t\) inverted. The identity
\[
 (\gamma-1)v=z
\]
holds there by its power series identity
at zero after clearing powers of \(\lambda\), so \(e_0-v\)
is \(\gamma\)-invariant. Its average over the tame
group is invariant under the full cyclotomic group and still
maps to the trivial basis. Berger's comparison
\cite[Theorem 0.2]{berger-differential} identifies these
invariants with \(D_{\rm cris}(V')\). The rank increases
by one from \(D_{\rm cris}(V)\), so \(V'\) is crystalline.
Thus the integral Herr class is finite. This criterion applies
to both lattices and is compatible with the same row map.

The quantitative estimate at an arc supplies a source vector
of this kind mapping to \(m_c w_c\) modulo \(XN_c\), where
\[
 m_c=u_c\ell^{2\sigma+2a_0+R+k_c},\qquad u_c\in O^\times,
                 \qquad 0\le k_c\le2\kappa.
\]
It applies the polynomials in \(\gamma\) to the target and then
to its lift to impose the required local exponent windows.
Modulo \(XN_c\) their actions reduce to their values at \(1\),
since \(w_c\) is \(\gamma\)-invariant modulo this submodule.

Let \(\omega\) be the cyclotomic value of \(\gamma\), with
\(v_\ell(\omega-1)=1\), and let \(E=a-1\le A\).
The polynomials have the form
\[
 P_{E,J}(Z)=\left(\sum_{v=0}^{\ell^E-1}Z^v\right)
            \prod_{s\in J\setminus\{0\}}
                 (Z^{\ell^E}-\omega^{s\ell^E}),
\]
At \(1\) the valuation is
\(E+\sum_{s\in J\setminus\{0\}}(1+E+v_\ell(s))\)
by the valuation of powers of a principal unit. The window
bound gives at most \(\kappa\) for each polynomial. Between the two
polynomials the comparison of Paragraph~\ref{par:wach-mixed-constant}
costs \(2\sigma+2a_0+R\) on the level inputs.
Paragraph~\ref{par:wach-cut-colength} gives the cut colength
on the two-sided Wach height bounds. The simultaneous level
realization and the two polynomial actions are those of
Paragraphs~\ref{par:wach-level-models}--\ref{par:wach-polynomial-actions}.

Multiply the finite source class for arc \(c\) by
\(u_c^{-1}\ell^{2\kappa-k_c}\), and sum over the cuts. Its image
is \(\ell^{2\sigma+2a_0+R+2\kappa}\) times the sum of the target
arc classes. Together with the initial decomposition this gives
\[
 e=2R+4\sigma+2a_0+2\kappa.
\]
The sum over the cuts costs no
additional power of \(\ell\).

\proofstep{Include the torsion generators and then take the normalized colength.}
The free part of the arc calculation bounds the failure of lifting along the orbit. Torsion in the generators must be included before the resulting length can enter the inert-prime table.


Let \(V_2=T_2\otimes_O\operatorname{Frac}(O)\) and
\(r_0=\dim V_2^{G_F}\). By the Bloch--Kato dimension formula the
free rank of the integral finite group is
\[
 \dim(D_{\mathrm{dR}}(V_2)/\mathrm{Fil}^0)+r_0\le dr+r_0.
\]

For a uniformizer \(\varpi\) of \(O\), the long exact sequence
for multiplication by \(\varpi\) gives
\[
 0\longrightarrow T_2^{G_F}/\varpi T_2^{G_F}
   \longrightarrow (T_2/\varpi T_2)^{G_F}
   \longrightarrow H^1(F,T_2)[\varpi]\longrightarrow0.
\]
The first term has dimension \(r_0\) and the middle term has
dimension at most \(r\). Thus at most \(r-r_0\) torsion generators
are needed. The rank of the free quotient and the number of torsion generators together cost at most
\[
 (dr+r_0)+(r-r_0)=(d+1)r.
\]
Since its quotient by the image of
the finite group of \(T_1\) is killed by \(\ell^e\), its length
divided by \(\operatorname{length}_O(O/\ell O)\) is at most
\((d+1)re\). This is the normalization used for the actual row
cokernel. It handles all coefficient ramification degrees.

Finally, for \(h\ge1\) use
\[
 \log_\ell h\le h/D_0+j+1,
\]
by splitting the logarithm at
\(D_0\) and using \(\log_\ell t\le t\) for \(t>0\).
When the word has a position with larger jumps, the bracket in the bound for
\(e\) is at most \(2H/D_0+2(G_P+1)(j+1)\). Otherwise \(h=g\le G_P\)
and it is at most \(3(G_P+1)(j+1)\).

In both cases the resulting
normalized colength is at most \(3C_2 U_B\), with \(U_B\) as in
Section~\ref{sec:cut-length-rows}. Choose a rational table entry \(u_B\) satisfying
\[
 3C_2U_B<u_B<(3C_2+1)U_B.
\]
The lower bound makes \(u_B\) valid for the cokernel at every good character. Such an entry exists since \(U_B>0\).
This gives case B of the ordinary congruence comparison on the
tied local data, with one constant and cutoff for every place,
coefficient field and marked character in the window.

Use the single word chosen for the large-orbit estimate.

Sum the arc bounds and convert the result to the rational entry in its marked row.
\label{cor:inert-long-orbit-1}

The filtration statement for the Wach module
of a crystalline representation over an unramified field is
Berger's filtered comparison \cite[Corollary III.4.5]{berger-wach}.
The exponent in \eqref{eq:long-orbit-exponent} belongs to the
submission's quantitative lifting argument.
\end{proof}


\subsubsection{Alternating sequences with larger jumps and the finite-class length}
\label{sec:inert-thick-word}
\hypertarget{source-inert-thick-word}{}

For an alternating sequence whose larger jumps dominate, the finite-class estimate needs information about the image of each successive lattice map. We follow a lifted class through the sequence and control the restrictions on its coefficients at the even positions.

\begin{established}[Marked alternating word]\label{hyp:alternating-thick-word}
Let \(F/\Q_\ell\) be unramified of even degree \(n\), with
\(\ell\ge11\). Let \(f:T_1\to T_2\) be the marked crystalline surjection of
free stable lattices over the coefficient ring \(O\).
Take its field \(K_O\) to contain the embeddings of \(F\) and
the primitive \(\ell\)-th roots. The descent
of the finite-class bound to the original coefficients uses the same
faithfully flat base change as in Section~\ref{sec:inert-long-orbit}.
Write \(r=\rk_O T_2\). Index the embeddings cyclically and absorb the
parity shift in the choice of origin. If \(J_i\) denotes the source jumps at letter \(i\), their boxes are
\[
 J_i\subseteq[-\eta,\eta]\cup\begin{cases}[-b_i-S,-b_i],&i\text{ even},\\ {}[b_i,b_i+S],&i\text{ odd}.\end{cases}
\]
At most \(m\) target jumps, counted
with multiplicity, lie in the central interval at each embedding.
The target jumps are among the source jumps by strictness of the
crystalline filtration.

Put
\[
\begin{gathered}
 h=\max_i b_i+S,\quad H_{\rm loc}=h+\eta,\quad
 L_0=1+\log_\ell(2H_{\rm loc}),\\
 A_0=\left\lceil\frac{(2\ell-1)H_{\rm loc}}{(\ell-1)^2}
                +3+L_0\right\rceil .
\end{gathered}
\]
In the retained alternating case one has
\(b_i\ge h/2\) and
\(b_i\ge 2(A_0^{\rm real})+\eta+1\), where \(A_0^{\rm real}\) is
the expression inside the ceiling. Define
\[
 c_D=1+\frac{H_{\rm loc}+1}{\ell-1},\qquad
 C_D=rS+2mc_D .
\]
For an \(O\)-module of finite length, we use its length divided by
\(k_O=\operatorname{length}_O(O/\ell O)\).
\end{established}
\begin{proof}
At an inert prime the two embeddings of the distinguished quadratic field are exchanged by Frobenius. On an orbit carrying larger jumps at every position, Lemma~\ref{lem:marked-signed-jumps} therefore gives the displayed signs after a cyclic reindexing. The central multiplicity is the number of central columns in the marked rank-\(r\) weights. The target inherits the source support by strict exactness. On the retained word
\(b_0\le b_i\le b_0+\eta\) with \(b_0=b+O(t)\), whereas
\(S+\eta+\log H=o(b_0)\) uniformly in the host. Thus
\(h\le b_0+o(b_0)\), and the coefficient
\(2(2\ell-1)/(\ell-1)^2\le 42/100\) for \(\ell\ge11\)
gives the second inequality beyond a common threshold.
In particular \(b_i>\eta\).
\end{proof}

\paragraph{Passing from a chain to a length.}
Set \(M_2=H^1_f(F,T_2)\) and \(I=f_*H^1_f(F,T_1)\).
Choose a free \(O\)-module \(A\) of rank at most \((n+1)r\)
with a surjection \(A\to M_2\), using the generator argument in
Section~\ref{sec:inert-long-orbit}. Representatives for a basis
give a linear map to the finite cocycles themselves.
\begin{lemma}[Length from the successive restrictions]
Let \(\ell^{e'}M_2\subset I\). Suppose the final submodule of a
chain in \(A\) maps into \(I\) after multiplication by \(\ell^e\).
Let \(N_g,N_d\) be the numbers of the generator-bounded and length-bounded successive quotients, respectively. If each of the former requires at most \(m\) generators and each of the latter has normalized length at most \(C_D\), the numerical hypotheses
\[
 N_g\le4n,\qquad N_d\le3n/2
\]
give the bound \eqref{eq:thick-chain-mass} on \(\operatorname{len}(M_2/I)\).
\end{lemma}
\begin{proof}
Suppose \(A=A^{(0)}\supset\cdots\supset A^{(k)}\) and
\(\ell^e A^{(k)}\) maps into \(I\).
Suppose also that \(\ell^{e'}M_2\subset I\).
If a successive quotient of the chain is generated by at most
\(m\) elements, its image in the filtration of \(M_2/I\) costs at
most \(me'\) in normalized length. If a successive quotient has
normalized length at most \(C_D\), the corresponding image costs
at most \(C_D\). Indeed these successive images are quotients of
the successive quotients of the chain.

At the last term the image
is killed by \(\ell^e\) and has at most \(\rk_O A\) generators,
since \(A^{(k)}\) is free of rank at most \(\rk_O A\).
Thus \(4n\) steps of the first kind and \(3n/2\) of the second
give normalized length at most
\begin{equation}\label{eq:thick-chain-mass}
           (n+1)r e+4nm e'+\tfrac32 n C_D .
\end{equation}
For later use, a different sufficient condition is
\(\ell^e M_2\subset I+\langle Z\rangle_O\), where \(Z\) is a finite
set of classes in \(M_2\). The normalized length is then at most
\((n+1)r e+|Z|e'\). Take the quotient by the image of
\(\langle Z\rangle_O\) first, and use \(\ell^{e'}M_2\subset I\)
on that image. These arguments apply to the full integral finite
group, including its torsion.

\end{proof}
\paragraph{Wach coordinates and the initial finite classes.}
\label{par:thick-initial-classes}
Use the tame-invariant parameter \(X\) and \(q=X+\ell\) of
Paragraph~\ref{par:wach-cut-correction}, and put \(p_j=\phi^j(q)\).
Write \(N_i(T)\) for the tame-invariant lattice, \(N(T)\) for their
product, and \(P_i^{(j)}(T)=\langle\Phi^jN_i(T)\rangle\) for its
integral Frobenius span at \(i+j\). Let \(a=A_0\),
\(B_i=b_i+S\) and \(H=H_{\rm loc}\). The two signed height bounds
at a destination \(i\) are
\[
 q^{M_i}P_{i-1}^{(1)}(T)\subset N_i(T),\qquad
 q^{H_i}N_i(T)\subset P_{i-1}^{(1)}(T),\qquad
 (M_i,H_i)=
 \begin{cases}(B_i,\eta)&i\text{ even},\\
 (\eta,B_i)&i\text{ odd}.\end{cases}
\]
Use the cyclotomic lattice comparison over the unramified field
and the filtered comparisons at the primitive roots
\cite[Proposition III.2.1 and Corollary III.4.5]{berger-wach}.
\begin{lemma}[Annihilators at the initial letters]
For the given surjection, and for the crystalline extension of
the trivial lattice by the target with its quotient map, let \(C_i\) be the cokernel on the \(i\)-th tame-invariant letter. Then
\[
 q^a C_i=\ell^R C_i=0,\qquad R=\left\lfloor\frac{\ell H}{(\ell-1)^2}+3+\log_\ell(H+1)\right\rfloor.
\]
\end{lemma}
\begin{proof}
The following image bound uses the maximum \emph{letter width}
\(H=\max_i(M_i+H_i)\). Do the calculation before tame invariants
in the ordinary parameter \(\pi\).
The common character with labelled filtration jump \(+M_i\) shifts both lattices by
\[
 [-M_i,H_i]+M_i=[0,M_i+H_i]\subseteq[0,H].
\]

The jumps of the decreasing filtration are now in \([0,H]\),
as required in Berger's positive comparison. Such a character is a product
of powers of the labeled Lubin--Tate characters of the
unramified field, with values in the coefficient field.
\cite[Proposition 2.3]{berger-lubin-tate}
Tensoring by this rank-one lattice preserves the annihilator
of the Wach cokernel. Rational Wach comparison is tensor
compatible. Integrally, multiplication by a generator of the
rank-one Wach lattice identifies the integral \'etale modules,
since the generator is a basis after \'etale scalar extension.
Taking their intersections with the rational Wach modules gives
the tensor identification on integral Wach lattices.
\cite[Proposition III.4.2]{berger-wach}

Apply Berger's positive comparison and the faithfully flat
descent in Paragraph~\ref{par:wach-cut-colength} after the common
twist. For the untwisted Wach cokernel \(C\), the common rank-one twist preserves the annihilator and hence gives
\[
 \pi^H C=0.
\]

The argument of that paragraph which lowers the \(\pi\)-degree
by the divided cyclotomic action gives a scalar exponent at most
\(\lfloor\ell H/(\ell-1)^2\rfloor\). The two input lattices also
occur as the extension of the trivial lattice by the target
and its trivial quotient. The jumps of the extension are those
of the target together with zero, so the same boxes apply.

After tame averaging the annihilators on each letter include
\(X^{\lfloor H/(\ell-1)\rfloor+1}\) and \(\ell^R\), with the following integer exponent
\[
 R=\left\lfloor \frac{\ell H}{(\ell-1)^2}
                 +3+\log_\ell(H+1)\right\rfloor .
\]
Set \(d=\lfloor H/(\ell-1)\rfloor\), so the two annihilators are
\(X^{d+1}\) and \(\ell^R\). In the expansion of \((X+\ell)^a\),
each term with \(X\)-degree \(k\le d\) has \(\ell\)-exponent
\(a-k\ge R\); every other term is divisible by \(X^{d+1}\).
Thus \(q^a\) annihilates the same cokernel.

The bounds are uniform in the number of letters,
the rank and the ramification of the coefficient ring.
\label{cor:inert-thick-initial-1}

\end{proof}

\begin{lemma}[Integral Herr representatives for the basis extensions]
For every basis cocycle \(y\) the two annihilators give vectors \(\mu_i(y),\nu_i(y)\) satisfying \eqref{eq:thick-extension-data} on each letter. On the free span, the choice is extended by
\[
 (\mu_i,\nu_i)\!\left(\sum_j b_j y_j\right)=\sum_j b_j\big(\mu_i(y_j),\nu_i(y_j)\big).
\]
\end{lemma}
\begin{proof}
For each basis cocycle in the free family take its extension
\(0\to T_2\to Y\to O\to0\). Its rational extension is crystalline
by the definition of the finite group. Use exactness of the
integral \((\varphi,\Gamma)\)-module functor and the Wach
intersection with the kernel. Tame averaging, of order
\(\ell-1\), preserves integral lifts.

Choose a tame-invariant
lift \(E\) of the trivial generator in the \'etale module of
\(Y\), and take the Herr pair \(y=(\Phi-1)E\),
\(z=(\gamma-1)E\) in that of \(T_2\). Denote their letter
components by \(y_i,z_i\), and use the convention
\((\gamma-1)y=(\Phi-1)z\).
The two annihilators provide Wach lifts \(E_i,E_i'\) of
\(\ell^R\) and \(q^a\) times the trivial generator in letter \(i\).
Set
\[
 \mu_i=E_i-\ell^RE|_i,\qquad \nu_i=E_i'-q^a E|_i,\qquad
 A_i=\ell^R y_i+\Phi\mu_{i-1}-\mu_i .
\]
Both differences defining \(\mu_i,\nu_i\) belong to the target
\'etale module. The Wach intersection and the first signed height
bound give
\begin{equation}\label{eq:thick-extension-data}
\begin{gathered}
 \ell^R z_i+(\gamma-1)\mu_i\in XN_i(T_2),\qquad
 q^{M_i} A_i\in N_i(T_2),\\
 \ell^R q^{H_i} A_i\in P_{i-1}^{(1)}(T_2),\qquad
 \ell^R\nu_i-q^a\mu_i\in N_i(T_2).
\end{gathered}
\end{equation}

For the third inclusion, \(A_i=\Phi E_{i-1}-E_i\) in the
extension, so \(q^{H_i}A_i\) belongs to its Frobenius span.
Expand in a free basis of \(O[[X]]\) over \(\phi(O[[X]])\),
as in the proof of the arc estimate. Each quotient coefficient
of the resulting Wach vectors vanishes by independence.
Thus the Frobenius span meets the target in the target span.
This even gives the inclusion before multiplication by \(\ell^R\).
The data in \eqref{eq:thick-extension-data} now extend linearly in the free family,
even though the individual extensions \(Y\) were chosen only
on its basis.

\end{proof}
\paragraph{First restriction at an even letter.}
\begin{lemma}[The first restriction on the even letters]
At each even letter one can impose a single kernel condition \(z\in\ker\pi_i\) for which
\[
 \operatorname{im}\pi_i\text{ is generated by at most }m\text{ elements},\qquad z\in\ker\pi_i\ \Longrightarrow\ u_i(z)\in N_i(T_2),
\]
where \(u_i(z)\) is the corrected vector constructed in the proof. If \(m=0\) the assertion holds on
the full free family.
\end{lemma}
\begin{proof}
For \(x\in q^{-a}N_i(T_2)\), let
\(\partial x=X^{-1}(\gamma-1)x\) in the \'etale module.
It preserves \(q^{-a}N_i(T_2)\) and \(N_i(T_2)\).
Indeed the action on the invariant Wach lattice modulo \(X\)
is trivial, and \(\gamma(q)/q\) is a unit congruent to one
modulo \(X\). At a primitive root, put
\[
             c_j=\frac{1-\chi(\gamma)^{-j}}{\ell}.
\]
On the associated Laurent coefficient of degree \(-j\) the
divided operator acts by \(c_j\) since \(X=-\ell\) there.
For a Hodge jump \(u\), the quotient of the allowed pole lattice
by the unshifted lattice has degrees \(-u-k\), \(1\le k\le a\).

To see why one factor per distinct shift suffices, choose a basis
of the crystalline space adapted to the filtration. In each
direction the Laurent action is triangular, with distinct
diagonal entries indexed by the degrees in that direction.
The action of \(1/X(t)\) is scalar and the cyclotomic action
in the crystalline basis is substitution
\[
 t\mapsto\chi(\gamma)t.
\]

Thus the product vanishes separately on each direction even
when a shift repeats in another direction. This is a computation
over the coefficient field. The Hodge criterion at the cyclotomic
roots and saturation for \(\ell\) in \(q^{-a}N_i(T_2)/N_i(T_2)\)
give the same inclusion integrally.

Let \(J_i\) be the nonzero shifts that occur. They belong to
\([1-B_i,a-b_i]\cup[1-\eta,a+\eta]\).
By the valuation of powers of a principal unit,
\[
 v_\ell(c_j)=v_\ell(j)
\]
for \(j\ne0\). In the first
interval there are at most \((a+S)/\ell^k+1\) multiples of
\(\ell^k\), and no such nonzero multiples for
\(k>\log_\ell(2H)\). For the central interval split at zero
and apply the factorial bound on both sides.

Taking the maximum
over even letters, choose \(s_1\) and polynomials \(R_i(U)\)
proportional over \(O\) to \(\prod_{j\in J_i}(U-c_j)\) with
\[
 R_i(0)=\ell^{s_1},\qquad
 s_1\le(2a+2\eta+S)/(\ell-1)+L_0,\qquad
 \partial R_i(\partial)q^{-a}N_i(T_2)\subset N_i(T_2).
\]
Including \(\partial\) kills degree zero if that shift occurs.
The scaling of each polynomial is integral by the choice of the
maximum. Write \(R_i(U)=\ell^{s_1}+U R_i^*(U)\), and set
\[
 m_i=R_i(\partial)(q^{-a}\nu_i)+R_i^*(\partial)(z_i/X),
 \qquad u_i=\ell^R m_i-\ell^{s_1}\mu_i .
\]
Put \(v_i=\ell^R q^{-a}\nu_i-\mu_i\) and
\(d_i=\ell^R z_i/X+\partial\mu_i\).

By \eqref{eq:thick-extension-data},
\(q^a v_i,d_i\in N_i(T_2)\). The polynomial identity gives
\(u_i=R_i(\partial)v_i+R_i^*(\partial)d_i\). Thus
\(q^a u_i\in N_i(T_2)\) and \((\gamma-1)u_i\in XN_i(T_2)\).
Also
\[
 \ell^R\bigl(\ell^{s_1}z_i+(\gamma-1)m_i\bigr)
       =X\ell^{s_1}d_i+(\gamma-1)u_i .
\]
The bracket belongs to \(XN_i(T_2)\) by cancellation
of \(\ell^R\) in the integral \'etale module, where the Wach
lattice and its \(X\)-multiple are \(\ell\)-saturated. In coordinates this follows from the
description of the \'etale module as the completed localization
of the lattice, with injective reduction of the lattice modulo
the coefficient uniformizer.

At a primitive root, the only Laurent coefficient of \(u_i\)
that can obstruct membership in the Wach lattice is degree zero.
Writing \([t^j]\) for coefficient extraction at the root gives
\[
 [t^j](\gamma-1)u_i=(\chi(\gamma)^j-1)[t^j]u_i\qquad(j\ne0).
\]
The scalar on the right is nonzero. The degree-zero obstruction is therefore
the class in \(\mathrm{Fil}^{-a}D_i/\mathrm{Fil}^{0}D_i\).
All jumps in \([-a,0)\) are central since \(a<b_i\).
The quotient has dimension at most \(m\). Each obstruction
map is linear in the free family and its image is an
\(O\)-module generated by at most \(m\) elements.
Successively take these kernels for even \(i\). This leaves
a family on which every \(u_i\) belongs to \(N_i(T_2)\), using
one rank restriction per even letter. When \(m=0\) the entire
family survives.
\end{proof}

\paragraph{Cyclotomic correction at fixed precision.}
\label{par:thick-integral-correction}
Put
\[
 M=1+\lfloor h/(\ell-1)\rfloor,\qquad
 \sigma_1=\left\lceil 1+(M+\eta)/(\ell-1)\right\rceil .
\]
For a vector \(v\in XN_i(T)\) there is an integral correction
\(\tau(v)\in P_i^{(1)}(T)\) well-defined modulo
\(X^Mq^\eta P_i^{(1)}(T)\), with
\[
  \ell^{\sigma_1}\Phi v-(\gamma-1)\tau(v)
             \in X^Mq^\eta P_i^{(1)}(T).
\]
We use it at an even \(i\), where the odd destination gives
\(q^\eta P_i^{(1)}(T)\subset N_{i+1}(T)\).
There is also a two-step correction with precision \(Xq^{2H}\)
on \(P_i^{(2)}(T)\) and exponent
\(\sigma_*=\lceil1+2H/(\ell-1)^2\rceil\).

Here are the divisions that give these two claims. In either
case solve \((\gamma-1)d=v\) modulo a large power \(X^K\),
over the coefficient field. The degree-by-degree solution
of Paragraph~\ref{par:wach-cut-correction} puts the term in
degree \(y\) in \(\ell^{-e_y}X^yN_i(T)\), where
\[
 e_y=\lceil\ell y/(\ell-1)\rceil.
\]
Take \(d\) with zero
constant coefficient. Write \(J\) for the exponent of \(q\)
in the precision ideal and \(a_y=(\ell-1)y\). In one step
\(\phi(X)^y\) is a unit times \(X^yq^{a_y}\).

If \(a_y<J=\eta\), then \(e_y\le\sigma_1\), since
\(J\le (\ell-1)M\). If \(a_y\ge J\) and \(y<M\),
expand \(q^{a_y-J}=(X+\ell)^{a_y-J}\).
The terms which survive modulo \(X^Mq^J\) are divisible by
\(\ell^{\max(0,\ell y-J-M+1)}\). For \(e_y>\sigma_1\)
we have \(\ell y>J+M+\ell-1\); hence
\[
 \ell y-J-M+1\ge e_y-\sigma_1.
\]

For \(a_y\ge J\) and \(y\ge M\) the entire term vanishes.

For two steps use \(M=1\), \(J=2H\), and
\(\phi^2(X)^y\) associated to
\(X^yq^{a_y}p_1^{a_y}\).
If \(a_y\ge J\) the terms with \(y>0\) vanish. Otherwise
\(p_1\in(\ell,q^\ell)\), by reduction modulo \(\ell\).
After factoring out \(X^yq^{a_y}\), the terms of the binomial
expansion that can survive have \(q^\ell\) used at most
\(\lfloor(J-a_y-1)/\ell\rfloor\) times. If
\(e_y>\sigma_*\), then
\(a_y>J/\ell+(\ell-1)^2/\ell\), and
\[
 a_y-\frac{J-a_y-1}{\ell}
      \ge \frac{\ell a_y-J}{(\ell-1)^2}
      \ge e_y-\sigma_* .
\]
The last inequality uses the defining roundings. Thus the
required division is integral. Here one may take \(K>\max(M,J)\), since
\(\phi^j(X)^K\) then belongs to the precision ideal in both
cases.

The integral coefficient quotient by \(X^Mq^J\) is free
over \(O\) by monic Weierstrass division. For either free Frobenius span \(P\) and its precision submodule \(P_{\rm prec}\), saturation gives
\[
 (P_{\rm prec}[1/\ell])\cap P=P_{\rm prec}.
\]

The resulting representative of \(\ell^{\sigma_1}\Phi d\)
or \(\ell^{\sigma_*}\Phi^2d\) satisfies the claimed
cyclotomic congruence, by commutation and saturation.
Uniqueness of \(d\) modulo \(X^K\) over the coefficient field
follows from the nonzero diagonal \(\chi(\gamma)^{(\ell-1)y}-1\)
in positive degrees. Hence the corrections are linear modulo
the precision and natural in maps of crystalline lattices.
They may be made linear on a free family by choosing
representatives on its basis. In one step the correction to a
vector of \(X^MN_i(T)\) can be taken zero. The calculation of
\(d\) in that case starts at degree \(M\), and
\[
 (\ell-1)M\ge\eta.
\]
For \(v=(\gamma-1)w\) with
\(w\in XN_i(T)\) one can use \(d=w\).

\paragraph{Second restriction and the alternating sum.}
Write \(N_i=N_i(T_2)\) for this construction. For even \(c\)
set
\[
 \xi_c=\ell^{s_1}z_c+(\gamma-1)m_c\in XN_c.
\]

Choose \(t_c=\tau(\xi_c)\) linearly on the restricted family.
At the odd letter \(c+1\) put
\[
 L_c=\ell^{\sigma_1}\Phi m_c+
             \ell^{s_1+\sigma_1}y_{c+1}-t_c,\qquad
 \Xi_c=\ell^{s_1+\sigma_1}z_{c+1}+(\gamma-1)L_c .
\]
The Herr relation gives
\(\Xi_c=\ell^{\sigma_1}\Phi\xi_c-(\gamma-1)t_c\).

Consequently \(\Xi_c\in X^MN_{c+1}\). It lies more precisely
in \(X^Mq^\eta P_c^{(1)}(T_2)\).
To make the next difference integral, write
\[
 U_{c+1}=\ell^R L_c-\ell^{s_1+\sigma_1}\mu_{c+1}
        =\ell^{\sigma_1}\Phi u_c+\ell^{s_1+\sigma_1}A_{c+1}
           -\ell^R t_c .
\]
The signed height and \eqref{eq:thick-extension-data} show that \(q^\eta U_{c+1}\in N_{c+1}\).
Applying \(\gamma-1\) gives \(\ell^R\Xi_c\) minus
\(\ell^{s_1+\sigma_1}(\ell^R z_{c+1}+(\gamma-1)\mu_{c+1})\).
This is in \(N_{c+1}\). Its only possible Laurent obstruction
is again degree zero, this time in
\(\mathrm{Fil}^{-\eta}D_{c+1}/\mathrm{Fil}^0D_{c+1}\).
Only central jumps contribute, so taking one kernel of rank
at most \(m\) per odd destination makes each \(U_{c+1}\) integral.
There are at most \(n\) restrictions including the first ones.

On the resulting family define at each even \(c\)
\[
 w_c=\Phi L_{c-2}+\ell^{s_1+\sigma_1}y_c-\ell^{\sigma_1}m_c,
 \qquad Y_c=\Phi w_c+t_c .
\]
All indices are cyclic, including when the word has length two.
There are three useful tests on \(w_c\). First,
\[
 \ell^R w_c=\Phi U_{c-1}-\ell^{\sigma_1}u_c
                        +\ell^{s_1+\sigma_1}A_c .
\]
The height bounds place the right side, after multiplication
by \(q^{B_c}\), in \(N_c\). Multiplication instead by \(q^\eta\)
places it in \(P_{c-1}^{(1)}(T_2)\) after a possible
additional factor \(\ell^R\), by \eqref{eq:thick-extension-data}.
Both lattices are \(\ell\)-saturated inside the integral
\'etale module.

For the Frobenius span, reduce in a Wach basis
and extend scalars by the finite free Frobenius coefficient map.
The original basis stays independent over the residue Laurent
field because the linearized Frobenius on the \'etale module
is invertible. Cancel these powers of \(\ell\).
Second, \((\gamma-1)w_c=\Phi\Xi_{c-2}-\ell^{\sigma_1}\xi_c\)
belongs to \(XN_c\), since
\(\Phi(X^MN)\subset q X^MN\). Third,
\[
                (\gamma-1)Y_c=\Phi^2\Xi_{c-2}-\Xi_c
                                 \in X^MN_{c+1}.
\]
These calculations give exactly the pole, first pullback
and cyclotomic tests on \(w_c\) used by the even-letter estimates.

Define \(Z_c=(\Phi-1)^{-1}(\gamma-1)Y_c\) in \(X^MN\),
where the inverse exists by contraction on the product of all
letters. Then the sum of \((Y_c,Z_c)\) over even \(c\)
represents \(\ell^{s_1+\sigma_1}(y,z)\). Indeed let
\[
  B=\sum_{c\text{ even}}\bigl(L_c+\Phi L_c
                                  +\ell^{s_1+\sigma_1}y_{c+2}\bigr).
\]
Summing first components gives
\(\sum_cY_c=\ell^{s_1+\sigma_1}y+(\Phi-1)B\)
by invariance of the even labels under the shift by two.
For the second components
\(\ell^{s_1+\sigma_1}z+(\gamma-1)B
 =\sum_c(\Xi_c+\Phi\Xi_c)\in X^MN\).
The Herr relation and uniqueness of the inverse prove the same
coboundary identity for the whole pair.
\label{cor:inert-thick-herr-1}

\paragraph{Integral comparison with the lifted vectors.}
For an admissible \(w\) at an even \(c\), use the first
correction to form
\[
             V(w)=\ell^{\sigma_1}\Phi w
                          -\tau((\gamma-1)w).
\]
Its cyclotomic defect lies in \(X^MN\). The second component is uniquely determined at this precision by
\[
 z=(\Phi-1)^{-1}(\gamma-1)V(w),\qquad z\in X^MN.
\]

Its first component has denominators in finitely many
cyclotomic factors, by the height bounds, and the second
belongs to \(X^MN\). The crystalline criterion proved in
Section~\ref{sec:inert-long-orbit} therefore makes its
integral Herr class finite. A change of representative of
the correction changes both components by a coboundary.
Indeed if the first component of a pair is in \(X^MN\)
along with its second component, the first is the
\((\Phi-1)\)-image of a vector in \(X^MN\). Commutation
and injectivity on \(X^MN\) give the other component of
the coboundary. This also shows that the resulting class
depends linearly on \(w\) and is natural in the marked map.
Writing \([V(w)]\) for the resulting class, the last property of the correction gives
\[
 w\in XN_c\quad\Longrightarrow\quad[V(w)]=0.
\]
Thus integral congruence modulo \(XN_c\) is
sufficient for comparing classes.

For \(w=w_c\), the defect at letter \(c\) is
\(\Phi\Xi_{c-2}-\ell^{\sigma_1}\xi_c\).
Its first summand lies in \(X^MN_c\), on which the
correction is zero. Its second summand has correction
\(-\ell^{\sigma_1}t_c\) modulo the precision.
Thus
\[
 V(w_c)\equiv\ell^{\sigma_1}Y_c\pmod{X^MN}.
\]

The resulting finite classes, denoted by \(\kappa_c\), sum
to \(\ell^{s_1+2\sigma_1}\) times the original class.
The displayed cocycle uses \(\Phi-1\) and \(\gamma-1\)
in this order and has the displayed plus sign. If the
opposite convention for the extension class is used, replace
the sum of the source lifts by its negative.

In particular, if
\[
 \ell^e w_c-f(u_c')\in XN_c(T_2)
\]
for an admissible vector \(u_c'\) in the source, the
finite class attached to \(u_c'\) maps to
\(\ell^e\kappa_c\). All these comparisons take place on
the original rank block and the same map; they do not
require exactness of the Wach functor on the right or a
choice of integral splitting of the surjection.
\label{cor:inert-thick-herr-2}

\paragraph{Filtered comparison at a positive letter.}
\label{par:thick-filtered-comparison}
Write \(i\) for an odd letter, \(b=b_i\) and \(A=B_{i-1}\).
Suppose \(y\in N_{i-2}(T_2)\) has
\(z=\Phi^2\operatorname{sp}_{i-2}(y)\in\mathrm{Fil}^bD_i\).
\begin{proposition}[Integral preimage of the positive filtered class]
In the two-step height and jump range at \(i-2,i-1,i\), put
\[
 E'=\left\lfloor\frac{3a+S+2\eta-1}{\ell-1}
                         +\log_\ell H\right\rfloor .
\]
There exists a vector of the source Wach lattice at \(i\) whose
specialization maps to \(\ell^{E'+\sigma_*+A-b}z\).
The exponent \(E'\le E_G\) is independent of \(y\).
\end{proposition}
\begin{proof}

Use the two-step correction
of Paragraph~\ref{par:thick-integral-correction} to form
\[
 \lambda=\ell^{\sigma_*}\Phi^2 y-
                      \tau((\gamma-1)y)\in P_{i-2}^{(2)}(T_2).
\]
It has cyclotomic defect in \(Xq^{2H}P_{i-2}^{(2)}(T_2)\).
At a primitive root the span is regular in \(t\) and the
constant of \(\lambda\) is \(\ell^{\sigma_*}z\).
The correction has zero constant there because each transported
positive \(X\)-power vanishes at the root and its division modulo
\(Xq^{2H}\) preserves this. By the cyclotomic defect all
positive coefficients below \(2H\) vanish. The correction
likewise has zero specialization at \(X=0\).
These conclusions hold for the same integral representative
at every primitive root.

The signed height at \(i-1\) puts \(p_1^A\lambda\) in
\(P_{i-1}^{(1)}(T_2)\). Set \(x=q^{-b}p_1^A\lambda\).
Then
\[
 q^{b+\eta}x\in N_i(T_2)
\]
by the height at \(i\).
At a primitive root, the coefficients of \(x\) of degree
\(k<-b\) vanish. In degrees \(-b\le k<2H-b\) they lie
in the span of \(z\), hence in \(\mathrm{Fil}^{-k}D_i\).
The filtration is full for \(k\ge 2H-b\ge\eta\), since
all jumps at the odd letter are at least \(-\eta\).
The Hodge lattice test therefore gives \(x\in N_i(T_2)\).
The same signed height at \(i\) applied to the defect
gives
\[
           p_1^A(\gamma-1)\lambda\in Xq^{a+b}N_i(T_2),
\]
because \(2H\ge a+b+\eta\).
In particular
\(\operatorname{sp}_i(x)=\ell^{\sigma_*+A-b}z\).

It remains to lift a controlled multiple of \(x\) in the
integral Wach lattice. Let \(V_q=\gamma(q)/q\),
\(V_1=\gamma(p_1)/p_1=\phi(V_q)\).
Then \(V_q\) is one modulo \(X\) and \(\chi(\gamma)\) modulo \(q\).
The last equality follows by expanding at a primitive root
in the parameter \(t\) of degree one. Also
\(V_1\equiv1\pmod{Xq}\) by the formula for \(\phi(X)\).
The preceding defect equation implies
\[
 \gamma x=V_q^{-b}V_1^A x
                    \pmod{Xq^aN_i(T_2)}.
\]
Thus \(\partial x=d x\bmod q^aN_i(T_2)\) with
\(d=(V_q^{-b}V_1^A-1)/X\) integral and \(d\equiv c_b\bmod q\).

For \(g\in O[[X]]\), \(\gamma(g)\equiv g\bmod q\), and
the divided operator preserves \(q^aN_i(T_2)\).
Applying it to \(q^k g x\) shows that
\[
    (\partial-c_{b-k})q^k g x
           \in q^{k+1}O[[X]]x+q^aN_i(T_2)\qquad(0\le k<a).
\]
Indeed the multiplier on the surviving degree is
\(\chi(\gamma)^k c_b+(\chi(\gamma)^k-1)/(-\ell)=c_{b-k}\).
It follows that \(P_-(\partial)x\in q^aN_i(T_2)\), where
\(P_-(U)=\prod_{j=b-a+1}^{b}(U-c_j)\).
The image input of Paragraph~\ref{par:thick-initial-classes}
gives an integral source preimage of this vector.

For a second filter, first lift \(q^a x\) integrally and divide
the lift by \(q^a\) in the source \'etale module.
The resulting vector \(v\in q^{-a}N_i(T_1)\) maps to \(x\).
The shifted-jump calculation of the first restriction, now
on the odd source, shows that \(P_+(\partial)v\in N_i(T_1)\),
where
\[
 P_+(U)=\prod_{j\in[b+1,b+S+a]\,\cup\,[1-\eta,\eta+a]}(U-c_j).
\]
The two sets of indices for \(P_-\) and \(P_+\) are disjoint
because \(b\ge 2a+\eta\). This integer inequality follows
from \(b\ge 2A_0^{\rm real}+\eta+1\) and
\(a<A_0^{\rm real}+1\), with equality of \(a\) and the
real bound allowed when the latter is integral.

For a root index \(j\) of \(P_-\), the valuation of
\(P_+(c_j)P_-'(c_j)\) is the sum of \(v_\ell(j-j')\)
over the other indices. Indeed, if \(I_-\) and \(I_+\) denote the two index sets, factoring
the polynomials at the root gives
\[
 v_\ell(P_+(c_j)P_-'(c_j))
 =\sum_{j'\in(I_-\cup I_+)\setminus\{j\}} v_\ell(c_j-c_{j'})
 =\sum_{j'\in(I_-\cup I_+)\setminus\{j\}} v_\ell(j-j') .
\]
The indices in or above the first interval give factorials
with total weight \(2a+S-1\). In the central interval the
differences form an interval of length \(a+2\eta\) with
largest absolute value at most \(H\).
Thus the valuation is at most
\[
           \frac{3a+S+2\eta-1}{\ell-1}+\log_\ell H .
\]
Let \(E'\) be the floor of the displayed bound.

For the index set \(J_-\) of \(P_-\) take
\[
 B(U)=\sum_{j\in J_-}
 \frac{\ell^{E'}}{P_+(c_j)P_-'(c_j)}\,\frac{P_-(U)}{U-c_j}.
\]
The valuation bound gives \(B\in O[U]\). Since \(P_-\) is monic
and the numerator vanishes at every \(c_j\), there is \(C\in O[U]\)
with \(BP_++CP_-=\ell^{E'}\). Choose integral source preimages
\(v_+,v_-\) of \(P_+(\partial)x,P_-(\partial)x\), respectively,
as above. Stability under \(\partial\) then gives an integral vector
\(B(\partial)v_++C(\partial)v_-\). By commutation with \(f\), its
image is \(\ell^{E'}x\) and its specialization maps to
\(\ell^{E'+\sigma_*+A-b}z\).

The exponent \(E'\le E_G\) is independent of \(y\)
and of any choice of an integral splitting. This is the
filtered estimate used at each odd letter in the two
even-letter comparisons.
\label{cor:inert-thick-filtered-1}
\end{proof}


\paragraph{The zeroth coefficient at an even letter.}\label{par:thick-constant-data}
Use \(N_i\) for the target Wach lattice at \(i\), in the tame-invariant
coordinate ring. Put \(B_i=b_i+S\) at each letter.
Use the associate \(q=X+\ell\) of the first cyclotomic factor
and \(t\) for the logarithmic parameter at a primitive \(\ell\)-th root.
The parameter \(X\) of the
tame-invariant Wach ring is a unit in the completed de Rham ring at
that root. Write \(D_i\) for the \(i\)-th crystalline space and
\(\operatorname{sp}_i:N_i\to D_i\) for the reduction at \(X=0\).

Let \(\Phi\) also denote the induced Frobenius \(D_i\to D_{i+1}\).
The one-step Wach height at \(c-2\) gives
\[
\begin{aligned}
 P_{c-2}&=XN_{c-2}+q^{B_{c-2}}\langle\Phi N_{c-3}\rangle
                 \ \subset N_{c-2},\\
 \operatorname{sp}_{c-2}(P_{c-2})
       &=\ell^{B_{c-2}}\Phi\,\operatorname{sp}_{c-3}(N_{c-3}),
\end{aligned}
\]
Here the bracket denotes the span over the tame-invariant ring.
The equality follows by reduction at \(X=0\), since \(q(0)=\ell\)
and all scalar constants in \(O\) occur as reductions.
Take
\[
  C=\mathrm{Fil}^{b_{c-1}}D_{c-1},\qquad
  L_\eta=\mathrm{Fil}^{-\eta}D_c,\qquad L_0^{\rm fil}=\mathrm{Fil}^0 D_c .
\]
Both \(D_{c-1}/C\) and \(L_\eta/L_0^{\rm fil}\) have dimension at most
\(m\). The first quotient keeps only jumps in the central interval at
the odd letter. The second keeps only the negative central jumps at
the even letter.

We give the coefficient comparison with its quantitative inputs.
On a submodule \(A'\) of the free family of cocycles, suppose that each
\(w(a)\) is a tame-invariant vector at \(c\), linear in \(a\), with
\((\gamma-1)w(a)\in XN_c\) and
\[
 q^{\eta+B_c}w(a)\in N_c.
\]

Its normalized Laurent coefficients \(d_j(a)\) are obtained by
expanding at one primitive \(\ell\)-th root after dividing by the
expansion of \(q^{\eta+B_c}\), which is \(t^{\eta+B_c}\) times a
unit. They are linear in \(a\). The action of the principal generator
fixes the root and on the \(j\)-th logarithmic coefficient acts by
\(\chi(\gamma)^j\). For \(j\ne 0\), the coefficient of
\((\gamma-1)w(a)\) is in \(\mathrm{Fil}^{-j}D_c\) and
\(\chi(\gamma)^j-1\ne0\). Hence
\[
 d_j(a)\in\mathrm{Fil}^{-j}D_c
\]
for \(j\ne0\).

\paragraph{Transport to the first cyclotomic level.}
Use either Wach module and choose an integral vector \(z_1\)
at \(c-1\). For a positive precision \(D\), put
\(J=\lceil D/(\ell-1)\rceil\) and
\[
 \sigma_D=\lceil 1+(D+1)/(\ell-1)\rceil.
\]

On a Wach basis at \(c-1\), take a polynomial vector \(v\)
of degree \(<J\) with constant coefficient equal to that of
\(z_1\) and with \((\gamma-1)v\) in \(X^J\) times the Wach
module after inverting \(\ell\).
The degree-by-degree construction of
Paragraph~\ref{par:wach-cut-correction} bounds the denominator
in degree \(a\) by \(\ell^{A_a}\), \(A_a=a+v_\ell(a!)\).

Consider \(\ell^{\sigma_D}\Phi v\)
in the Frobenius span at \(c\) modulo \(Xq^D\).
In degree \(a\ge1\) its scalar before the coefficient denominators
is a unit times \(\ell^{\sigma_D}X^a q^{(\ell-1)a}\).
In the expansion \(X^a=X(q-\ell)^{a-1}\), the terms retained
modulo \(Xq^D\) have \(\ell\)-content at least
\(\max(0,\ell a-D)\). Since \(A_a\le\ell a/(\ell-1)\),
\[
 \sigma_D+\max(0,\ell a-D)\ \ge\ A_a .
\]
For \(\ell a\le D\) this follows from
\(\sigma_D\ge D/(\ell-1)\). For \(\ell a\ge D\) it
follows because \(\ell a/(\ell-1)-\ell a+D\) decreases
as \(a\) increases. Thus the transported vector has an
integral representative \(y\) in the Frobenius span modulo
\(Xq^D\) after inverting \(\ell\).

Since \(\phi(X)^J\) is divisible by \(Xq^D\), the representative
has \((\gamma-1)y\) in \(Xq^D\) times this span after inverting
\(\ell\). For the free Frobenius span \(P\), the quotient by the distinguished factor \(Xq^D\) is \(O\)-torsion-free. Consequently
\[
 (Xq^D P[1/\ell])\cap P=Xq^D P.
\]
This returns the rational divisibility to the integral span.
The construction is in the tame-invariant span. The coordinates
of \(y\) at \(X=0\) on the Frobenius basis are
\(\ell^{\sigma_D}\) times the coordinates of \(z_1\) at \(X=0\)
on the preceding basis.

At each primitive \(\ell\)-th root the
Frobenius span evaluates in the crystalline space with
coefficients in \(K_O[[t]]\). Its constant term is obtained
by \(\Phi\) from specialization at the preceding letter.
The \(\gamma-1\) congruence gives zero coefficients for
\(0<j<D\), because
\(\chi(\gamma)^j-1\ne0\) at each such degree.
Thus through degree \(D-1\) the expansion consists of the one
constant \(\ell^{\sigma_D}\Phi\operatorname{sp}_{c-1}(z_1)\).

The filtered crystalline comparison at the intervening odd letter
provides a number \(E'\le E_G\), where
\[
\sigma_2=1+\frac{2H_{\rm loc}}{(\ell-1)^2},\qquad
E_G=\frac{3A_0+S+2\eta}{\ell-1}+L_0 .
\]
For \(z\in \Phi^2\operatorname{sp}_{c-3}(N_{c-3})\cap C\) there
is \(z_1\in N_{1,c-1}\) whose specialization maps to \(\ell^k z\),
where \(k=\lceil E'+\sigma_2\rceil+B_{c-2}-b_{c-1}\).
Here \(N_{1,i}\) denotes the source Wach lattice.
Apply the transport in the source with \(D=B_c\).
The signed height gives \(q^{B_c}y\in N_{1,c}\) and
\((\gamma-1)y\in XN_{1,c}\). Its image has constant coefficient
\(\ell^\kappa\Phi z\) and all nonzero coefficients below \(B_c\)
vanish, with
\[
 \kappa\le (B_c+1)/(\ell-1)+3+E_G+\sigma_2+
                   B_{c-2}-b_{c-1},\qquad
 \kappa=\sigma_{B_c}+k.
\]

\paragraph{The lattice at an even letter.}
\begin{lemma}[Covolume after retaining the central directions]
\label{lem:thick-even-covolume}
At an even letter \(i\) of this word, start with
\(P_i=XN_i+q^{B_i}\langle\Phi N_{i-1}\rangle\).
Adjoining at most \(m\) vectors to \(P_i\) leaves normalized
colength at most \(C_D=rS+2mc_D\) in \(N_i\).
In particular
\[
 m=0\quad\Longrightarrow\quad\ell^S N_i\subseteq P_i.
\]
\end{lemma}

\begin{proof}
At an even \(i\), choose a Wach basis \(e_j\) at \(i-1\) whose
Frobenius specializations give an \(O\)-basis of the specialized
Frobenius lattice adapted to the Hodge filtration at \(i\).
Such a basis exists because the intersections of a lattice
with the subspaces of a flag are saturated. Each successive
quotient is free over \(O\). Lift the resulting basis from
the specialization to the Wach module, where it is a basis
by Nakayama's lemma. Write \(-\alpha_j\) for the jump of the
corresponding crystalline vector. Divide the indices according to the signed jumps of the adapted basis:
\[
 J_0=\{j:|\alpha_j|\le\eta\},\qquad |J_0|\le m,\qquad b_i\le\alpha_j\le B_i\quad(j\notin J_0).
\]


Let \(D=B_i+\eta\). For each central direction \(j\), apply the
first-level transport with precision \(D\) to \(e_j\), obtaining
\(y_j\) in the Frobenius span. The vector \(g_j=q^{\alpha_j}y_j\)
belongs to \(N_i\). To see this, at each root its constant
direction starts in degree \(\alpha_j\) since \(q=t\) times
a unit, and its crystalline direction is in
\(\mathrm{Fil}^{-\alpha_j}\). Its error starts in degree
\[
 D+\alpha_j\ge B_i,
\]
where the filtration is full.
The \(q\)-power denominator is removed by the same integral
Weierstrass test as above.

For a Wach basis \((e_j)\) at \(i\), factoring out the common powers from the columns gives
\[
 \det(q^{B_i}\Phi e_1,\ldots,q^{B_i}\Phi e_r)
       =u q^{rB_i-\sum_j\alpha_j},\qquad u\in O[[X]]^\times .
\] The exponent follows at the
cyclotomic level because the Hodge lattice in the adapted
directions has orders \(\alpha_j\).
The two-sided signed height makes this matrix invertible
after inverting \(q\) over \(O[[X]]\), and \(q=X+\ell\)
generates a prime ideal. These facts give the asserted unit
and power over the integral ring. Leave the noncentral
columns as they are and replace each central one by \(g_j\).
Relative to the old columns the replacement is given by the
transport matrix with the central columns multiplied by
\(q^{\alpha_j-B_i}\). At \(X=0\), before the individual
multipliers the transport matrix has determinant
\(\ell^{m'\sigma_D}\),
where \(m'\le m\) is the number of central columns.

If \(C_{\rm old},C_{\rm new}\) denote the specialized
column matrices, then
\[
 \det C_{\rm new}=\det C_{\rm old}\cdot
       \ell^{m'\sigma_D+\sum_{j\ {\rm central}}(\alpha_j-B_i)}.
\]
Its normalized valuation is therefore
\[
\sum_{j\ {\rm noncentral}}(B_i-\alpha_j)+m'\sigma_D
 \le rS+2mc_D.
\]
Indeed \(D\le H_{\rm loc}\) gives
\(\sigma_D\le\lceil c_D\rceil\le2\lfloor c_D\rfloor\).
The noncentral columns belong to \(P_i\), and adding the
at most \(m\) central classes to \(P_i\) thus leaves a quotient
of normalized length at most \(C_D\).
When \(m=0\) the completed Hodge lattice is contained in
\(q^{b_i}\) times the completed Frobenius span. Weierstrass
saturation gives this inclusion integrally. Multiplication
by \(q^S\) then lands in \(P_i\), and
\(q^S\equiv\ell^S\bmod X\) gives the sharper
\(\ell^S N_i\subset P_i\).
\end{proof}

\paragraph{Recovering the target coefficient integrally.}
\begin{lemma}[An integral representative of the constant coefficient]
Use Lemma~\ref{hyp:alternating-thick-word} at an even letter \(c\).
Let \(w(a)\) be the tame-invariant vector with
\((\gamma-1)w(a)\in XN_c\) and \(q^{\eta+B_c}w(a)\in N_c\),
with the coefficient \(d_0(a)\) and filtered space \(L_\eta\)
of Paragraph~\ref{par:thick-constant-data}. There is
\(x\in N_{c-2}\) such that
\[
 \ell^{e_c}d_0(a)-\Phi^2\operatorname{sp}_{c-2}(x)\in L_\eta,\qquad
 e_c\le\sigma_2+B_{c-1}-b_c/\ell+L_0+3.
\]
\end{lemma}
\begin{proof}
Write \(\overline d\) for the class of \(d_0(a)\) modulo \(L_\eta\)
and \(y_0=B_{c-1}\). The coefficients of \(w(a)\)
in nonzero degrees \(j<b_c\) lie in \(L_\eta\). For \(j\le\eta\)
this follows from \(\mathrm{Fil}^{-j}\subset L_\eta\).
For \(\eta<j<b_c\) the two filtrations agree, by the gap
between the negative and central jumps at the even letter.
Take \(D=2b_c+S+\eta\), \(P=\phi(q)\), and
\(\sigma_H=\lceil 1+D/(\ell-1)^2\rceil\).
We first obtain the integral coefficient relation
\[
 \ell^{\sigma_H}q^\eta P^{y_0}\overline d
       \equiv \sum_j a_j(q)\mu_j \pmod{q^{b_c+\eta}} .
\]
The coefficients \(a_j\) are in \(O[[q]]\) and each \(\mu_j\) is
the image of \(\Phi^2\operatorname{sp}_{c-2}(x_j)\) for some
\(x_j\in N_{c-2}\).

The pole input \(q^\eta w(a)\in\langle\Phi N_{c-1}\rangle\)
and the opposite height at \(c-1\) give
\(q^\eta P^{y_0}w(a)\in\langle\Phi^2 N_{c-2}\rangle\).
Transport an entire truncated horizontal basis and its inverse
from \(c-2\), now by two Frobenius steps. We give the scalar
integrality for both matrices. Truncate at degree
\(J=\lceil D/(\ell-1)\rceil\) so the first omitted degree is at
least the degree needed for the \(D\)-jet test. The inverse has the
same denominator bound \(A_a=a+v_\ell(a!)\) in degree \(a\),
since
\[
 A_{i+j}-A_i-A_j=v_\ell\binom{i+j}{i}\ge0.
\]
Iterating this inequality controls every product in the inverse matrix series.

The degree-\(a\) scalar \((\phi^2 X)^a\) is a unit times
\[
        X^a q^{(\ell-1)a}P^{(\ell-1)a}.
\]
In the \(q\)-parameter \(P=\ell\epsilon+q^\ell v\) with
\(\epsilon,v\in O[[q]]^\times\), by
\(\phi(X)=Xq^{\ell-1}\) times a unit and \(X=q-\ell\).
In the binomial expansion of \(P^{(\ell-1)a}\), an index \(j\)
contributes a factor \(q^{\ell j}\ell^{(\ell-1)a-j}\).
For \(a\ge1\), the term with index \(j\) is divisible by \(Xq^D\) as soon as \((\ell-1)a+\ell j\ge D\). Thus a surviving term has
\[
 j<\frac{D-(\ell-1)a}{\ell},\qquad
 (\ell-1)a-j>(\ell-1)a(1+1/\ell)-D/\ell .
\]
The constant term in the transported matrix has no denominator. The retained terms
have \(\ell\)-content at least
\[
\max\bigl(0,(\ell-1)a(1+1/\ell)-D/\ell\bigr).
\]

Put \(a_*=D/((\ell-1)(\ell+1))\). For \(a\le a_*\) the cost bound
\(\ell a/(\ell-1)\) increases. For \(a\ge a_*\) one subtracts the
linear content term, so the slope becomes
\[
 \frac{\ell}{\ell-1}-(\ell-1)(1+1/\ell)<0\qquad(\ell\ge3).
\]
Extending the cost bound from integral degrees to real arguments, its maximum is attained at the junction of the two linear bounds:
\[
 \max_{a\in\mathbb R_{\ge0}}\left\{\frac{\ell a}{\ell-1}
       -\max\left(0,(\ell-1)a(1+1/\ell)-D/\ell\right)\right\}
 =\frac{\ell a_*}{\ell-1}
 =\frac{\ell D}{(\ell+1)(\ell-1)^2}
 \le \frac{D}{(\ell-1)^2}.
\]

Thus \(\ell^{\sigma_H}\) clears both transported matrices
modulo \(Xq^D\). Denote their integral representatives by \(Y,Y'\),
in the order of the horizontal basis and its inverse. The first
omitted degree contributes nothing because \(Xq^D\mid\phi^2(X)^J\).
For the integral coordinate vector \(\mathbf v\) of
\(q^\eta P^{y_0}w(a)\) in the two-step Frobenius span, put
\begin{equation}\label{eq:thick-two-step-inverse}
 \mathbf a=Y'\mathbf v,\qquad
 Y\mathbf a=YY'\mathbf v\equiv\ell^{2\sigma_H}\mathbf v
                                              \pmod{Xq^D}.
\end{equation}

Evaluate at a primitive \(\ell\)-th root. Two Frobenius steps take
the initial cyclotomic argument to the identity. The columns of
\(Y\) in the Frobenius basis are \(\gamma\)-invariant to precision
\(Xq^D\), as in the one-step construction. Their jets through
order \(D-1\) therefore equal the constants
\(\ell^{\sigma_H}\Phi^2\operatorname{sp}_{c-2}(x_j)\).
In the quotient by \(L_\eta\) these become
\(\ell^{\sigma_H}\mu_j\).

Write a bar for projection to this quotient.
The jump gap proved above gives
\[
 \overline{w(a)}\equiv\overline d\pmod{q^{b_c}},\qquad
 q^\eta P^{y_0}\overline{w(a)}
    \equiv q^\eta P^{y_0}\overline d\pmod{q^{b_c+\eta}} .
\]
These congruences use \((q)=(t)\) after evaluation.
Divide the evaluation of \eqref{eq:thick-two-step-inverse} by \(\ell^{\sigma_H}\) over
the coefficient field. Since \(D\ge b_c+\eta\), it gives the claimed
relation with the integral coordinates \(a_j(q)\) of \(\mathbf a\).
The translation \(q=X+\ell\) in the integral series is valid by
completeness of \(O\).

Let \(L\) be the \(O\)-span of the \(\mu_j\). It is finite free.
The coefficient of \(q^\eta\) shows
\[
 \overline d\in K_O L
\]
because \(P(0)\ne0\). Pass to an \(O\)-basis of \(L\), write
one coordinate of \(\overline d\) as \(\theta\), and divide the
resulting congruence by \(q^\eta\). This gives
\(\ell^{\sigma_H}\theta P^{y_0}\) integral modulo \(q^{b_c}\).
Put \(j_*=\lfloor(b_c-1)/\ell\rfloor\le y_0\), the inequality
following from \(b_{c-1}\ge h/2\). Among
\(j=0,\ldots,j_*\), choose the least index \(j_1\) minimizing
\(v_\ell\binom{y_0}{j}+y_0-j\).

Larger indices contribute only at degrees \(>\ell j_1\),
and smaller ones have strictly larger valuation by the minimal
choice of \(j_1\). The coefficient from \(j_1\) is a unit
multiple of \(\binom{y_0}{j_1}\ell^{y_0-j_1}\), since the
constant terms of the two units are units. Thus
\[
 \begin{aligned}
 v_\ell([q^{\ell j_1}]P^{y_0})
   &= y_0-j_1+v_\ell\binom{y_0}{j_1}\\
   &\le y_0-j_*+v_\ell\binom{y_0}{j_*}
    \le y_0-b_c/\ell+1+\log_\ell y_0 .
 \end{aligned}
\]
The last inequality uses the factorial floor sum and
\(j_*=\lfloor(b_c-1)/\ell\rfloor\).
It follows that every coordinate of \(\ell^{e_c}\overline d\)
is integral for
\[
e_c=\sigma_H+\lceil y_0-b_c/\ell+2+\log_\ell y_0\rceil .
\]
Lifting the resulting \(O\)-linear combination of the \(\mu_j\)
gives \(x\in N_{c-2}\) with
\[
 \ell^{e_c}d_0(a)-\Phi^2\operatorname{sp}_{c-2}(x)\in L_\eta .
\]
Since \(\sigma_H\le\sigma_2+1\) and \(y_0\le H_{\rm loc}\),
this exponent satisfies
\[
 e_c\le \sigma_2+B_{c-1}-b_c/\ell+L_0+3 .
\]

\end{proof}
\paragraph{Restrictions on the coefficients at the even letters.}
\begin{lemma}[Reduction to a Wach residual]
At an even letter let \(N_g,N_d\) count the restrictions by generator number and by length, respectively. To reduce to an integral Wach residual it suffices to take
\[
 N_g=3,\qquad N_d=1,
\]
with each of the three generator quotients generated by at most \(m\) elements. For \(m=0\) the reduction holds on the
whole family at exponent \(E_{\rm lev}+S\).
\end{lemma}
\begin{proof}
The height cost cancels before taking a bound over the letters.
Indeed \(e_c+\kappa-B_{c-2}\le B_{\rm lev}\), where
\[
 B_{\rm lev}=\frac{h}{\ell(\ell-1)}+\frac{S+1}{\ell-1}
              +S+E_G+2\sigma_2+6+L_0 .
\]
The calculation uses
\((B_c+1)/(\ell-1)-b_c/\ell
 =b_c/(\ell(\ell-1))+(S+1)/(\ell-1)\).
Put \(E_{\rm lev}=\lfloor B_{\rm lev}\rfloor\).
As \(e_c+\kappa-B_{c-2}\) is integral, the bound gives
\[
  E_{\rm lev}-e_c-\kappa+B_{c-2}\ \ge0,\qquad
  \ell^{E_{\rm lev}}d_0(a)
   \equiv \ell^{\kappa-B_{c-2}}\Phi^2\operatorname{sp}_{c-2}(x)
                      \pmod{L_\eta}
\]
for \(x\in N_{c-2}\). The correction is obtained by multiplying
the earlier choice of \(x\) by the nonnegative power
\(\ell^{E_{\rm lev}-e_c-\kappa+B_{c-2}}\).

Let \(K_\eta\) be the kernel of
\(\Phi^2\operatorname{sp}_{c-2}\bmod L_\eta\) on \(N_{c-2}\).
Multiplication by \(\ell^{\kappa-B_{c-2}}\) is invertible
on the coefficient-field quotient. Hence two choices of \(x\)
differ by \(K_\eta\), and addition of congruences gives an
\(O\)-linear map to \(N_{c-2}/K_\eta\). Composing with the
further quotient gives
\[
 A'\ \longrightarrow\ N_{c-2}/K_\eta\
       \longrightarrow\ N_{c-2}/(P_{c-2}+K_\eta).
\]
Write \(Q=N_{c-2}/(P_{c-2}+K_\eta)\) and \(I\subset Q\) for
the image of \(A'\). The kernel consists of the elements for which
one can choose \(x\in P_{c-2}\).

By Lemma~\ref{lem:thick-even-covolume}, there is a submodule
\(U\subset Q\) generated by at most \(m\) distinguished classes
with normalized length \(\ell_{\rm norm}(Q/U)\le C_D\).
The two restrictions on \(A'\) have successive quotients read from
\begin{equation}\label{eq:thick-restriction-quotients}
 0\longrightarrow I\cap U\longrightarrow I
       \longrightarrow (I+U)/U\longrightarrow0,\qquad
 \ell_{\rm norm}((I+U)/U)\le C_D .
\end{equation}
For the first term of \eqref{eq:thick-restriction-quotients}, the preimage of \(I\cap U\) in a free cover \(O^m\to U\) is free
of rank at most \(m\). Its image generates \(I\cap U\); killing that
image costs one generator restriction.

On the resulting kernel the new coefficients satisfy
\[
 \ell^{E_{\rm lev}}d_0(a)\equiv\ell^\kappa\Phi z\pmod{L_\eta},
 \qquad z\in Z:=\Phi^2\operatorname{sp}_{c-3}(N_{c-3}).
\]
Set \(K'=\ker(Z\xrightarrow{\,\Phi\,}D_c/L_\eta)\), so \(z\) is determined modulo
\(K'\). Requiring a representative in \(C\) costs at most \(m\)
generators: the quotient \(Z/((Z\cap C)+K')\) is an image of
\[
 Z/(Z\cap C)\hookrightarrow D_{c-1}/C,\qquad
 \dim_K(D_{c-1}/C)\le m ,
\]
where \(K\) denotes the coefficient field. The source of the
injection is finitely generated over \(O\), since \(Z\) comes from
a Wach lattice modulo \(X\). It is therefore free of rank at most
\(m\), and submodules of its quotient also need at most \(m\)
generators. For \(m=0\) the inclusion
\[
 \ell^S N_{c-2}\subset P_{c-2}
\]
allows \(\ell^S x\) as representative and
\(C=D_{c-1}\).

Choose the source vectors \(y\) from the first-level construction on
an \(O\)-basis of the restricted family, then extend linearly.
The constant of \(\ell^{E_{\rm lev}}w(a)-f(y(a))\) lies in
\(L_\eta\). Requiring it to lie in \(L_0^{\rm fil}\) costs one more
generator restriction of size at most \(m\). When \(m=0\) these
two filtration steps agree. For \(j\ne0\) the difference of
the coefficients lies in \(\mathrm{Fil}^{-j}D_c\), since the
source image has vanishing coefficients for \(j<B_c\) and
\[
 \mathrm{Fil}^{-j}D_c=D_c
\]
for \(j\ge B_c\).
Tame transport carries this property from the chosen root to
all primitive \(\ell\)-th roots. The rational cyclotomic
evaluation test now puts the vector difference in \(N_c\).

For the integral test, start with the known denominator
\(q^{\eta+B_c}\) and divide its coordinates in a Wach basis by the
distinguished polynomial associated to that power. At each root the
required coefficient vanishing follows from membership in the
completed Hodge lattice. The remainder is therefore zero by
Weierstrass division. This proves the reduction to a Wach vector
residual with three generator restrictions and one length restriction.
For \(m=0\) it works on the whole family with exponent
\(E_{\rm lev}+S\).

\end{proof}
\paragraph{The residual modulo \(X\).}
\begin{lemma}[Lifting the remaining class modulo \(X\)]
Three generator restrictions and two length restrictions lift the residual modulo \(XN_c\) after multiplication by \(\ell^{e_6}\), where
\[
 e_6=S+\lceil E_G+\sigma_2\rceil.
\] If \(m=0\), exponent
\(e_6+2S\) suffices on the entire family.
\end{lemma}
\begin{proof}
Use the same filtered crystalline comparison and
Lemma~\ref{lem:thick-even-covolume} at both even letters \(c,c-2\). Let \(v\) be the
integral Wach residual on the family.
Restrict first to \(v\in P_c\), costing one length and one
generator restriction, as above. There is a vector \(v_1\in N_{c-1}\)
with \(\operatorname{sp}_c(v)=\ell^{B_c}\Phi\operatorname{sp}_{c-1}(v_1)\).
At the odd letter the opposite Wach height gives
\(q^{B_{c-1}}N_{c-1}\subset\langle\Phi N_{c-2}\rangle\).
Specializing at \(X=0\) gives \(v_2\in N_{c-2}\) with
\[
 \ell^{B_{c-1}}\operatorname{sp}_{c-1}(v_1) =\Phi\operatorname{sp}_{c-2}(v_2).
\]

The specializations of \(v_1,v_2\) are linear in the family,
because Frobenius is invertible on the crystalline spaces and the
scaling factors are nonzero. The kernel of specialization is
\(XN_{c-2}\), so the map \(v_2\bmod P_{c-2}\) is well defined.
Restricting it to zero costs one length and one generator
restriction. Then for some \(v_3\in N_{c-3}\) one has
\(\operatorname{sp}_{c-2}(v_2)
=\ell^{B_{c-2}}\Phi\operatorname{sp}_{c-3}(v_3)\).
Restrict to the kernel of the image of
\(z=\Phi^2\operatorname{sp}_{c-3}(v_3)\) in \(D_{c-1}/C\).
This is one generator restriction of size at most \(m\).

Now \(e_6=S+\lceil E_G+\sigma_2\rceil\) satisfies
\[
 e_6+B_{c-2}=u_0+k+B_{c-1},\qquad
 u_0=\lceil E_G+\sigma_2\rceil-\lceil E'+\sigma_2\rceil\ge0 .
\]
Choose \(z_1\in N_{1,c-1}\) with
\(f(\operatorname{sp}_{c-1}z_1)=\ell^k z\).
Combining the three specialized equalities gives
\[
 \operatorname{sp}_c(\ell^{e_6}v)
   = f\bigl(\operatorname{sp}_c(\ell^{u_0}q^{B_c}\Phi z_1)\bigr).
\]
The vector \(\ell^{u_0}q^{B_c}\Phi z_1\) belongs to \(N_{1,c}\)
by the source height, so it is tame invariant and has
\(\gamma-1\) difference in \(XN_{1,c}\).
The residual is therefore lifted modulo \(XN_c\) at the cost of
three generator and two length restrictions. When \(m=0\), one
uses \(\ell^S N_i\subset P_i\) twice and \(C=D_{c-1}\).
This gives a lift on the entire family with exponent \(e_6+2S\).

\end{proof}
\begin{proposition}[Finite-class loss on an alternating word of larger jumps]\label{prop:thick-inert-colength}
For the marked crystalline surjection and signed jump word of
Lemma~\ref{hyp:alternating-thick-word}, the chain uses at most \(4n\) generator restrictions of
size \(m\) and \(3n/2\) length steps. Let \(A_{\rm end}\) be the last submodule in the restricted chain and
\(\pi:A\to M_2=H^1_f(F,T_2)\) the chosen surjection. There is an exponent \(e\) such that
\[
 \ell^e\pi(A_{\rm end})\subseteq f_*H^1_f(F,T_1),\qquad
 e\le \frac{32H_{\rm loc}}{\ell^2}+4S+\frac{10\eta}{\ell-1}+25L_0 .
\]
For \(m=0\) the construction gives a lift on the full finite
group with the additional allowance \(3S\). With the separately
available coarse lift, \eqref{eq:thick-chain-mass} converts these
losses to the normalized place colength.
\end{proposition}

\begin{proof}
Paragraphs~\ref{par:thick-initial-classes} and
\ref{par:thick-integral-correction} decompose the target finite
classes after \(s=s_1+2\sigma_1\), with at most \(n\) generator
restrictions of size \(m\). They realize the vectors \(w_c\)
on the target and compare their finite classes under the
same marked map modulo \(XN_c\).

\proofstep{Restrict the same free family at all even letters, then sum the cut classes.}
Now apply the letter
comparisons successively on the submodules obtained so far.
They are stable under restriction
of the family. The first lift \(y\) and the lift of the residual
combine by adding \(\ell^{e_6}y\) to the latter source vector. There
are \(n/2\) even letters. Counting at every other letter gives the respective maxima
\[
 n+6(n/2)=4n\quad\text{generator steps},\qquad 3(n/2)=3n/2\quad\text{length steps}.
\]
Each cut class has a finite lift after the same power.
Summing their classes lifts the scaled original class, since the
decomposition holds modulo coboundaries. This costs
\(e=s+E_{\rm lev}+e_6\), independently of the number of letters.
For \(m=0\) the extra powers on the two comparisons give \(e+3S\)
on the entire finite group.

\proofstep{Combine the decomposition and residual losses, using \(\ell\ge11\).}
The successive restrictions have been imposed on one family of lifts. Their losses can now be added without counting a different family at each even position.

Here is the numerical bound from this construction. Put
\(u=1/(\ell-1)\) and
\(\sigma=1+(\lfloor hu\rfloor+1+\eta)u\).
The resulting bounds for the decomposition and its remainder are
\[
\begin{aligned}
 s&\le(2A_0+2\eta+S)u+L_0+2\sigma+2,\\
 e_6&\le S+E_G+\sigma_2+L_0 .
\end{aligned}
\]
Adding these bounds and the bound for \(E_{\rm lev}\), using
\(\sigma\le1+hu^2+u+\eta u\), gives
\[
 e\le 8A_0u+8\eta u+(4S+3)u+9H_{\rm loc}u^2+2S+5L_0+13 .
\]
Since \(u\le1/10\) and
\(A_0u\le (2+u)H_{\rm loc}u^2+(4+L_0)u\), we obtain
\[
 e\le25.8H_{\rm loc}u^2+8\eta u+2.4S+5.8L_0+16.5
   \le \frac{32H_{\rm loc}}{\ell^2}+4S+
              \frac{10\eta}{\ell-1}+25L_0 ,
\]
using \(L_0\ge1\) and \(\ell/(\ell-1)\le 11/10\).
With the common coarse finite-class lift \(e'\),
\eqref{eq:thick-chain-mass} gives the bound for sequences with larger jumps.

The target coefficient bound uses the estimate on the rank of the kernel. The filtered crystalline estimate gives the realization at the first level. Apply the two lattice bounds for even letters in their respective cases.
 Combine the squared Frobenius bound on the residual lattice for an even letter with the coefficient lift.
 The resulting lattice bounds give the length at the chosen prime.
The local argument applies the integral Wach and Herr comparisons
to the marked crystalline models and uses the corresponding
coarse class lift for the generator restrictions.
\end{proof}


The common entry in cases C and D is
Proposition~\ref{prop:ordinary-inert-place-entry}.

\subsubsection{Smaller jumps and the two estimates at ordinary inert primes}
\label{sec:inert-row-table}
\hypertarget{source-inert-row-table}{}

When smaller jumps occur, the uniform restriction used for the larger jumps does not cover the whole sequence. Decomposing it into short arcs gives the missing estimate; it can then be combined with the larger-jump case at the same inert place.

Here \(\ell\ge11\) and \(n=n_\lambda\) is the local unramified degree. In the
case with some smaller jumps, let \(s_i\) be the sizes in the local word.
Use \(b_0,w,S,\eta\) for its threshold for the larger jumps, bound for the smaller jumps,
spread and central width. Each
\[
 s_i>\eta.
\]
A size in the larger range is
in \([b_0,b_0+\eta]\), and a size in the smaller range is at most \(w\).
Use the signed jump intervals of Section~\ref{subsec:inert-rows}.

At every letter at most \(m\)
target jumps, counted with multiplicity, lie in \([-\eta,\eta]\).
On the marked row the signed windows come from the source lattice
in the sum of twisted duals. The bound on the central multiplicity
comes from the twisted dual target at each labeled embedding,
applied with the same enumeration of the embeddings as the word.
Both statements therefore apply to a size in the smaller range.
\label{cor:inert-row-table-1}
There is a position with smaller jumps, the signs alternate cyclically,
and there are at most two consecutive positions with larger jumps.
The numerical hypotheses are
\[
 20w,20S,20\eta\le b_0.
\]

Put \(h=b_0+\eta+S\) if there is a position with larger jumps and \(h=w+S\)
otherwise, and put
\[
 \gamma_h=\log_\ell(2h),\qquad
 S_g=h/\ell^2+(w+S+\eta)/\ell+1+\gamma_h .
\]
Let \(B\) be the indicator that some position with smaller jumps \(i\)
has a neighbor with larger jumps and satisfies
\[
 \abs{\ell s_i-b_0}\le 320
              (\ell\eta+\ell S+b_0/\ell+s_i+\ell\gamma_h).
\]
The common place indicator in \(F_{\ell,\lambda}\) is the
existence of such a word on the marked data, over all good
characters and compatible local coefficient data. Therefore every
local \(B\) is at most this place indicator.

Cut before each position with smaller jumps, numbering the cut by
its preceding letter. Begin the enumeration at one such position, so
the construction includes the arc that crosses the end of the cycle.
Let \(\mathcal C\) be the set of cuts. The arc ending at \(c\)
consists of one smaller letter followed by \(k_c\) larger letters.
Consequently
\[
 a_c=1+k_c,\qquad 0\le k_c\le2,\qquad
 \sum_{c\in\mathcal C}a_c=n .
\]

Suppose first that a larger letter occurs and \(40S_g>b_0/4\).
The two jump bounds of the coarse lift are at most \(h\); its
maximum width is at most \(h+\eta\). The coarse exponent calculation
in the proof of Proposition~\ref{prop:ordinary-inert-place-entry}
therefore bounds the exponent \(e'\) of a full finite lift by
\(7.1 h/\ell+7.1+2\gamma_h\). Since \(h\le1.1b_0\),
\(b_0<160S_g\), and \(\ell\ge11\),
\[
 7.1 h/\ell <113.6 S_g,\qquad
 7.1+2\gamma_h\le7.1S_g,\qquad e'<160 S_g .
\]

In the remaining range, use Lemmas~\ref{lem:two-letter-wach-lift},
\ref{lem:three-letter-negative-lift} and~\ref{lem:three-letter-positive-lift}
on arcs of length two and three. To compare the arc lift with the cut
decomposition we need the exponents \(R+2\sigma\) and \(R+3\sigma\).
Put \(g=w+S+\eta\), \(q_0=\ell-1\), and
\(M=\lfloor2h/q_0\rfloor+1\). There is one smaller letter
immediately after the cut. The corresponding correction satisfies
\[
 \begin{aligned}
 \sigma
 &=\left\lceil1+\frac{\ell(2(w+S)+M)}{q_0^2}\right\rceil\\
 &\le 2+\frac{\ell}{q_0^2}
       +\frac{2\ell g}{q_0^2}+\frac{2\ell h}{q_0^3}
 \le 3\left(1+\frac{g}{\ell}+\frac{h}{\ell^2}\right).
 \end{aligned}
\]
Indeed \(\ell/q_0\le11/10\) bounds the two variable coefficients
after normalization by \(2(11/10)^2\) and \(2(11/10)^3\).
For the recurrence of Section~\ref{sec:inert-long-orbit}, take
symmetric widths at most \(2h\) throughout and at most \(2g\)
at the first letter after the cut. With
\(c_H=\lceil1+\log_\ell(1+\max_i H_i)\rceil\) one has
\[
 R\le c_H + \frac{\ell}{q_0(q_0-1)}
                    (2g+c_H+3+2h/q_0).
\]

Here \(c_H\le 2+\tfrac15+\gamma_h\) by
\(1+2h\le3h\) and \((3/2)^5<11\).
The coefficients of \(g/\ell\) and \(h/\ell^2\) in the bound
are at most \(2\cdot11^2/(10\cdot9)<3\) and
\(2\cdot11^3/(10^2\cdot9)<3\), respectively. Since
\(\ell/(q_0(q_0-1))\le11/90\), the remaining terms
are at most \(3(1+\gamma_h)\). Combining these estimates gives
\begin{equation}\label{eq:inert-short-cut-budget}
 R\le3S_g,\qquad \sigma\le3S_g,\qquad
 R+2\sigma\le9S_g,\quad R+3\sigma\le12S_g .
\end{equation}
With the cut costs of \eqref{eq:inert-short-cut-budget}, the arc lift at \(c\) itself has exponent at most
\[
 \sigma+\kappa_c+\lfloor 40S_g+\eta B\rfloor ,
 \qquad
 \kappa_c\le 13S+3(1+\gamma_h).
\]
It lands in the finite image modulo a span of at most
\(m r_c(a_c+1)\) target finite classes, with \(r_c\le a_c\).
We derive the coordinate lift and the fixed exceptional span below.

\paragraph{Coordinates on a short arc.}\label{par:short-wach-frame}
\begin{proposition}[Common coordinates for short-arc comparisons]\label{prop:short-arc-coordinates}
For a short arc, use the same integral vertical and
horizontal frames and the corresponding cyclotomic levels.
For an arc ending at \(c\) of length \(a=a_c\), the offset \(0\le j<a\)
has size \(u_j=s_{c-j}\) and cyclotomic factor \(P_j=\phi^j(X+\ell)\). The horizontal and vertical comparisons use the same positive pole and negative zero factors
\[
 D=\prod_{j\ {\rm positive}}P_j^{u_j},\qquad V=\prod_{j\ {\rm negative}}P_j^{u_j+S}.
\]
The two ordinary lifting rules use these coordinates. In the single-letter edge cases the
positive rule applies to every input, and the negative rule
uses the empty pole set and the initial constant lift.
\end{proposition}

\begin{proof}
\proofstep{Fix the same cut coordinate and ordered cyclotomic levels for both rules.}
Write \(a=a_c,\ c'=c-a,\ G=w+S,\ u_j=s_{c-j}\), and
\(P_j=\phi^j(q)\), so \(P_0=q=X+\ell\).
Let \(\mathcal R=O[[X]]=O[[q]]\). All ideals of fractions below
are in the component at \(c\).
Call an offset positive or negative according to its large Hodge
jumps. This sign is used even at the first letter of small size.
The full denominator for transport over \(a\) steps and the two
smaller products determined by the jumps are
\[
 \Lambda=P_{a-1}^{G}\prod_{j<a-1}P_j^h,\qquad
 D=\prod_{j\ {\rm positive}}P_j^{u_j},\qquad
 V=\prod_{j\ {\rm negative}}P_j^{u_j+S}.
\]
Choose a basis of the tame-invariant Wach lattice at \(c'\).
The unique formal \(\gamma\)-fixed basis with the same constants is
its horizontal basis.

Write \(U(X)=I+\sum_{r>0}U_rX^r\) for the change
to the horizontal basis, and let \(F_r(U_1,\ldots,U_{r-1})\)
be the degree-\(r\) residual computed with \(U_r=0\).
The recursion is
\[
       (\omega_X^r-1)U_r=-F_r(U_1,\ldots,U_{r-1}).
\]
Its new denominator has order \(1+v_\ell(r)\) at step \(r\).
Thus through that step the denominators divide
\(\ell^{r+v_\ell(r!)}\). This also holds for the inverse change
of basis, by the factorial inequality for products of coefficients.
Use the basis truncated before degree \(M\). These are coordinates
on the given lattice map: a map of Wach lattices acts on horizontal
columns by its matrix of constants at \(c'\), since both sides solve
the same equation with the same constant. Frobenius transports these
constants to the crystalline space at \(c\).

Here are details on the precision of these coordinates. Put
\[
 P(T)=\langle\Phi^aN_{c'}(T)\rangle_{\mathcal R}.
\]

Represent \(v\in\Lambda^{-1}P(T)\) modulo \(X\Lambda P(T)\)
in the transported horizontal basis, allowing powers of \(\ell\)
to clear coefficients in this comparison. Equivalently multiply
by \(\Lambda\) and compare modulo \(X\Lambda^2 P(T)\).
The tails beyond the truncation have a factor
\(\phi^a(X)^M\), divisible by \(X\Lambda^2\) since
\((\ell-1)M>2h\). The finite coefficients, and the inverse matrix,
can be transported integrally at cost \(\sigma\). Indeed on
the first step one reduces \(\ell^\sigma\phi(X)^r/\ell^{r+v_\ell(r!)}\)
modulo \(X^Mq^{2G}\).

The term \(X^r q^{(\ell-1)r}\) pays the
denominator after this precision, by the binomial division of
Paragraph~\ref{par:wach-cut-correction} with one forward step and
precision \(2G\). More explicitly, for \(k=r+v_\ell(r!)-\sigma>0\)
one needs
\[
 (\ell-1)r\ge 2G+k+M-r-1.
\]
Then modulo \(\ell^k\)
the binomial expansion of \((X+\ell)^{(\ell-1)r-2G}\) is divisible
by \(X^{M-r}\). If \(k>0\), the coefficient bound gives
\(r>(\ell-1)\sigma/\ell\). Together with the definition of
\(\sigma\) this yields the preceding power inequality.

In
the binomial expansion the terms of \(\ell\)-degree \(\ge k\)
vanish modulo \(\ell^k\); every other term has \(X\)-degree
at least \((\ell-1)r-2G-k+1\ge M-r\). Thus the termwise
division is legitimate at the first precision. The next
\(a-1\) Frobenius maps give
\[
             \phi^{a-1}\bigl((X^Mq^{2G})\bigr)
                                     \subseteq (X\Lambda^2)
\]
by the same level bounds defining \(\Lambda\). Thus the transported matrices are represented by
integral lifts of \(\ell^\sigma\phi^a(U)\) and
\(\ell^\sigma\phi^a(U^{-1})\) modulo \(X\Lambda^2\)
in the Frobenius basis. If \(A,B\) denote these two integral lifts in the displayed order, their common precision gives
\[
 AB\equiv BA\equiv\ell^{2\sigma}I\pmod{X\Lambda^2}.
\]


For clarity, the coefficient test in these coordinates is as
follows. If a constant direction \(d\) at the end of the arc
belongs to \(\mathcal F_j^{-\alpha_j}\) for each transported
filtration, then \(d\) with coefficient in
\((\prod P_j^{\alpha_j})\mathcal R\) passes the Wach test at
each level. Sums of such terms after scalar extension give the
rational test; the bounded integral numerators are required as
well.

To see the rational assertion, separate the levels by the
Chinese remainder theorem for their distinguished polynomials
and take a basis of the filtration at each level. For a basis vector \(d\) of jump \(u\), the corresponding generator and its Laurent order are
\[
 t^{-u}d,\qquad \operatorname{ord}_t(t^{-u})=-u,
\]
as in
Paragraph~\ref{par:wach-level-models}. Multiplying the other
level factors to sufficiently high orders realizes each of these
finite jets without imposing further conditions there.
Divisibility of bounded numerators by a product of these
polynomials descends to \(\mathcal R\) by monic division.

Let \(u\) denote a coordinate vector satisfying all the
local tests and let \(n\) be a vector in \(N_c\).
The forward comparison and its inverse give, in their
respective quotients by the transported precision,
\[
   f_{\rm hor}(u)=[\ell^\sigma u]_{\rm Wach},
   \qquad g_{\rm hor}(n)=[\ell^\sigma n]_{\rm coordinates},
   \qquad g_{\rm hor}f_{\rm hor}=\ell^{2\sigma}.
\]
Thus an admissible coordinate vector produces a Wach class,
and the restricted inverse sends a Wach vector back into the
tested coordinate module. Indeed the transported forward
and inverse matrices above give integral lifts with the bounded
denominator, and changes by \(X\Lambda P(T)\) lie in \(XN_c\)
by the height bounds.

For an admissible vector one tests
\((\gamma-1)/X\) in the same manner. This gives an integral
vector with \((\gamma-1)v\in XN_c\), since its horizontal
columns are fixed to the indicated precision. A product of the
\(P_j\) clears its denominators in \(N_c\). These statements
apply to subarcs beginning at \(c'\) with their own product
\(\Lambda\). \label{cor:inert-short-frame-1}

We will lift a coordinate \(f\) on a target horizontal column
\(d\) which belongs to \(\mathcal F_j^{u_j}\) at every positive
level. Its scalar hypotheses are
\begin{equation}\label{eq:short-scalar}
 Df\in\mathcal R,\qquad
 D(\ell^{k_j}f-x_j)\in P_j^{u_j+S}\mathcal R
 \quad(j\ {\rm negative})
\end{equation}
for some \(k_j\ge0,\ x_j\in O\). A lift at cost \(e\) represents
\(\ell^e f d\) by the image of an admissible source vector up to
\(XN_c(T_2)\) and the coordinate precision. There are two ways
to construct it.
\begin{enumerate}
\item If \(Df\in V\mathcal R\), then \(fd\) passes the target
Wach test. The poles at the positive levels are allowed by the
filtration of \(d\); the zeros at the negative levels accommodate
the lowest possible jump. The coordinate comparison and scalar image bound at \(c\) use successive multipliers
\[
 \ell^\sigma\ell^R=\ell^{R+\sigma}.
\]
Tame averaging preserves the lifted image, so no extra exponent enters this case.
\item Suppose the poles are confined to a subset of the positive
levels and an integral source direction maps to \(\ell^t d\) and
lies in \(\mathcal F_j^{u_j}\) there. At each omitted positive
level also require that \(f\) is constant modulo \(P_j^\eta\)
after clearing scalars and the retained denominator. Then the
coefficient \((\gamma f-f)/X\) times this source direction passes
the level tests. For each level factor, write \(\gamma P_j/P_j=1+Xh_j\) with \(h_j\in\mathcal R\). Then
\[
 (\gamma P_j-P_j)/X=h_jP_j.
\]

For every \(n\in\mathbb Z\) the ratio \((1+Xh_j)^n\) belongs to \(1+X\mathcal R\). Thus
\[
 (\gamma P_j^n-P_j^n)/X\in P_j^n\mathcal R ,
\]
including for negative powers. The constant disappears
at each omitted positive or negative level. Equation~\eqref{eq:short-scalar}
therefore supplies the required zeros at negative levels, and
\(\eta\) supplies them at omitted positive levels. The retained
positive poles pass by the source filtration. The admissible-vector
construction gives the lift at \(t+\sigma\). For an integral constant, transport the constant term of the initial target basis lift at \(c'\). Its coefficient differences vanish, so its cost is at most \(R+\sigma\).
\end{enumerate}

In the second rule the congruences use the denominator for the
retained positive levels. Costs of two summands may be padded
to their maximum, by linearity of the source map and of the
coordinate comparison. For a single positive letter the first
rule applies to every \(f\). For a single negative letter use
the second with empty pole set and the initial constant lift.
\end{proof}


\begin{auxlemma}[Separating a scalar by Weierstrass division]\label{lem:short-arc-content}
Let \(O\) be a coefficient DVR with fraction field \(K\),
\(v_p(p)=1\), and let \(J\subset\mathcal R=O[[q]]\) be generated by
a distinguished polynomial. Take \(D\in\mathcal R\) and
\(f\in\operatorname{Frac}(\mathcal R)\). Assuming that the residue
\(\bar D\in\mathcal R/J\) is nonzero, denote by \(c_J(D)\) the minimum of the
valuations of its Weierstrass remainder coefficients. Let \(k,d\)
be nonnegative integers with \(c_J(D)\le d\). If
\(Df\in\mathcal R\), \(x\in K\), and \(p^k Df-Dx\in J\), then there is an integral constant satisfying
\[
 \exists x'\in O:\qquad D(p^d f-x')\in J.
\]
\end{auxlemma}
\begin{proof}
Write \(H=\mathcal R/J\), a free \(O\)-module with the monomials
below \(\deg J\) as basis, where \(\deg J\) is the degree of the
distinguished generator. In \(H\otimes_O K\) the given equality
implies
\[
 x\bar D=p^k\overline{Df}\in p^k H,\qquad
 v_p(x)+c_J(D)\ge k .
\]
Thus \(x'=p^{d-k}x\) is integral, including when \(k>d\).
Both \(p^d Df\) and \(x'D\) lie in \(\mathcal R\), and their
images agree after tensoring \(H\) with \(K\). Freeness makes
\(H\to H\otimes_O K\) injective, which proves the assertion.
\end{proof}

\begin{lemma}[Two-letter lifting on a short arc]\label{lem:two-letter-wach-lift}
In the models of Proposition~\ref{prop:short-arc-coordinates},
write \(e_2\) for the exponent of a two-letter correction. One can achieve
\[
 e_2\le15L+1+\eta B,
\]
where \(L\) is the local size from the short-arc comparison and \(B\) its
exceptional indicator.
\end{lemma}

\begin{proof}
\label{par:short-two-levels}
Use \(u=s_{c-1}\) for the small size and \(v=s_c\) for the large
one. Put \(p=\ell,\ L=S_g\) within this calculation. Denote by
\(A_i\) the ceiling of the least ordinary Wach image exponent in
\(\pi\) at \(i\), divided by \(p-1\) before taking the ceiling.
Then \(X^{A_i}\) kills the tame-invariant cokernel. In signed
boxes of widths \(s_i+S+\eta\), the recurrence of
Paragraph~\ref{par:wach-cut-colength} gives
\begin{equation}\label{eq:short-A}
 A_i\le(s_{i+1}+S+\eta)/p+\tfrac52 L.
\end{equation}
Here one can replace \(5/2\) by \(2\) if \(i+2\) is small.

\proofstep{Compute the denominators of the two neighboring level evaluations.}
To check the constants, write \(d=p-1,\ H=h+\eta\),
\(W_j=s_j+S+\eta\) and \(c_H=\lceil1+\log_p(1+H)\rceil\).
The ordinary image exponents satisfy
\(s_i^{\rm im}\le W_{i+1}+c_H+3+s_{i+1}^{\rm im}/d\),
and are at most \(H\) by the signed twist. Thus
\[
 A_i-W_{i+1}/p\le 1+(c_H+3)/d+H/(pd)+H/d^2.
\]
The general bound uses \(H=h+\eta\), \(H+1\le2h\), and
\(p/d\le1.1\).

If \(i+2\) is small, iterating the recurrence
once replaces its last summand by
\[
 \frac{H}{d^2}\quad\rightsquigarrow\quad
 \frac{w+S+\eta+c_H+3}{d^2}+\frac{H}{d^3}.
\]
The respective coefficients of \(h/p^2\) are bounded by
\[
 \frac{p}{d}+\frac{p^2}{d^2}\le2.31,\qquad
 \frac{p}{d}+\frac{p^2}{d^3}\le1.24 .
\]
The \(\eta/p^2\) terms are absorbed in \((w+S+\eta)/p\).
Together with \(c_H\le2+\log_p(2h)\), these give the factors
\(5/2\) and \(2\) in the two versions of \eqref{eq:short-A}.
\label{cor:inert-short-two-1}

Apply Lemma~\ref{lem:short-arc-content} with
\(P=P_1=p\varepsilon+Uq^p\), where \(\varepsilon,U\in\mathcal R^\times\).
We estimate the contents for the two possible denominators.
For \(D=q^v,J=(P^y)\), put \(t=\lfloor v/p\rfloor+1\).
Multiply by \(q^{pt-v}\) and expand
\[
 (q^p)^t=U^{-t}(P-p\varepsilon)^t.
\]
Modulo \(P^y\) the
content of this power is the minimum over \(j<y,\ j\le t\)
of \(t-j+v_p\binom tj\). After division by the lowest
uniformizer power, the least \(j\) attaining it has strictly
smaller order in the residue parameter than the other surviving
terms and than \(P^y\). This proves the minimum formula even for
nonconstant units.

Since multiplication by \(q^{pt-v}\) can only
increase content in the first quotient, any admissible index gives
an upper bound for \(c_{(P^y)}(q^v)\).
For \(D=P^u,J=(q^C)\), use the binomial expansion of
\((p\varepsilon+Uq^p)^u\) modulo \(q^C\). The two bounds are
\begin{equation}\label{eq:two-letter-wach-contents}
 \begin{aligned}
 c_{(P^y)}(q^v)&\le t-i_1+v_p\binom{t}{i_1},
       & i_1&=\min(t,y-1),\\
 c_{(q^C)}(P^u)&\le u-i_2+v_p\binom{u}{i_2},
       & i_2&=\min(u,\lfloor(C-1)/p\rfloor).
 \end{aligned}
\end{equation}
For \(pu<C\) the second line has \(i_2=u\) and the content is zero.
For the binomial terms in \eqref{eq:two-letter-wach-contents}, the carry formula gives \(v_p\binom ti\le\log_p t\).
In particular the two bounds are at most
\[
 \max(0,v/p-y+2+\log_p v)+1,\qquad
 \lceil\max(0,u-C/p+2+\log_p u)\rceil ,
\]
respectively. In the first bound the case \(i=t\) contributes
zero. All contents are normalized by \(v_p(p)=1\).

Suppose first that the large terminal letter is positive. If
\(p(u+S+\eta)+\eta<v\), construct a filtered source constant.
On the target take the coordinate \(P^{u+S}/q^v\) in direction
\(d\), realize it at cost \(\sigma\) and lift at the cut at
cost \(R\). Express the lift in source horizontal coordinates
at the cost of a further \(\sigma\). Multiply these coordinates
by \(q^{v+S}P^\eta\), obtaining integral series. In fact every
nonzero filtered summand has exponents at least \(-(v+S)\)
and \(-\eta\) at the two levels. Divisibility after rational
expansion descends by saturation. In the coordinate \(q\) take
the coefficient of degree
\[
 j=S+p(u+S+\eta).
\]

A filtered summand with level-zero exponent \(\alpha\) can
contribute only if
\[
 v+S+\alpha\le j.
\]
This forces
\(-\alpha\ge\eta+1\); the gap puts its direction in
\(\mathcal F_0^v\). The constant matrix of the map sends this
coefficient to the same coefficient of
\(p^{R+2\sigma}q^S P^{u+S+\eta}d\). Indeed the coordinate error
after clearing these denominators is divisible by \(q^h\) and
\(j<h\). The coefficient of degree \(p(u+S+\eta)\) of
\(P^{u+S+\eta}\) is a unit, by reduction modulo the maximal
ideal. Dividing by this unit supplies the filtered preimage of \(p^{R+2\sigma}d\). The additional coordinate comparison in the second rule raises the exponent to
\[
 (R+2\sigma)+\sigma=R+3\sigma.
\]


In the complementary range take \(D=q^v\), \(y=u+S\). The
first content bound separates the constant at the negative level
from the remainder. Each is lifted at cost \(R+\sigma\) by the
two rules. Writing the extra cost as \(c_{\rm cont}\), the two
ranges give
\[
 c_{\rm cont}\le
 \begin{cases}
    3+\gamma_h,& v<py,\\
    \eta+\eta/p+3+\gamma_h,& py\le v\le p(y+\eta)+\eta .
 \end{cases}
\]
In the second range, combining the two bounds on \(v\) with
\(b_0\le v\le b_0+\eta\) gives
\[
 b_0/p-S-\eta-\eta/p\le u\le(b_0+\eta)/p-S,
\]
the prescribed band.

Now suppose the first small letter is positive. Put
\(A=A_{c-1}\). When \(u\ge A+\eta+1\), take the partial
arc of length one at this letter and the coordinate \(q^{-u}d\)
there, with \(d\) read before the last Frobenius transport.
Realize it at cost \(\sigma\), lift \(X^A\) times this integral
vector, and read the source coordinates at cost \(\sigma\).
Multiply by \(q^G\) and take the coefficient in degree \(G-u+A\).

The degree \(G-u+A\) is below \(G-\eta\), so the degree and gap test places the coefficient in \(\mathcal F^u\). The target coefficient in the chosen degree is one, as
\[
 [q^{A+G-u}](q-p)^A q^{G-u}=1.
\]
The precision includes \(q^{2G}\) after denominators are cleared, so the image of the coefficient is exactly \(p^{2\sigma}d\). Tame averaging costs a unit.
Transport forward and apply the second rule at pole set
\(\{1\}\). The cost is \(3\sigma\).

When \(u\le A+\eta\), apply the second content bound with
\(D=P^u,\ C=v+S\), and lift the separated remainder and
constant at cost \(R+\sigma\). There is no extra cost for
\(pu<C\). Otherwise \eqref{eq:short-A} with its improved
constant gives
\[
 u-C/p\le\eta+\eta/p+2L.
\]

If the band is absent, \(pu\ge C\ge b_0\) and
\(pu-b_0>2p\eta\). But
\[
 u-b_0/p\le \eta+(2\eta+S)/p+2L,
\]
so \(\eta<4L\) using
\((2\eta+S)/p\le L\) (since \(w>\eta\)).
Thus the two-letter exponent in all cases is at most
\(15L+1+\eta B\). \label{cor:inert-short-two-2}
\end{proof}


\begin{lemma}[Three-letter lifting with negative outer signs]\label{lem:three-letter-negative-lift}
With \(L=S_g\) and \(B\) the short-band indicator of Section~\ref{sec:inert-row-table},
for the chronological signs \(-,+,-\), one can choose a three-letter correction in the short-arc frames with exponent \(e_3\) satisfying
\[
 e_3\le22L+1+\eta B
\]
in the given DVR normalization.
\end{lemma}

\begin{proof}
\label{par:short-three-levels}
Put
\[
 u=s_{c-2},\ v=s_{c-1},\ z=s_c,
\]
so \(u\le w\)
and \(v,z\in[b_0,b_0+\eta]\). Here \(40L\le b_0/4\).
At a partial arc use level numbers relative to its own endpoint.

\proofstep{Start with the middle positive letter and retain the two negative neighbors.}
First consider signs \(-,+,-\) in chronological order.
If \(p(u+S+\eta)\le v\), work at the middle letter and
write \(A=A_{c-1}\), \(y=u+S+\eta\).
As in the two-letter coefficient extraction use the coordinate
\(P_1^{u+S}/q^v\) in direction \(d\). At this intermediate
point the scalar image bound at a cut is unavailable. Instead
lift \(X^A\) times the realized vector. After the two horizontal
comparisons and multiplication by \(q^{v+S}P_1^\eta\), the
scalar on the target, in the variable \(q\), is
\[
 F=p^{2\sigma}(q-p)^A q^S P_1^y .
\]
Below degree \(v+S-\eta\) every source coefficient lies in
\(\mathcal F^v\) by the gap. To recover an integral direction
it suffices that some such coefficient of \(F\) divides a small
power \(p^t\) in \(O\).

Take
\[
 t=1+\left\lfloor\tfrac{11}{10}
 (2\sigma+(\eta+A)/p+2+\gamma_h)\right\rfloor .
\]
Then \(2\sigma<t\le10L\), \(A+t+\eta\le v\).
For the first bound use \(\sigma\le3L\) and
\(A\le (h+\eta)/(p-1)+1\), giving
\[
 (\eta+A)/p\le 1.1(h+\eta)/p^2+(\eta+1)/p.
\]

Substitution in the floor bound gives \(t\le10L\).
For the second, \(A\le (b_0+2\eta+S)/(p-1)+1\),
\(10L\le b_0/16\) and \(20\eta,20S\le b_0\).
If all the indicated coefficients were divisible by \(p^t\),
put \(k=t-2\sigma\). Modulo \(p^k\) one has
\(q^{A+k-1}\in((q-p)^A)\) by the binomial expansion.
Cancel \(p^{2\sigma}\), remove the initial power of \(q\),
and use this divisibility. The result would be
\[
 P_1^y\in(p^k,q^{v-\eta-(A+t-1)}).
\]
The second content calculation, before rounding, gives
\[
 k\le\max\{0,\ y-(v-\eta-A-t+1)/p+2+\log_p y\}.
\]
Indeed using the index \(\min(y,\lfloor(C-1)/p\rfloor)\)
costs at most \(y-C/p+1+1/p+\log_p y\) when positive.

Since \(py\le v\) and \(p/(p-1)\le11/10\), the definition
of \(t\) contradicts the last display. Choose a coefficient
of valuation less than \(t\), and multiply its source vector by
the integral quotient of \(p^t\) by that coefficient.
Moving one step forward gives the positive-level preimage.
The sole pole of \(f\) is at this level, so the second lifting
rule gives cost
\[
 t+\sigma\le13L.
\]


Suppose next that \(v<p(u+S+\eta)\). The down-level constants
in \eqref{eq:short-scalar} can be different. Eliminate the one
at level two first. The denominator there is \(P_1^v\), the
modulus is \(P_2^{u+S}\). Put \(y=u+S\). Since
\(P_2=p\varepsilon_1+U_1P_1^p\) with unit coefficients, the
first content calculation (using \(P_1,P_2\) in place of
\(q,P_1\)) gives a splitting
\[
 t_2f=x_2+f_2,\quad D f_2\in P_2^{y}\mathcal R,\quad
 D(p^{k_0}f_2-x_0)\in q^{z+S}\mathcal R,\qquad D=P_1^v .
\]

If \(v<py\), the content of \(P_1^v\) modulo \(P_2^y\)
is zero by reduction and comparison of orders. Otherwise use
the ceiling multiple \(p\lceil v/p\rceil\) in its place.
For \(py\le v<p(y+\eta)\), the expansion in \(p,P_1^p\) of \(P_2^y\) costs
at most \(\eta+2+\gamma_h\), and the small letter belongs to
the band. Thus \(v_p(t_2)\le(\eta+2+\gamma_h)B\).

Write \(C=z+S\). For \(x_0\ne0\), the second congruence gives
\begin{equation}\label{eq:short-last-constant}
 v_p(x_0)-k_0\ge y-v+C/p-C/p^2-2-\gamma_h .
\end{equation}
Indeed \(P_2^y\) has content at least \(y-C/p^2\) modulo
\(q^C\). Expand first in \(p,P_1^p\), then in \(p,q^p\):
a term with less scalar valuation has \(q\)-degree at least
\(C\). The content of \(P_1^v\) modulo \(q^C\) is at most
\(v-C/p+2+\gamma_h\) by the second calculation, since
\(C\le pv\). Comparing contents of
\[
 x_0P_1^v\equiv p^{k_0}P_2^y U\bmod q^C
\]
proves the claim.

Here is how to use this scalar inequality with the middle
filtered lattice. Transport the truncated horizontal column to
the middle point, paying \(\sigma\) at the first small letter,
and denote its integral representative by \(\lambda\). It
represents \(p^\sigma d\) there.

Set \(s=P_1^y/q^v\) at this
point. The vector \(s\lambda\) lies in the target Wach lattice.
The level test uses \(d\in\mathcal F^v\) at the middle and
the zero of order \(y\) at the initial negative letter.
In addition
\[
 \gamma(s)(\gamma-1)\lambda\in Xq^{a_*}N_{c-1}(T_2)
\]
where \(a_*=A_{c-1}+r_{c-1}\) and \(r_i\) is the least
scalar Wach image exponent at \(i\). Indeed the one-step
horizontal correction has defect in \(X^M q^{2G}\) in the
first Frobenius span. The next transport gives factors
\(X^M q^{(p-1)M}P_1^{2G}\); after multiplication by
\(\gamma(s)\) these pay \(Xq^{a_*}\) and the signed height losses
\(q^\eta P_1^{u+S}\), since \((p-1)M\ge v+a_*+\eta\).

The powers required in the ordinary parameter and the scalar
give the \(q^{a_*}\) image by the binomial theorem.
The two-filter argument of Paragraph~\ref{par:thick-filtered-comparison}
therefore lifts \(p^E s\lambda\) integrally, where
\[
 E=\left\lfloor\frac{3a_*+S+2\eta-2}{p-1}
               +\log_p(v+\eta-1)\right\rfloor_{\!+}\le8L .
\]
For this application the lower filter has indices \([v-a_*+1,v]\)
and the other has indices in \([v+1,v+S+a_*]\cup[1-\eta,\eta+a_*]\).
The factors adjacent to \(j\) in the first two intervals
give factorials of combined length \(2a_*+S-1\).

In the remaining
interval of length \(2\eta+a_*\), isolate a term of largest
valuation; the others are bounded by factorials of combined
length \(2\eta+a_*-1\). The isolated valuation is at most
\(\log_p(v+\eta-1)\). If \(a_*=0\) the map is already onto
at this letter. The global ordinary width bound \(H\le h+\eta\)
and degree lowering give
\[
 a_*\le H(2p-1)/(p-1)^2+\lceil1+\log_p(1+H)\rceil+1.
\]

These imply \(2a_*+\eta+1\le v\) under \(40L\le b_0/4\)
and give the displayed estimate for \(E\), by \(p\ge11\) and
\(H\le b_0+2\eta+S\). In particular
\[
 v+a_*+\eta\le2v\le2h<(p-1)M,
\]
as used in the defect estimate. The separation gives the
disjointness required by interpolation. \label{cor:inert-short-three-1}

Let \(n_1\) be this lift. Take its constant in a horizontal
basis at \(c-1\) and transport it to \(c\) with precision
\(Xq^C\). This gives \(y_1\) in the last Frobenius span,
with \((\gamma-1)y_1\in XN_c(T_1)\) and \(q^C y_1\in N_c(T_1)\).
Use the truncation order
\[
 R_0=\lceil C/(p-1)\rceil
\]
and
take \(\sigma_1=R_0\). In positive degree \(r<R_0\) the
denominator costs \(r+v_p(r!)\). The transported monomial is
a unit times \(X^r q^{(p-1)r}\). Keep one \(X\) and expand
the others as \((q-p)^{r-1}\).

Modulo \(Xq^C\) the surviving
terms contain at least \(\max(0,pr-C)\) powers of \(p\):
a term of extra degree \(j\) must satisfy \(j+(p-1)r<C\).
The loss is therefore at most
\[
 pr/(p-1)-\max(0,pr-C)\le C/(p-1);
\]
the maximum occurs at the crossing \(pr=C\). The transported
tail is zero to this precision. The last-level jet of \(y_1\)
is horizontal in the last Frobenius span: its difference is
zero modulo \(Xq^C\) there, by the truncation and the integral
correction. Over the coefficient field use the local logarithmic
parameter at the primitive level. On a horizontal column the
action on degree \(k\) has eigenvalue \(\chi(\gamma)^k\), so
an invariant jet has only its constant coefficient. This gives
the jet of \(F(y_1)\) from the transported constant of \(n_1\)
with multiplier \(p^{\sigma_1}\).

To compare with \(f_2\Phi\lambda\), use the common horizontal
column at the last level. At this level the transported multiplier
\(\phi(s)=P_2^y/P_1^v\) has constant term
\[
 (P_2^y/P_1^v)(0)=p^{y-v},
\]
since \(P_j(0)=p\) for \(j\ge1\) in the \(q\)-coordinate.
Thus the two scalar constants to be matched, after multiplying
by \(p^{k_0+v}\), are \(x_0p^v\) and
\(p^{k_0+\sigma_1+E+y}\). In the second expression the
common horizontal normalization cancels; its \(p^\sigma\) is
present in both comparisons.

Choose \(t_3,B'\in O\) with
\[
 t_3 x_0p^v=B' p^{k_0+\sigma_1+E+y},\qquad v_p(t_3)\le14L .
\]
For \(x_0=0\) use \((1,0)\). Otherwise take \(t_3=1\) if
\(v_p(x_0)+v\ge k_0+\sigma_1+E+y+\sigma\). In the other
case choose equality with \(B'=p^\sigma\). Equation
\eqref{eq:short-last-constant} bounds \(v_p(t_3)\) by
\[
 \sigma+E+\sigma_1-C/p+C/p^2+2+\gamma_h\le14L .
\]
In fact \(\sigma_1\le C/(p-1)+1\), \(C\le h\), and
\(1/(p(p-1))+1/p^2\le 2.1/p^2\).

For the marked map \(F\), the required Frobenius-span congruence is
\[
 t_3 f_2\Phi\lambda-B' F(y_1)\in q^C\langle\Phi N_{c-1}(T_2)\rangle.
\]

For the first term the numerator at level two supplies the
height loss from the first negative letter; equivalently
\(f_2\Phi\lambda\) is an integral multiple of
\(\Phi(s\lambda)\). For the divisibility assertion compare
the two horizontal jets after inverting \(p\) and the higher
level factors. Modulo \(q^C\) the first has constant multiplier
\(x_0/p^{k_0}\) on \(\Phi\lambda\), by the last congruence;
the horizontal jet of the second has constant multiplier
\(p^{\sigma_1+E+y-v}\). Higher level factors are units at
\(q=0\) over the fraction field, and the bounded integral
Frobenius span is free.

For the free bounded integral Frobenius span \(P\), monic division yields
\[
 P\cap q^C P[1/p]=q^C P.
\]
Thus the congruence of rational horizontal jets is integral. The negative terminal
height puts this difference in \(N_c(T_2)\). Lift \(p^R\)
times the difference there and add \(p^R B' y_1\).
The result is admissible and maps to \(p^Rt_3 f_2\Phi\lambda\);
its denominators in \(N_c(T_1)\) are cleared by \(q^C\).
Since \(\Phi\lambda\) represents \(p^\sigma d\), this
lifts \(t_3f_2\) at cost \(R+\sigma\). The summand \(t_3x_2\)
admits the constant lift at the same cost. The scalar clearing cost is
\[
 \lceil v_p(t_3t_2)\rceil=\lceil v_p(t_3)+v_p(t_2)\rceil.
\]
Adding the lifting costs gives the bound \(22L+1+\eta B\).
\end{proof}


\begin{lemma}[Three-letter lifting with positive outer signs]\label{lem:three-letter-positive-lift}
For the same \(L=S_g\) and short-band indicator \(B\) as in
Lemma~\ref{lem:three-letter-negative-lift}, and for chronological signs \(+,-,+\), the adjacent-level correction can be chosen with exponent \(e_3\) satisfying
\[
 e_3\le34L+1+\eta B.
\] With the other short-arc branches, the
single exponent \(E=\lfloor40S_g+\eta B\rfloor\) works on
every eligible column of the chosen horizontal frames,
in powers of \(\ell\) in the original DVR.
\end{lemma}

\begin{proof}
\proofstep{Use the middle congruence with the two positive neighbor denominators.}
Now consider \(+,-,+\). The denominator of \(f\) is
\(D=q^zP_2^u\) and its middle congruence is modulo
\(P_1^{v+S}\). Write \(C=v+S\), \(a_0=\lceil z/p\rceil\).
We first show that the middle rational constant \(x/p^k\)
has denominator at most \(p^s\), with
\begin{equation}\label{eq:three-denominator}
 s\le\max\{1,u-C/p+16L/5\}+1 .
\end{equation}
Put \(Q=u-(C-a_0+1)/p+2+\log_p u\), and use
\(s=1+\lfloor\max(0,Qp/(p-1))\rfloor\).

If \(D\in(p^s,P_1^C)\), multiply by \(q^{pa_0-z}\).
Expanding \(P_1=p\varepsilon+Uq^p\) gives
\(P_1^{a_0+s-1}\in(q^{pa_0},p^s)\).
Hence \(P_1^{a_0+s-1}P_2^u\in(p^s,P_1^C)\).
Cancelling this power of \(P_1\) modulo \(p^s\), and descending
by \(\phi\), gives
\[
 P_1^u\in(p^s,q^{C-a_0-s+1}).
\]

Cancellation holds by reduction to a distinguished polynomial,
and descent holds since the ring is free over its Frobenius
image (Weierstrass division by \(\phi(q)-Y\)).
The content inequality from the preceding paragraph would force
\[
 s\le\max(0,u-(C-a_0-s+1)/p+2+\log_p u),
\]
contrary to the choice. Here \(C-a_0-s+1>0\):
\(a_0\le z/p+1\), \(20u\le b_0\) and the bound on \(s\)
in \eqref{eq:three-denominator} gives \(a_0+s<C+1\) using
\(160L\le b_0\).

To verify that bound directly, use
\((a_0-1)/p\le h/p^2\), \(u\le w\) and \(\log_p u\le\gamma_h\).
Expanding the factor \(p/(p-1)=1+1/(p-1)\) charges at most
\((w+h/p^2+2+\gamma_h)/(p-1)\) in addition to
\(u-C/p+h/p^2+2+\gamma_h\), which fits the displayed allowance.
This proves the bound on the content of \(D\); the middle
congruence then gives
\[
 p^sx/p^k\in O.
\]


Put \(A=A_{c-2}\). If \(u<A+\eta+2\sigma+2\), separate the
constant at cost \(s\) and lift both summands at \(R+\sigma\).
Indeed the new remainder has numerator divisible by \(P_1^C\).
We check that the size-dependent loss in
\eqref{eq:three-denominator} fits the band allowance.
By \eqref{eq:short-A},
\[
 u-C/p<\eta+\eta/p+\tfrac52 L+2\sigma+2.
\]
Inside the band this gives \(s\le\eta+14L+1\). Outside it,
if \(pu\le b_0\) the excess \(u-C/p\) is nonpositive.
If \(pu>b_0\), absence of the band gives
\(u-b_0/p>2\eta\), while the last upper bound and
\((C-b_0)/p\le(\eta+S)/p\) imply \(\eta<12L\).
These bounds give
\[
 s\le 28L+1+\eta B
\]
on every arc.

If \(u\ge A+\eta+2\sigma+2\), construct a filtered constant
at the initial positive letter. The one-letter calculation with
the coordinate \(q^{-u}d\) gives scalar comparison
\(p^{2\sigma}(q-p)^A q^{G-u}\) modulo \(q^{G-\eta}\).
The degree \(G-u+A\) is in the comparison range. Its leading coefficient is
\[
 [q^{G-u+A}]\,p^{2\sigma}(q-p)^A q^{G-u}=p^{2\sigma}.
\] Equivalently, some coefficient
in that range divides \(p^{2\sigma+1}\): expanding
\(q^{A+2\sigma}\) in \((q-p)+p\) shows that divisibility of all
the low coefficients by \(p^{2\sigma+1}\) is impossible.
All these source coefficients lie in \(\mathcal F^u\). Thus
after transport a source direction at positive level two maps
to \(p^{2\sigma+1}d\). Only its filtration at level two has
been prescribed. We choose the part of \(f\) to which it applies.

Let
\[
 j=\max(0,\lfloor u-C/p-\eta/p^2\rfloor).
\]
There is a
decomposition
\[
 P_2^u=p^j U_0+P_1^C q^\eta U_1,\qquad U_0,U_1\in\mathcal R .
\]
In the expansion of \((p\varepsilon_1+U'_1P_1^p)^u\), the
term with \(P_1^{pr}\) is paid by its scalar coefficient if
\(r\le C/p+\eta/p^2\). Otherwise factor off \(P_1^C\) and
expand the remaining \(pr-C\) powers of \(P_1\).
Any resulting term with \(q\)-degree below \(\eta\) has
scalar valuation at least
\(u-r+(pr-C)-\eta/p\ge u-C/p-\eta/p^2\).
This proves the decomposition. Put \(e_4=(s-j)_+\) and
\(x'=p^s x/p^k\).

Take
\[
 f_a= p^{e_4+j-s}x' U_0/P_2^u.
\]
Then \(p^sf_a-p^{e_4}x'\) belongs to
\(P_1^Cq^\eta P_2^{-u}\mathcal R\).
The second rule lifts \(f_a\) at cost \(3\sigma+1\), using
the source filtration at level two and the constant congruence
at the omitted positive level zero. The remainder
\(p^{e_4}f-f_a\) has numerator divisible by \(P_1^C\) over
\(D\), by the middle congruence and saturation. It costs
\(R+\sigma\). Since \(j\ge u-C/p-\eta/p^2-1\),
\eqref{eq:three-denominator} gives
\[
 e_4\le 16L/5+\eta/p^2+2
\]
when the second term in the maximum is used; otherwise
\(e_4\le2\). In particular \(R+3\sigma+1+e_4\le18L+1\).
The two initial-positive cases therefore both fit
\(34L+1+\eta B\). \label{cor:inert-short-positive-1}

All the bounds for the individual scalar lifts are at most
\(40S_g+\eta B\), since \(S_g\ge1\).
They hold on every eligible column of every chosen horizontal
frame; consequently the integer
\(E=\lfloor40S_g+\eta B\rfloor\) can be used simultaneously.
The exponents are in powers of \(\ell\) in the original DVR,
including when its uniformizer has fractional normalized valuation.
\end{proof}


\begin{lemma}[A fixed exceptional span for a short arc]\label{lem:short-arc-exceptional-span}
For the short arc of length \(a\) in the cut comparison,
one may choose an exceptional set \(Z_c\) of target classes, independent of the input, with
\[
 \#Z_c\le(a+1)ma.
\] Outside its span the scaled lift has
exponent \(\sigma+E+\kappa_c\). The signed count is
\(r_c=a\) because each letter has an off-central window.
\end{lemma}

\begin{proof}
\label{par:short-exceptions}
On the marked target over the coefficient field let \(U\) be the
intersection of the transported \(\mathcal F_j^{u_j}\) over
positive levels. Let \(V_-\) be the sum of
\(\mathcal F_j^{-\eta}\) over negative levels. By the bound on
central multiplicities,
\[
 \operatorname{codim}U+\dim V_-\le ma .
\]
In the initial integral basis, the constants form the lattice
\(D_O\) with the transport of scalars to the terminal crystalline
space understood. Split the saturated sublattice \(D_O\cap U\)
off with complement \(K\), and project \(V_-\) along \(K\)
to \(U\). The intersection of the projected space with the
lattice of \(U\) is again saturated. Use adapted bases of these
two splittings.

Let \(e_j\) denote the adapted basis and let \(J\) index the complement \(K\) and this last intersection. Then
\[
 |J|\le ma,\qquad V_-\subseteq\operatorname{span}_{K_O}\{e_j:j\in J\}.
\]
Every other column belongs to \(U\).
The two changes of basis are integral and invertible, so they
give a Wach frame at \(c'\). Take the unique horizontal columns
with these initial values.

Apply the cut construction of Section~\ref{sec:inert-long-orbit}
to a crystalline extension representing a finite target class.
Its cut vector is the difference in that extension of a vector
in the incoming Frobenius span and a terminal Wach vector.
At each cyclotomic level the incoming span is regular; a pole
of order \(k>0\) therefore has coefficient in the extension's
\(\mathcal F^k\), and its image lies in the kernel of projection
to the trivial representation. Strictness identifies this
intersection with \(\mathcal F^k\) of the target.

For a polar term \(a_{-k}t^{-k}\) with \(k>0\), strictness therefore gives the two respective tests
\[
 a_{-k}\in V_-\quad(\text{negative letter}),\qquad k\le u_j+S\quad(\text{positive letter}).
\]
The statement for an entire polar part is obtained
coefficient by coefficient in the Laurent expansion, followed
by Chinese remaindering on the finite jets. At negative levels
all negative degrees thus disappear on indices outside \(J\).
Likewise \((\gamma-1)w\in XN_c\) imposes the Wach filtration
on the divided differences. Their projection outside \(J\)
vanishes at a negative level in degrees below \(u_j\):
for lower degrees the required filtration is contained in
\(\mathcal F_j^{-\eta}\), or equals it by the gap.

Write \(\mathcal J_c\) for the union of signed degree intervals
used in \(\kappa_c\).
Apply to these horizontal coordinates the integral operator
\[
 K(\gamma)=\left(\sum_{j=0}^{p^a-1}\gamma^j\right)
     \prod_{k\in\mathcal J_c\setminus\{0\}}
            (\gamma^{p^a}-\chi(\gamma)^{kp^a}) ,
\]
On a Laurent coefficient of
degree \(k\) at level \(t<a\), the principal action raised to
\(p^a\) is the scalar \(\chi(\gamma)^{kp^a}\). The sum
projects degree zero to a common value on the primitive roots
after tame invariants. Outside \(J\) this removes the excess
positive-level poles and the negative-level degrees from \(u_j\)
through \(u_j+S-1\).

At a negative level the divided difference condition on a nonconstant coefficient \(a_k\) below \(u_j\) reads
\[
 (\chi(\gamma)^{kp^a}-1)a_k=0,\qquad k\ne0.
\]
The scalar is nonzero, hence \(a_k=0\).
Thus the transformed scalar coefficients outside \(J\) satisfy
\eqref{eq:short-scalar}. The remaining constant at each negative
level is in the coefficient field, so a scalar power clears it.
The operators preserve the bounded integral denominator
\(\Lambda\); the improved pole orders descend to \(D\) by
division by the distinguished factors. Outside \(J\), the only poles allowed by \(D\) are positive with directions in \(U\); the negative zeros have the required orders. Write \(f=f_J+f_{J^c}\) for the support decomposition and \(\Delta\) for divided difference. The Wach test is linear, so
\[
 \Delta f_J=\Delta f-\Delta f_{J^c}
\]
passes as well.

For the latter part let \(M_J\) consist of the tuples with
denominator \(\Lambda\) supported on \(J\); let \(I_J\) be
the submodule whose divided differences pass the Wach test.
Let \(P_J\) consist of \(X\) times the tuples passing the Wach test
with bounded denominator, in \(M_J\). For numerators supported on \(J\),
a change by \(X\Lambda^2h\) induces the following change of tuple:
\[
 \frac{f+X\Lambda^2 h}{\Lambda}-\frac{f}{\Lambda}
       =X\Lambda h\in P_J .
\]
Weierstrass division reduces each numerator modulo \(X\Lambda^2\);
hence \(M_J/P_J\) is finite over \(O\).

The quotient \(M_J/P_J\) is also torsion-free.
Suppose \(0\ne c\in O\) and \(c(f/\Lambda)\in P_J\). Monic division by
\(X\) gives \(cf(0)=0\), hence \(f(0)=0\). The tuple
\(f/(X\Lambda)\) now has bounded denominator. Its rational filtration
tests hold since \(c f/(X\Lambda)\) satisfies them and the filtration
spaces are vector spaces over the coefficient field. Thus
\[
 f/\Lambda\in P_J.
\]


Let \(E_J\) be the image of \(I_J\) in \(M_J/P_J\).
After inverting \(p\), Chinese remainders for \(X\) and the
distinct cyclotomic factors of \(\Lambda\) split the numerator
jets at \(0\) and at the \(a\) levels. At \(0\) the quotient by
\(X\) leaves at most one entry per coordinate in \(J\).
At a cyclotomic level \(X\) is invertible, so the quotient is
by the filtered Laurent lattice.

In coefficient degree \(k\),
the principal power fixing the roots acts on that quotient by
\(\chi(\gamma)^{k p^a}\). A nonzero fixed class therefore needs
\(k=0\), since \(\chi(\gamma)\) has infinite order. In degree
zero the tame and principal root actions are transitive. Invariants
have at most one copy of each supported coordinate, and passing
to the quotient cannot increase this invariant dimension by
the characteristic-zero averaging projection. Thus
\[
    \dim_{\operatorname{Frac}(O)}(E_J[1/p])
          \le |J|+\sum_{t=0}^{a-1}|J|=(a+1)|J| .
\]
Since \(M_J/P_J\) is finite torsion-free over \(O\), this is
also a bound on the number of generators of \(E_J\).

Take an \(O\)-basis of the image and choose its representatives
in \(I_J\) once for the cut. Realize \(p^\sigma\) times each
by the horizontal comparison and take its cut Herr class. The
resulting set \(Z_c\) consists of finite classes, since a
product of cyclotomic factors carries the representative into
the Wach lattice. Using one representative per basis vector of \(E_J\) gives the uniform bound
\[
 |Z_c|\le\operatorname{rank}_O E_J\le (a+1)ma.
\]

\proofstep{Reduce every class against the same chosen representatives.}
For any input the exceptional coordinates are an integral
combination of these representatives modulo \(P_J\).
Terms of \(P_J\) represent elements of \(XN_c\) after the
same comparison and give zero under the cut class map.
The other coordinates lift at the common cost \(E\). To see the
exponent accounting, the inverse horizontal comparison reads
coordinates of \(p^\sigma w\). Realizing the split coordinates
requires another \(p^\sigma\), while on the outside part the
vector lifting relation already represents \(p^E\) times the
coordinates.

Since \(E\ge\sigma\), pad the realized split parts
by \(p^{E-\sigma}\). Their sum now matches
\[
 p^{\sigma+E} K(\gamma)w
\]
modulo \(XN_c\). Naturality and the vanishing on
\(XN_c\) follow from Frobenius contraction in the Herr complex,
as in the cut construction. Write \(K(T)-K(1)=(T-1)Q(T)\). Stability and the cut-vector condition give
\[
 K(\gamma)w-K(1)w=Q(\gamma)(\gamma-1)w\in XN_c.
\]
Thus the two vectors give the same cut class. In the scalar reading
\[
        p^{\sigma+E}K(1)w
             = u_c p^{\sigma+E+\kappa_c}w,\qquad u_c\in O^\times
\]
since \(v_p(K(1))=\kappa_c\). This gives the exponent
\(\sigma+E+\kappa_c\) with the same set \(Z_c\) for every
input class. For the present signed arcs every letter has a
positive or negative off-central window, so the signed level count
\(r_c=a\). \label{cor:inert-short-exceptions-1}
\end{proof}



\paragraph{Adding the short-arc colengths.}
\begin{lemma}[The common colength from the short cuts]
\label{lem:ordinary-inert-short-cuts}
In either ordinary inert range with a position having smaller jumps, use its
marked arcs. The normalized finite-class colength is bounded by
\eqref{eq:nonthick-colength}; in the range requiring exceptional representatives, there is one set \(Z\) for all the comparisons such that
\[
 \#Z\le4nm.
\]
\end{lemma}

\begin{proof}
Apply Lemma~\ref{lem:short-arc-exceptional-span} with \(p=\ell\) and \(a=a_c\).
On the finite jets at levels \(t<a_c\), the action of \(\gamma^{\ell^{a_c}}\)
on a Laurent coefficient \(c_k\) of degree \(k\) is multiplication by
\(\chi(\gamma)^{k\ell^{a_c}}\). For each specified nonzero degree
the corresponding factor of \(K(\gamma)\) therefore satisfies
\[
 \bigl(\gamma^{\ell^{a_c}}-\chi(\gamma)^{k\ell^{a_c}}\bigr)c_k=0 .
\]

To estimate \(\kappa_c\), count these degrees with their signs. For a positive
jump window these degrees lie in \([-s_i-S,-s_i)\) and for a
negative window in \([s_i,s_i+S)\). Thus there are at most
\(a_c S\) degrees, each nonzero with absolute value at most \(h\).
In any one of these intervals of \(S\) integers the count of
multiples of \(\ell^j\) is at most \(S/\ell^j+1\).
For such an interval \(I\) excluding zero the resulting valuation bound is
\[
 \sum_{k\in I}v_\ell(k)=\sum_{j=1}^{\lfloor\log_\ell h\rfloor}\#\{k\in I:\ell^j\mid k\}\le S/(\ell-1)+\log_\ell h.
\]

The principal trace contributes at most \(a_c\), and the
nonzero degrees \(k\) contribute at most \(a_c+1+v_\ell(k)\) each,
by the principal-unit valuation used in
Section~\ref{sec:inert-long-orbit}. Hence
\[
 \kappa_c\le a_c+(a_c+1)a_c S
                +a_c(S/(\ell-1)+\log_\ell h)
            \le13 S+3(1+\gamma_h).
\]
Given the cut decomposition and the arc lift, pad the powers to
their common maximum and sum the lifted cut classes.
The set \(Z\) of exceptional classes
is contained in the union of the sets from the cuts, so
\[
 |Z|\le\sum_c m r_c(a_c+1)\le4m\sum_c a_c=4nm .
\]
The cost of the scaling on the original class is at most
\(55S_g+13S+\eta B\), since \(1+\gamma_h\le S_g\).
Thus the length argument of Section~\ref{sec:inert-thick-word}
applies to the same integral finite group.

In the larger range
one uses the full lift and takes \(Z\) empty. In both ranges the
normalized colength is at most
\begin{equation}\label{eq:nonthick-colength}
 320(n+1)r\bigl(h/\ell^2+(w+\eta)/\ell+
                      S+1+\gamma_h+\eta B\bigr)+320nm e' .
\end{equation}
Indeed \(S/\ell\le S\), all the terms in the parentheses
are nonnegative, and \(55+13\le320\). Lemma~\ref{lem:short-arc-exceptional-span}
supplies the fixed representatives for this estimate. The coarse exponent
\(e'\) is for the actual marked surjection and is used only on the
image of the span of the exceptional finite classes.
\end{proof}

\begin{proposition}[One normalized entry for all ordinary inert rows]\label{prop:ordinary-inert-place-entry}
In the ordinary inert range, with the marked word data and
\(U_H,U_F,U_M\) defined in Section~\ref{sec:cut-length-rows},
there is one rational place entry \(c_w\) for every good character at the place. Its range is
\[
 0\le c_w-(368U_H+2082U_F+22365U_M)\le U_F.
\]
Here the indicator entering \(U_F\) is the common place indicator,
and every normalized length uses the same coefficient valuation.
\end{proposition}

\begin{proof}
The common unramified coarse finite-class lifting theorem uses the
negative and positive bounds \(M_i,H_i^+\) at each embedding.
Put \(M_*=\max_i M_i\), \(H'=\max_i(M_i+H_i^+)\)
and \(a_*=\lfloor M_*/(\ell-1)\rfloor+1\).
If \(J\) is the set of nonzero degrees formed by adding
\(1,\ldots,M_*\) to jumps and also including \(1,\ldots,M_*\),
the exponent for this theorem is
\[
  2\left\lfloor\frac{H'\ell}{(\ell-1)^2}+3+
                        \log_\ell(H'+1)\right\rfloor
       + a_*+v_\ell(a_*!) +\sum_{j\in J}v_\ell(j).
\]
The theorem supplies a finite lift for every target class after
multiplication by this power of \(\ell\).

In the general boxes \(J\subset[-H',2H']\setminus\{0\}\).
Legendre's formula bounds the valuation sum by \(3H'/(\ell-1)\).
Also
\(a_*+v_\ell(a_*!)\le(\ell/(\ell-1))(H'/(\ell-1)+1)\).
Using \(\ell/(\ell-1)\le11/10\) gives
\[
 e'\le23\bigl(2+H'/\ell+\log_\ell(H'+1)\bigr).
\]
For the large-\(S_g\) boxes above, substitute \(J\subset[-h,2h]\),
\(M_*\le h\), and \(H'\le h+\eta\le21h/20\) into the same
expression. Adding the coefficients of \(h/\ell\) from the three terms gives
\[
 2(21/20)(11/10)^2+(11/10)^2+33/10<7.1.
\]

The constants other than the logarithm total at most \(7.1\),
and \(H'+1\le2h\) since \(H'\) is integral and \(h\ge1\).
This proves the exponent estimate used in that range.

Use the global \(H,w_P,S_c,\eta_3,z\) and the quantities
\(U_H,U_F,U_M,\mathcal P_\ell\) of
Section~\ref{sec:cut-length-rows}. The local transport gives
\(S\le S_c\), \(\eta\le\eta_3\), \(w\le w_P\), and \(h\le H\)
in both branches. In the case in which every position has larger jumps this last inequality
means \(\max_i b_i+S\le H\). The polynomial size bounds give
\(1\le H\le2e^z\), \(z\ge1\) and \(\eta_3\le H\).
When some jumps are in the smaller range, \(H'\le2H\). In the
case in which all jumps are in the larger range
\[
 H'\le H_{\rm loc}\le H+\eta_3\le2H.
\]

As \(H'+1\le3H\le\ell H\) for the logarithm's argument, it follows
that \(e'\le69\mathcal P_\ell\).

For \eqref{eq:nonthick-colength}, \(\gamma_h\le3+z\) follows from
\(2h\le4e^z\) and \(\log\ell\ge1\).
Consequently \(1+\gamma_h\le5z\). The row length is at most
\(320U_H+1600U_F+22080U_M\).

For Proposition~\ref{prop:thick-inert-colength},
\(L_0\le8+z\le9z\), since
\(2H_{\rm loc}\le8e^z\). Further,
\[
 32H_{\rm loc}/\ell^2\le32H/\ell^2+3\eta_3/\ell,\qquad
 10\eta/(\ell-1)\le11\eta_3/\ell .
\]
The exponent \(e\) is therefore at most
\(32H/\ell^2+225F_{\ell,\lambda}\).
Also \(H_{\rm loc}+1\le3H\), which gives
\(c_D\le4\mathcal P_\ell\).
Substituting in \eqref{eq:thick-chain-mass} bounds the colength
by \(32U_H+227U_F+288U_M\).
Both branches thus have common majorant
\[
                       368U_H+2082U_F+22365U_M .
\]
The place entry is a rational number between this value and its
sum with \(U_F\). Here \(U_F>0\) by \(z\ge1\) and \(r>0\).
All comparisons use the same place indicator and coefficient
normalization. Hence the entry bounds the row for every
good character at once.
\end{proof}


